\chapter{Point-to-Point Communication}
\mpitermtitleindexsubmain{point-to-point}{communication}
\label{sec:pt2pt}
\label{chap:pt2pt}

\section{Introduction}
\label{sec:pt2pt-intro}

Sending and receiving of \mpitermni{messages} by \MPI/ processes is the basic \MPI/
communication mechanism.
The basic point-to-point communication operations are \mpiterm{send} and
\mpiterm{receive}.  Their use is illustrated in Example~\ref{ex:pt2pt-helloworld}.

\begin{example}
\label{ex:pt2pt-helloworld}
\exindex{Point-to-point!Hello world}%
\Exindex{C}{Hello world}{MPI\_Init,MPI\_Comm\_rank,MPI\_Send,MPI\_Recv}%
A simple `hello world' example usage of point-to-point communication.

%% FIXME: Use MPI_STATUS_IGNORE instead of declare, ignore status?
%%HEADER
%%LANG: C
%%DECL:#include <string.h>
%%ENDHEADER
\begin{lstlisting}[language={[MPI]C}]
#include "mpi.h"
int main(int argc, char *argv[])
{
  char message[20];
  int myrank;
  MPI_Status status;
  MPI_Init(&argc, &argv);
  MPI_Comm_rank(MPI_COMM_WORLD, &myrank);
  if (myrank == 0)    /* code for process zero */
  {
      strcpy(message,"Hello, there");
      MPI_Send(message, strlen(message)+1, MPI_CHAR, 1, 99,
               MPI_COMM_WORLD);
  }
  else if (myrank == 1)  /* code for process one */
  {
      MPI_Recv(message, 20, MPI_CHAR, 0, 99, MPI_COMM_WORLD, &status);
      printf("received :%s:\n", message);
  }
  MPI_Finalize();
  return 0;
}
\end{lstlisting}
\end{example}

In Example~\ref{ex:pt2pt-helloworld}, process zero (\code{myrank = 0}, strictly `the \MPI/ process with rank \code{0} in communicator \mpiconst{MPI\_COMM\_WORLD}') sends a \mpiterm{message} to process one
using the
\mpiterm{send} operation \mpifunc{MPI\_SEND}. The
operation specifies a \mpitermni{send buffer}\mpitermindex{send!buffer} in the sender memory from which the
\mpitermni{message data}\mpitermindex{message!data} is taken.  In the example above, the send buffer consists of the
storage containing the variable \code{message} in the memory of process zero.
The location, size and type of the send buffer are specified by the first three
parameters of the send operation.  The message sent will contain the 13
characters of this variable.
In addition, the send operation associates an \mpiterm{envelope}\mpitermindex{message!envelope} with the
message.  This \mpiterm{envelope}\mpitermindex{message!envelope} specifies the message destination and contains
distinguishing information that can be used by the \mpiterm{receive} operation to
select a particular message.
The last three parameters of the send operation, along with the rank of the
sender,
specify the \mpiterm{envelope}\mpitermindex{message!envelope} for the message sent.

Process one (\code{myrank = 1}, strictly `the \MPI/ process with rank \code{1} in communicator \mpiconst{MPI\_COMM\_WORLD}') receives this message with the
\mpiterm{receive} operation \mpifuncshort{MPI\_RECV}.
The message to be received is selected according to the value of its
\mpiterm{envelope}\mpitermindex{message!envelope}, and the \mpitermni{message data}\mpitermindex{message!data} is stored into the
\mpitermni{receive buffer}\mpitermindex{receive!buffer}.
In the example above, the receive buffer consists of the storage
containing the string \code{message} in the memory of process one.
The first three parameters of the receive operation specify the location, size
and type of the receive buffer.  The next three
parameters are used for selecting the incoming message.  The last parameter is
used to return information on the message just received.

\begin{users}
 Colloquial usage commonly permits references to ``rank 0'' or ``process 0'',
 which are strictly ambiguous and ideally should be qualified by including
 the relevant context, for example, the \MPI/ communicator in the case above.
\end{users}

The next sections describe the blocking send and receive
operations.  We discuss send, receive, blocking communication semantics,
type matching requirements, type conversion in
heterogeneous environments, and more general communication modes.
Nonblocking communication is addressed next, followed by
probing and cancelling a message,
channel-like constructs and send-receive operations,
ending with a description of the
``dummy'' \MPI/ process, \mpiconst{MPI\_PROC\_NULL}.

\section{Blocking Send and Receive Operations}
\mpitermtitleindex{point-to-point communication!blocking}
\mpitermtitleindex{blocking!point-to-point}
\label{sec:pt2pt-basicsendrecv}

\subsection{Blocking Send}
\mpitermtitleindex{send!blocking}
\mpitermtitleindex{point-to-point communication!send operation}
\label{subsec:pt2pt-basicsend}

The syntax of the \mpitermdefni{blocking send}\mpitermdefindex{send!blocking} procedure is given below.

\cdeclmainindex{MPI\_Comm}\cdeclmainindex{TYPE(MPI\_Comm)}
\begin{mpi-binding}
    function_name("MPI_Send")

    parameter(
        "buf", "BUFFER", desc=r"initial address of send buffer", constant=True
    )
    parameter(
        "count",
        "POLYXFER_NUM_ELEM_NNI",
        desc=r"number of elements in send buffer",
    )
    parameter(
        "datatype", "DATATYPE", desc=r"datatype of each send buffer element"
    )
    parameter("dest", "RANK", desc=r"rank of destination")
    parameter("tag", "TAG", desc=r"message tag")
    parameter("comm", "COMMUNICATOR")
\end{mpi-binding}

The blocking semantics of this call are described in Section~\ref{sec:pt2pt-modes}.

\subsection{Message Data}
\mpitermtitleindexmainsub{message}{data}
\label{subsec:pt2pt-messagedata}

The send buffer specified by the \mpifunc{MPI\_SEND} procedure consists of
\mpiarg{count} successive entries of the type indicated by
\mpiarg{datatype}, starting
with the entry at address \mpiarg{buf}.  Note that we specify the
message length
in terms of number of \emph{elements}, not number of \emph{bytes}.  The former is
machine independent and closer to the application level.

The data part of the message consists of a sequence of \mpiarg{count}
values, each of the type indicated by \mpiarg{datatype}.
\mpiarg{count} may be zero,
in which case the data part of the message is empty.
The \mpitermdef{basic datatypes} that can be specified
for message data values correspond to
the basic datatypes of the host language.
Possible values of this argument for Fortran and the
corresponding Fortran types are listed
in Table~\ref{table:pttopt:datatypes:fortran}.
%
Possible values for this argument for C and the corresponding C
types are listed
in Table~\ref{table:pttopt:datatypes:c}.

\begin{table}[t]
\caption{Predefined \MPI/ datatypes corresponding to Fortran datatypes\mpitermdefindex{basic datatypes!Fortran}}
\label{table:pttopt:datatypes:fortran}
\begin{center}
\begin{tabular}{|l|l|}
\hline
\textbf{\MPI/ datatype} & \textbf{Fortran datatype} \\
\hline
\typemain{MPI\_INTEGER} & \ftype{INTEGER} \\
\typemain{MPI\_REAL} & \ftype{REAL} \\
\typemain{MPI\_DOUBLE\_PRECISION} & \ftype{DOUBLE PRECISION} \\
\typemain{MPI\_COMPLEX} & \ftype{COMPLEX} \\
\typemain{MPI\_LOGICAL} & \ftype{LOGICAL} \\
\typemain{MPI\_CHARACTER} & \ftype{CHARACTER(1)} \\
\type{MPI\_BYTE} & \\
\type{MPI\_PACKED} & \\
\hline
\end{tabular}
\end{center}
\end{table}

\begin{table}[t!]
\caption{Predefined \MPI/ datatypes corresponding to C
 datatypes\mpitermdefindex{basic datatypes!C}}
\label{table:pttopt:datatypes:c}
\begin{center}
\begin{tabular}{|l|l|}
\hline
\textbf{\MPI/ datatype} & \textbf{C datatype} \\
\hline
\typemain{MPI\_CHAR} & \ctype{char} \\
                       &
                         (treated as printable character) \\
\typemain{MPI\_SHORT} & \ctype{signed short int} \\
\typemain{MPI\_INT} & \ctype{signed int} \\
\typemain{MPI\_LONG} & \ctype{signed long int} \\
\typemain{MPI\_LONG\_LONG\_INT}
                       &
                         \ctype{signed long long int} \\
\typemain{MPI\_LONG\_LONG} (as a synonym)
                       &
                         \ctype{signed long long int} \\
\typemain{MPI\_SIGNED\_CHAR} & \ctype{signed char} \\
                       &
                         (treated as integral value) \\
\typemain{MPI\_UNSIGNED\_CHAR} & \ctype{unsigned char} \\
                       &
                         (treated as integral value) \\
\typemain{MPI\_UNSIGNED\_SHORT} & \ctype{unsigned short int} \\
\typemain{MPI\_UNSIGNED} & \ctype{unsigned int} \\
\typemain{MPI\_UNSIGNED\_LONG} & \ctype{unsigned long int} \\
\typemain{MPI\_UNSIGNED\_LONG\_LONG}
                       &
                         \ctype{unsigned long long int} \\
\typemain{MPI\_FLOAT} & \ctype{float} \\
\typemain{MPI\_DOUBLE} & \ctype{double} \\
\typemain{MPI\_LONG\_DOUBLE} & \ctype{long double} \\
\typemain{MPI\_WCHAR}
                       &
                         \ctype{wchar\_t} \\
                       & (defined in \code{<stddef.h>}) \\
                       &
                         (treated as printable character) \\
\typemain{MPI\_C\_BOOL} & \ctype{\_Bool} \\
\typemain{MPI\_INT8\_T} & \ctype{int8\_t} \\
\typemain{MPI\_INT16\_T} & \ctype{int16\_t} \\
\typemain{MPI\_INT32\_T} & \ctype{int32\_t} \\
\typemain{MPI\_INT64\_T} & \ctype{int64\_t} \\
\typemain{MPI\_UINT8\_T} & \ctype{uint8\_t} \\
\typemain{MPI\_UINT16\_T} & \ctype{uint16\_t} \\
\typemain{MPI\_UINT32\_T} & \ctype{uint32\_t} \\
\typemain{MPI\_UINT64\_T} & \ctype{uint64\_t} \\
\typemain{MPI\_C\_COMPLEX} & \ctype{float \_Complex} \\
\typemain{MPI\_C\_FLOAT\_COMPLEX} (as a synonym) & \ctype{float \_Complex} \\
\typemain{MPI\_C\_DOUBLE\_COMPLEX} & \ctype{double \_Complex} \\
\typemain{MPI\_C\_LONG\_DOUBLE\_COMPLEX} & \ctype{long double \_Complex} \\
\type{MPI\_BYTE} & \\
\type{MPI\_PACKED} & \\
\hline
\end{tabular}
\end{center}
\end{table}

\begin{table}[t]
\caption{Predefined \MPI/ datatypes corresponding to both C and Fortran datatypes\mpitermdefindex{basic datatypes!C and Fortran}}
\label{table:pttopt:datatypes:c_f}
\begin{center}
\begin{tabular}{|l|l|l|}
\hline
\textbf{\MPI/ datatype} & \textbf{C datatype} & \textbf{Fortran datatype} \\
\hline
\typemain{MPI\_AINT} & \mpictype{MPI\_Aint} & \mpiftypekind{ADDRESS}\\
\typemain{MPI\_OFFSET} & \mpictype{MPI\_Offset} & \mpiftypekind{OFFSET}\\
\typemain{MPI\_COUNT} & \mpictype{MPI\_Count} & \mpiftypekind{COUNT}\\
\hline
\end{tabular}
\end{center}
\end{table}

The datatypes \type{MPI\_BYTE}\mpitermdefindex{basic datatypes!byte} and \type{MPI\_PACKED}\mpitermdefindex{basic datatypes!packed} do not correspond to a
Fortran or C datatype.  A value of type \mpidtypemain{MPI\_BYTE} consists of a byte
(8 binary digits).  A
byte is uninterpreted and is different from a character.
Different machines may have
different representations for characters, or may use more than one
byte to represent characters.  On the other hand, a byte has the same
binary value on all machines.
The use of the type \type{MPI\_PACKED} is explained in
Section~\ref{sec:pt2pt-packing}.

\MPI/ requires support of
these datatypes,
which match the basic
datatypes of
Fortran and ISO C.
Additional \MPI/ datatypes\mpitermdefindex{basic datatypes!additional host language} should be provided if the host language has
additional datatypes\footnote{These types, such as \code{DOUBLE COMPLEX} and \code{INTEGER*4}, are not specified by any Fortran standard but are extensions commonly accepted by Fortran compilers.}:
\mpidtypemain{MPI\_DOUBLE\_COMPLEX} for double precision complex in
Fortran declared to be of type \ftype{DOUBLE COMPLEX};
\mpidtypemain{MPI\_REAL2},
\mpidtypemain{MPI\_REAL4}, \mpidtypemain{MPI\_REAL8}, and \mpidtypemainshort{MPI\_REAL16} for Fortran reals, declared to be of
type \ftype{REAL*2}, \ftype{REAL*4}, \ftype{REAL*8}, and \ftype{REAL*16}, respectively;
\mpidtypemain{MPI\_INTEGER1}, \mpidtypemain{MPI\_INTEGER2}, \mpidtypemain{MPI\_INTEGER4},
\mpidtypemainshort{MPI\_INTEGER8}, and \mpidtypemain{MPI\_INTEGER16} for Fortran integers, declared to be of type
\ftype{INTEGER*1}, \ftype{INTEGER*2}, \ftype{INTEGER*4}, \ftype{INTEGER*8}, and \ftype{INTEGER*16} respectively;
\mpidtypemain{MPI\_LOGICAL1}, \mpidtypemain{MPI\_LOGICAL2}, \mpidtypemain{MPI\_LOGICAL4},
\mpidtypemainshort{MPI\_LOGICAL8}, and \mpidtypemain{MPI\_LOGICAL16} for Fortran integers, declared to be of type
\ftype{LOGICAL*1}, \ftype{LOGICAL*2}, \ftype{LOGICAL*4}, \ftype{LOGICAL*8}, and \ftype{LOGICAL*16} respectively;
\mpidtypemain{MPI\_COMPLEX4},
\mpidtypemain{MPI\_COMPLEX8},
\mpidtypemain{MPI\_COMPLEX16},
and
\mpidtypemain{MPI\_COMPLEX32}
for complex numbers in Fortran declared to be of type
\ftype{COMPLEX*4},
\ftype{COMPLEX*8},
\ftype{COMPLEX*16},
and
\ftype{COMPLEX*32},
respectively;
etc.

\begin{rationale}
One goal of the design is to allow for \MPI/ to be implemented as a library,
with no need for additional preprocessing or compilation.
Thus, one cannot assume that a communication call has information on the
datatype of variables in the communication buffer; this information must be
supplied by an explicit argument.
The need for such datatype information will become clear in
Section~\ref{subsec:pt2pt-conversion}.
\end{rationale}

The datatypes \type{MPI\_AINT},
\type{MPI\_OFFSET}, and \type{MPI\_COUNT}
correspond to the
\MPI/-defined C types \mpictype{MPI\_Aint},
\mpictype{MPI\_Offset}, and \mpictype{MPI\_Count}
and their Fortran
equivalents \mpiftypekind{ADDRESS},
\mpiftypekind{OFFSET},
and \mpiftypekind{COUNT}.
This is described in Table~\ref{table:pttopt:datatypes:c_f}.
All predefined datatype handles are available in all language bindings.
See
Sections~\ref{subsec:mpiopaqueobjects}
and~\ref{subsec:interlanguage-communication}
on page~\pageref{subsec:mpiopaqueobjects}
and~\pageref{subsec:interlanguage-communication}
for information on interlanguage communication with these types.

If there is an accompanying C++ compiler then the
datatypes in Table~\ref{table:pttopt:datatypes:c++} are also supported in C and Fortran.

\begin{table}[t]
\caption{Predefined \MPI/ datatypes corresponding to C++ datatypes\mpitermdefindex{basic datatypes!C++}}
\label{table:pttopt:datatypes:c++}
\begin{center}
\begin{tabular}{|l|l|}
\hline
\textbf{\MPI/ datatype} & \textbf{C++ datatype} \\
\hline
\mpidtypemain{MPI\_CXX\_BOOL} & \ctype{bool} \\
\mpidtypemain{MPI\_CXX\_FLOAT\_COMPLEX} & \ctype{std::complex$<$float$>$} \\
\mpidtypemain{MPI\_CXX\_DOUBLE\_COMPLEX} & \ctype{std::complex$<$double$>$} \\
\mpidtypemain{MPI\_CXX\_LONG\_DOUBLE\_COMPLEX} & \ctype{std::complex$<$long double$>$} \\
\hline
\end{tabular}
\end{center}
\end{table}


\subsection{Message Envelope}
\mpitermtitleindex{message!envelope}
\mpitermtitleindex{envelope}
\label{subsec:pt2pt-envelope}

In addition to the data part, messages carry information that can be used to
distinguish messages and selectively receive them.  This information consists
of a fixed number of fields, which we collectively call
the \mpitermdefni{message envelope}\mpitermdefindex{envelope}\mpitermdefindex{message!envelope}.  These fields are
% See description of tuple for communicators in context.tex. Use same format
% for consistency.
\begin{obeylines}
\textbf{source}
\textbf{destination}
\textbf{tag}
\textbf{communicator}
\end{obeylines}

The \mpitermni{message source} is implicitly determined by the identity of the
message sender.  The other fields are specified by arguments in the send
procedure.

The \mpitermni{message destination} is specified by the \mpiarg{dest}
argument.

The integer-valued \mpitermni{message tag} is specified by the \mpiarg{tag} argument.
This integer can be used by the program to distinguish different types of
messages.
The range of valid tag values is $0,\ldots,\mpicode{UB}$, where the value of
\mpicode{UB} is
implementation dependent. It can be found by querying the
value of the attribute \mpiconst{MPI\_TAG\_UB}, as
described in Chapter~\ref{chap:environment}.  \MPI/ requires that
\mpicode{UB} be no less than 32767.

The \mpiarg{comm} argument specifies the \mpiterm{communicator} that is used for
the send operation.
Communicators are explained in Chapter~\ref{chap:context}; below is a brief
summary of their usage.

A communicator specifies the
communication context for a communication operation.
Each communication context provides a separate ``communication
universe'': messages are always received within the context they were
sent, and messages sent in different contexts do not interfere.

The communicator also specifies the group of \MPI/ processes that share this
communication context.  This \MPI/ \mpitermni{process group}\mpitermindex{group}
is ordered and \MPI/ processes are identified by their
rank within this group.  Thus, the range of valid values for \mpiarg{dest} is
$0, \ldots, n-1 \cup \{\mpiconst{MPI\_PROC\_NULL}\}$, where $n$ is the number of
\MPI/ processes in the group.  (If the communicator is an inter-communicator\mpitermindex{inter-communicator!point-to-point}, then
destinations are identified by their rank in the remote group.  See
Chapter~\ref{chap:context}.)

An \MPI/ process may have a different rank in each group in which it is a member.

When using the \wpm{} (see Section~\ref{sec:dynamic:initialization}),
a predefined communicator \mpiconst{MPI\_COMM\_WORLD} is provided by \MPI/. It
allows communication with all \MPI/ processes that are accessible after \MPI/
initialization and \MPI/ processes are identified by their rank in the
group of \mpiconst{MPI\_COMM\_WORLD}.

\begin{users}
Users that are comfortable with the notion of a flat name space for
\MPI/ processes, and a single communication context, as offered by most
existing communication libraries, need only use the \wpm{} for \MPI/ initialization, and the predefined
variable \mpiconst{MPI\_COMM\_WORLD} as the \mpiarg{comm} argument.
This will allow communication with all the \MPI/ processes available at
initialization time.

Users may define new communicators, as explained in Chapter~\ref{chap:context}.
Communicators provide an important encapsulation mechanism for libraries and
modules.  They allow modules to have their own disjoint communication universe
and their own \MPI/ process numbering scheme.
\end{users}

\begin{implementors}
The \mpitermni{message envelope}\mpitermindex{envelope}\mpitermindex{message!envelope}
would normally be encoded by a fixed-length message header.
However, the actual encoding
is implementation dependent. Some of the information (e.g.,
source or destination) may be implicit, and need not be
explicitly carried by messages. Also, \MPI/ processes may be identified by
relative ranks, or absolute ids, etc.
\end{implementors}

\subsection{Blocking Receive}
\mpitermtitleindex{receive!blocking}
\mpitermtitleindex{point-to-point communication!receive operation}
\label{subsec:pt2pt-basicreceive}

The syntax of the \mpitermdefni{blocking receive}\mpitermdefindex{receive!blocking} procedure is given below.

\cdeclmainindex{MPI\_Status}\cdeclmainindex{TYPE(MPI\_Status)}
\begin{mpi-binding}
    function_name("MPI_Recv")
    parameter(
        "buf",
        "BUFFER",
        desc=r"initial address of receive buffer",
        direction="out",
    )
    parameter(
        "count",
        "POLYXFER_NUM_ELEM_NNI",
        desc=r"number of elements in receive buffer",
    )
    parameter(
        "datatype", "DATATYPE", desc=r"datatype of each receive buffer element"
    )
    parameter(
        "source", "RANK", desc=r"rank of source or \mpiconst{MPI_ANY_SOURCE}"
    )
    parameter("tag", "TAG", desc=r"message tag or \mpiconst{MPI_ANY_TAG}")
    parameter("comm", "COMMUNICATOR")
    parameter("status", "STATUS", direction="out")
\end{mpi-binding}

The blocking semantics of this call are described in Section~\ref{sec:pt2pt-modes}.

The receive buffer consists of
the storage containing \mpiarg{count} consecutive elements of the type
specified by \mpiarg{datatype}, starting at address \mpiarg{buf}.
The length of the received message must be less than or equal to the length of
the receive buffer. An overflow error\mpitermdefindex{semantics!point-to-point communication!overflow error on trunction} occurs if all incoming data does
not fit, without truncation, into the receive buffer.

If a
message that is shorter than the receive buffer arrives, then only those
locations corresponding to the (shorter) message are modified.

\begin{users}
The \mpifunc{MPI\_PROBE} function described in
Section~\ref{sec:pt2pt-probe} can be used to receive messages of
unknown length.
\end{users}

\begin{implementors}
Even though no specific behavior is mandated by \MPI/ for \mpitermni{erroneous programs}\mpitermindex{erroneous program}, the
recommended handling of overflow situations is to return in
\mpiarg{status} information about the
source and tag of the incoming message.  The receive procedure will return an
error code.  High-quality implementations will also ensure that no memory
that is outside the receive buffer will ever be overwritten.

In the case of a message shorter than the receive buffer, \MPI/ is
quite strict in that it allows no modification of the other locations.
A more lenient statement would allow for some
optimizations but this is not allowed.  The implementation must be ready to
end a copy into the receiver memory exactly at the end of the receive
buffer, even if it is an odd address.
\end{implementors}


The selection of a message by a receive operation is governed by
the value of the \mpitermni{message envelope}\mpitermindex{envelope}\mpitermindex{message!envelope}.
A message can be received by a receive operation
if its \mpiterm{envelope}\mpitermindex{message!envelope} matches\mpitermdefindex{matching rules!envelope} the \mpiarg{source}, \mpiarg{tag} and
\mpiarg{comm} values specified by the
receive operation.  The receiver may specify a \mpitermdef{wildcard}\mpitermdefindex{message!wildcard}\mpitermdefindex{semantics!point-to-point communication!wildcard receive}\mpitermdefindex{matching rules!wildcard}
\mpiconstmain{MPI\_ANY\_SOURCE}
value for \mpiarg{source}, and/or
a wildcard \mpiconstmain{MPI\_ANY\_TAG}
value for \mpiarg{tag},
indicating that any source and/or tag are acceptable.  It cannot specify a
wildcard value for \mpiarg{comm}.
Thus, a message can be received by a receive
operation only if it is addressed
to the receiving \MPI/ process, has a matching communicator, has
matching source unless
\mpiarg{source}\mpicode{ = }\mpiconst{MPI\_ANY\_SOURCE}
in the pattern, and has a matching tag unless
\mpiarg{tag}\mpicode{ = }\mpiconst{MPI\_ANY\_TAG}
in the pattern.

The message tag is specified by the \mpiarg{tag} argument of the
receive operation.
The argument \mpiarg{source}, if different from \mpiconst{MPI\_ANY\_SOURCE},
is specified as a rank within the \MPI/ process group associated
with that same communicator (remote \MPI/ process group, for inter-communicators\mpitermindex{inter-communicator!point-to-point}).
Thus, the range of valid values for the
\mpiarg{source} argument is
\{$0,\ldots,n-1\}\cup\{\mpiconst{MPI\_ANY\_SOURCE}\}\cup\{\mpiconst{MPI\_PROC\_NULL}\}$, where
$n$ is the number of \MPI/ processes in this group.

Note the asymmetry between send and receive operations: A receive
operation may
accept messages from an arbitrary sender, on the other hand, a send operation
must specify a unique receiver.  This matches a ``push'' communication
mechanism, where data transfer is effected by the sender (rather than a
``pull'' mechanism, where data transfer is effected by the receiver).

Source = destination is allowed, that is, an \MPI/ process can send a
message to itself.
However, it is unsafe to do so with the blocking send and receive
operations described above, since this may lead to deadlock\mpitermindex{semantics!deadlock!send to self}. See
Section~\ref{sec:pt2pt-semantics}.

\begin{implementors}
Message context and other communicator information
can be implemented as an additional tag field.  It differs
from the regular message tag in that wild card matching is not allowed on this
field, and that value setting for this field is controlled by communicator
manipulation functions.
\end{implementors}

The use of \mpiarg{dest}\mpicode{ = }\mpiconst{MPI\_PROC\_NULL} or \mpiarg{source}\mpicode{ = }\mpiconst{MPI\_PROC\_NULL} to define a ``dummy''
destination or source in any send or receive call is described in
\sectionref{sec:pt2pt-nullproc}.

\subsection{Return Status}
\mpitermtitleindex{status}
\label{subsec:pt2pt-status}

The source or tag of a received message may not be known if wildcard
values were used in the receive operation.
Also, if multiple requests
are completed by a single \MPI/ function (see
Section~\ref{subsec:pt2pt-multiple}), a distinct error code may need to be
returned for each request.
The information is returned by the \mpiarg{status}
argument of \mpifunc{MPI\_RECV}.  The type of \mpiarg{status} is
\MPI/-defined.
Status variables need to be explicitly
allocated by the user, that is, they are not system objects.


In C, \mpiarg{status}\mpitermdefindex{point-to-point communication!status} is a structure\mpitermdefindex{status!structure in C} that contains three fields named
\mpiconst{MPI\_SOURCE}, \mpiconst{MPI\_TAG}, and \mpiconst{MPI\_ERROR}; the
structure may contain additional fields.  Thus, \mpiarg{status.MPI\_SOURCE},
\mpiarg{status.MPI\_TAG}, and \mpiarg{status.MPI\_ERROR} contain the
source, tag, and error code, respectively, of
the received message.

In Fortran with \code{USE mpi} or (deprecated) \code{INCLUDE 'mpif.h'},
\mpiarg{status} is an array\mpitermdefindex{status!array in Fortran} of \ftype{INTEGER}s of size
\mpiconstmain{MPI\_STATUS\_SIZE}.
The constants \mpiconstmain{MPI\_SOURCE},
\mpiconstmain{MPI\_TAG}, and
\mpiconstmain{MPI\_ERROR}
are the indices of the entries that store the
source, tag, and error fields.  Thus, \mpiarg{status(MPI\_SOURCE)},
\mpiarg{status(MPI\_TAG)}, and \mpiarg{status(MPI\_ERROR)}
contain, respectively, the source,
tag, and error code of the received message.

With Fortran \code{USE} \code{mpi\_f08}, status is defined as the
Fortran \ftype{BIND(C)} derived type\mpitermdefindex{status!derived type in Fortran 2008} \mpiftype{TYPE(MPI\_Status)}\cdeclindex{MPI\_Status}
containing three public \ftype{INTEGER} fields named \mpiconst{MPI\_SOURCE},
\mpiconst{MPI\_TAG}, and \mpiconst{MPI\_ERROR}.
\mpiftype{TYPE(MPI\_Status)} may contain additional, implementation-specific fields.
Thus, \code{status\%MPI\_SOURCE}, \code{status\%MPI\_TAG},
and \code{status\%MPI\_ERROR} contain the source,
tag, and error code of a received message respectively.
Additionally, within both the \code{mpi} and the \code{mpi\_f08} modules,
the constants \mpiconst{MPI\_STATUS\_SIZE},
\mpiconst{MPI\_SOURCE}, \mpiconst{MPI\_TAG}, \mpiconst{MPI\_ERROR},
and \mpiftype{TYPE(MPI\_Status)}\cdeclindex{MPI\_Status}
are defined to allow conversion between both status representations.
Conversion routines are provided in
\sectionref{sec:conversion:status}.

\begin{rationale}
The Fortran \mpiftype{TYPE(MPI\_Status)}\cdeclindex{MPI\_Status}
is defined as a \ftype{BIND(C)} derived type
so that it can be used at any location where
the status integer array representation can be used,
e.g., in user defined common blocks.
\end{rationale}

\begin{rationale}
It is allowed to have the same name (e.g., \mpiconst{MPI\_SOURCE})
defined as a constant (e.g., Fortran parameter) and as a field
of a derived type.
\end{rationale}

In general, message-passing calls do not modify the value of the error
code field of status variables.  This field may be updated only by
the functions in Section~\ref{subsec:pt2pt-multiple} that return multiple
statuses.  The field is updated if and only if such function returns with
an error code of \errormain{MPI\_ERR\_IN\_STATUS}\mpitermindex{status!error in status}.

\begin{rationale}
The error field in status is not needed for calls that return only one
status, such as \mpifunc{MPI\_WAIT}, since that would only duplicate the
information returned by the function itself.  The current design
avoids the additional overhead of setting it, in such cases.  The field
is needed for calls that return multiple statuses, since each request
may have had a different failure.
\end{rationale}

The status argument also returns information on the length\mpitermindex{status!message length} of the message
received.  However, this information is not directly available as a field
of the status variable and a call to \mpifunc{MPI\_GET\_COUNT} is required
to ``decode'' this information.

\cdeclindex{MPI\_Status}
\begin{mpi-binding}
    function_name("MPI_Get_count")
    parameter(
        "status",
        "STATUS",
        desc=r"return status of receive operation",
        constant=True,
    )
    parameter(
        "datatype", "DATATYPE", desc=r"datatype of each receive buffer entry"
    )

    # MR/DH: this is intentionally left as XFER_NUM_ELEM_SMALL, not _NNI, because
    # that does not replicate the MPI 3.1. count may be negative with MPI_UNDEFINED.
    parameter(
        "count",
        "POLYXFER_NUM_ELEM",
        desc=r"number of received entries",
        direction="out",
    )
\end{mpi-binding}

Returns the number of entries received.  (Again, we count \emph{entries}, each of
type \mpiarg{datatype}, not \emph{bytes}.)
The \mpiarg{datatype} argument should match
the argument provided by the receive call that set the \mpiarg{status} variable.
If the number of entries received exceeds the limits of the \mpiarg{count}
parameter, then \mpifunc{MPI\_GET\_COUNT} sets the value of \mpiarg{count} to
\mpiconst{MPI\_UNDEFINED}.
There are other situations where the value of
\mpiarg{count} can be set to \mpiconst{MPI\_UNDEFINED}; see Section~\ref{subsec:pt2pt-datatypeuse}.

\begin{rationale}
Some message-passing libraries use \INOUT\ \mpiarg{count},
\mpiarg{tag} and \mpishortarg{source} arguments, thus using them both to specify
the selection criteria for incoming messages and return the actual \mpiterm{envelope}\mpitermindex{message!envelope}
values of the received message.
The use of a separate status argument prevents errors that are often attached
with \INOUT\ argument (e.g., using the
\mpiconst{MPI\_ANY\_TAG} constant as
the tag in a receive).
Some libraries use calls that refer implicitly to the ``last message
received.''  This is not thread safe.

The \mpiarg{datatype} argument is passed to \mpifunc{MPI\_GET\_COUNT} so as to
improve performance. A message might be received without counting the number
of elements it contains, and the count value is often not needed.
Also, this allows the same function to be used after a call to
\mpifunc{MPI\_PROBE} or \mpifunc{MPI\_IPROBE}. With a status from \mpifunc{MPI\_PROBE} or \mpifunc{MPI\_IPROBE},
the same datatypes are allowed as in a call to \mpifunc{MPI\_RECV} to receive this message.
\end{rationale}

%  3.2.11 MPI\_GET\_COUNT with zero-length datatypes
The value returned as the \mpiarg{count} argument of
\mpifunc{MPI\_GET\_COUNT} for a datatype of length zero where zero bytes
have been transferred is zero.  If the number of bytes transferred is
greater than zero, \mpiconst{MPI\_UNDEFINED} is returned.

\begin{rationale}
Zero-length datatypes may be created in a number of cases.
An important case is
\mpifunc{MPI\_TYPE\_CREATE\_DARRAY}, where the
definition of the particular
darray
results in an empty block on some
\MPI/ process.  Programs written in an SPMD style will not check for
this special case and may want to use \mpifunc{MPI\_GET\_COUNT} to check
the status.
\end{rationale}

\begin{users}
The buffer size required for the receive can be affected by data conversions and
by the stride of the receive datatype.
In most cases, the safest approach is to use the same datatype with \mpifunc{MPI\_GET\_COUNT} and the receive.
\end{users}

All send and receive operations use the \mpiarg{buf}, \mpiarg{count},
\mpiarg{datatype}, \mpiarg{source}, \mpiarg{dest}, \mpiarg{tag},
\mpiarg{comm}, and \mpiarg{status} arguments in the same way as the blocking
\mpifunc{MPI\_SEND} and
\mpifunc{MPI\_RECV} procedures described in this section.


While the \mpiconst{MPI\_SOURCE}, \mpiconst{MPI\_TAG}, and \mpiconst{MPI\_ERROR}
status values are directly accessible by the user,
for convenience in some contexts, users can also access
them via procedure calls, as described below.

\cdeclindex{MPI\_Status}%
\begin{mpi-binding}
    function_name("MPI_Status_get_source")
    parameter(
        "status",
        "STATUS",
        direction="in", constant=True,
        desc=r"status from which to retrieve source rank",
    )
    parameter(
        "source",
        "RANK",
        desc=r"rank set in the \mpiconst{MPI\_SOURCE} field",
        direction="out",
    )
\end{mpi-binding}

Returns in \mpiarg{source} the value of the \mpiconst{MPI\_SOURCE} field in the \mpiarg{status} object.

\cdeclindex{MPI\_Status}%
\begin{mpi-binding}
    function_name("MPI_Status_get_tag")
    parameter(
        "status",
        "STATUS",
        direction="in", constant=True,
        desc=r"status from which to retrieve tag",
    )
    parameter(
        "tag",
        "TAG",
        desc=r"tag set in the \mpiconst{MPI\_TAG} field",
        direction="out",
    )
\end{mpi-binding}

Returns in \mpiarg{tag} the value in the \mpiconst{MPI\_TAG} field of the \mpiarg{status} object.


\cdeclindex{MPI\_Status}%
\begin{mpi-binding}
    function_name("MPI_Status_get_error")
    parameter(
        "status",
        "STATUS",
        direction="in", constant=True,
        desc=r"status from which to retrieve error",
    )
    parameter(
        "err",
        "ERROR_CODE",
        desc=r"error set in the \mpiconst{MPI\_ERROR} field",
        direction="out",
    )
\end{mpi-binding}

Returns in \mpiarg{err} the value in the \mpiconst{MPI\_ERROR} field of the \mpiarg{status} object.

Procedures for setting these fields in a status object are defined in Section~\ref{sec:ei-status}.

\subsection{Passing \texorpdfstring{\mpiconst{MPI\_STATUS\_IGNORE}}{MPI\_STATUS\_IGNORE} for Status}
\mpitermtitleindex{status!ignore}
\label{sec:pt2pt-status-ignore}

Every call to \mpifunc{MPI\_RECV} includes a \mpiarg{status} argument, wherein
the system can return details about the message received.
There are also a number of other \MPI/ calls
where \mpiarg{status} is returned.
An object of type
\mpictype{MPI\_Status} 
is not an \mpi/ opaque object; its structure is declared
in \code{mpi.h} and (deprecated) \code{mpif.h}, and it exists in the user's program.  In many
cases, application programs are constructed so that it is unnecessary for them
to examine the \mpiarg{status} fields.  In these cases, it is a waste for the user
to allocate a status object, and it is particularly wasteful for the \mpi/
implementation to fill in fields in this object.


To cope with this problem, there are two predefined constants,
\mpiconstmain{MPI\_STATUS\_IGNORE} and
\mpiconstmain{MPI\_STATUSES\_IGNORE},
which when
passed to a receive,
probe,
wait, or test function, inform the implementation that
the status fields are not to be filled in.  Note that
\mpiconst{MPI\_STATUS\_IGNORE} is not a special type of \mpictype{MPI\_Status}
object; rather, it is a special value for the argument.  In C one would expect
it to be \code{NULL}, not the address of a special \mpictype{MPI\_Status}.

\mpiconst{MPI\_STATUS\_IGNORE}, and the array version
\mpiconst{MPI\_STATUSES\_IGNORE}, can be used everywhere a status argument is
passed to a receive, wait, or test function.  \mpiconst{MPI\_STATUS\_IGNORE}
cannot be used when status is an \IN\ argument.
Note that in Fortran \mpiconst{MPI\_STATUS\_IGNORE} and \mpiconst{MPI\_STATUSES\_IGNORE} are
objects like \mpiconst{MPI\_BOTTOM} (not usable for initialization or
assignment), see Section~\ref{subsec:named-constants}.




In general, this optimization can apply to all functions for which
\mpiarg{status} or an array of \mpiarg{status}es is an \OUT\ argument.
Note that this converts \mpiarg{status} into an \INOUT\ argument.
The functions that can be passed \mpiconst{MPI\_STATUS\_IGNORE} are all the various forms of
\mpifunc{MPI\_RECV},
\mpifunc{MPI\_PROBE},
\mpifunc{MPI\_TEST}, and \mpifunc{MPI\_WAIT}, as well as
\mpifunc{MPI\_REQUEST\_GET\_STATUS}.
When an array is passed, as in the
\mpiskipfunc{\mpicode{MPI\_}\{\mpicode{TEST$|$WAIT}\}\{\mpicode{ALL$|$SOME}\}}
functions, a separate constant, \mpiconst{MPI\_STATUSES\_IGNORE}, is passed for
the array argument.
It is possible for an \MPI/ function to return \error{MPI\_ERR\_IN\_STATUS}
even when \mpiconst{MPI\_STATUS\_IGNORE} or \mpiconst{MPI\_STATUSES\_IGNORE} has
been passed to that function.



\mpiconst{MPI\_STATUS\_IGNORE} and \mpiconst{MPI\_STATUSES\_IGNORE} are not
required to have the same values in C and Fortran.

It is not allowed to have some of the statuses in an array of statuses for
\mpiskipfunc{\mpicode{MPI\_}\{\mpicode{TEST$|$WAIT}\}\{\mpicode{ALL$|$SOME}\}}
functions set to
\mpiconst{MPI\_STATUS\_IGNORE}; one either specifies ignoring \emph{all} of
the statuses in such a call with \mpiconst{MPI\_STATUSES\_IGNORE}, or \emph{none}
of them by passing normal statuses in all positions in the array
of statuses.

\subsection{Blocking Send-Receive}
\mpitermtitleindex{send-receive!blocking}
\mpitermtitleindex{point-to-point communication!send-receive operation}
\label{sec:pt2pt-sendrecv}

The \mpitermdefni{send-receive}\mpitermdefindex{send-receive!blocking} operations combine in one operation the sending of a
message to one destination and the receiving of another message, from
another \MPI/ process.  The two (source and destination) are possibly the same.
A send-receive operation is
very useful for executing a shift operation across a chain of
\MPI/ processes.  If blocking sends and receives are used for such a shift,
then one needs to order the sends and receives correctly (for
example, \MPI/ processes with even rank in the communicator
send, then receive, \MPI/ processes with odd rank in the communicator receive first, then send) so as to prevent
cyclic dependencies that may lead to \mpitermni{deadlock}\mpitermindex{semantics!deadlock!cyclic shift}.  When a send-receive
operation is used, the communication subsystem takes care of
these issues\mpitermindex{deadlock avoidance!cyclic shift}.  The send-receive operation can be used in conjunction
with the procedures described in Chapter~\ref{chap:topol} in order to
perform shifts on various logical topologies.
Also, a send-receive operation is useful for implementing
remote procedure calls.


A message sent by a
send-receive operation can be received by a regular receive
operation or probed by a probe operation; a send-receive operation can
receive a message sent by a regular send operation.


\cdeclindex{MPI\_Status}
\begin{mpi-binding}
    function_name("MPI_Sendrecv")
    parameter(
        "sendbuf",
        "BUFFER",
        desc=r"initial address of send buffer",
        constant=True,
    )
    parameter(
        "sendcount",
        "POLYXFER_NUM_ELEM_NNI",
        desc=r"number of elements in send buffer",
    )
    parameter("sendtype", "DATATYPE", desc=r"type of elements in send buffer")
    parameter("dest", "RANK", desc=r"rank of destination")
    parameter("sendtag", "TAG", desc=r"send tag")

    parameter(
        "recvbuf",
        "BUFFER",
        desc=r"initial address of receive buffer",
        direction="out",
    )
    parameter(
        "recvcount",
        "POLYXFER_NUM_ELEM_NNI",
        desc=r"number of elements in receive buffer",
    )
    parameter(
        "recvtype", "DATATYPE", desc=r"type of elements receive buffer element"
    )
    parameter(
        "source", "RANK", desc=r"rank of source or \mpiconst{MPI_ANY_SOURCE}"
    )
    parameter("recvtag", "TAG", desc=r"receive tag or \mpiconst{MPI_ANY_TAG}")
    parameter("comm", "COMMUNICATOR")
    parameter("status", "STATUS", direction="out")
\end{mpi-binding}

Execute a blocking send-receive operation.  Both send and receive
use the same communicator, but
possibly different tags.  The send buffer and receive buffers must be
disjoint, and may have different lengths and datatypes.

The semantics of a send-receive operation\mpitermdefindex{semantics!point-to-point communication!send-receive concurrency} is what would be obtained
if the caller forked two concurrent threads, one to execute the send,
and one to execute the receive, followed by a join of these two
threads.

\cdeclindex{MPI\_Status}
\begin{mpi-binding}
    function_name("MPI_Sendrecv_replace")
    parameter(
        "buf",
        "BUFFER",
        desc=r"initial address of send and receive buffer",
        direction="inout",
    )
    parameter(
        "count",
        "POLYXFER_NUM_ELEM_NNI",
        desc=r"number of elements in send and receive buffer",
    )
    parameter(
        "datatype",
        "DATATYPE",
        desc=r"type of elements in send and receive buffer",
    )
    parameter("dest", "RANK", desc=r"rank of destination")
    parameter("sendtag", "TAG", desc=r"send message tag")
    parameter(
        "source", "RANK", desc=r"rank of source or \mpiconst{MPI_ANY_SOURCE}"
    )
    parameter(
        "recvtag", "TAG", desc=r"receive message tag or \mpiconst{MPI_ANY_TAG}"
    )
    parameter("comm", "COMMUNICATOR")
    parameter("status", "STATUS", direction="out")
\end{mpi-binding}

Execute a blocking send and receive. The same buffer is used both for
the send and for the receive, so that the message sent is replaced by
the message received.

\begin{implementors}
Additional intermediate buffering\mpitermindex{message!intermediate buffering} is needed for the
``replace'' variant.
\end{implementors}

\section{Datatype Matching and Data Conversion}
\label{sec:pt2pt-typematch}

\subsection{Type Matching Rules}
\mpitermtitleindex{type!matching rules}
\mpitermtitleindex{matching rules!type}
\label{subsec:pt2pt-typematch}

One can think of message transfer as consisting of the following three phases.
\begin{enumerate}
\item
Data is pulled out of the send buffer and a message is assembled.
\item
A message is transferred from sender to receiver.
\item
Data is pulled from the incoming message and disassembled into the receive
buffer.
\end{enumerate}

\mpitermdefni{Type matching}\mpitermdefindex{matching rules!type} has to be observed at each of these three phases: The type
of each variable in the sender buffer has to match
the type specified for that entry by the send operation;
the type specified by the send operation has to match the type specified by
the receive operation; and
the type of each
variable in the receive buffer has to match the type specified for
that entry by the receive operation.
A program that fails to observe these
three rules is \mpitermni{erroneous}\mpitermindex{erroneous program!type matching}\mpitermindex{semantics!erroneous program!type matching}.

To define type matching more precisely, we need to deal with two issues:
matching of types of the host language with types specified in
communication operations;
and matching of types at sender and receiver.

The types of a send and receive match (phase two) if
both operations use identical names. That is, \type{MPI\_INTEGER}
matches \type{MPI\_INTEGER}, \type{MPI\_REAL} matches \type{MPI\_REAL},
and so on.
There is one exception to this rule, discussed in
Section~\ref{sec:pt2pt-packing}: the type \type{MPI\_PACKED} can match
any other type.

The type of a variable in a host program matches the type specified in the
communication operation
if the datatype name used by that operation corresponds
to the basic type of the host program variable.  For example, an entry with type
name \type{MPI\_INTEGER} matches a Fortran variable of type \ftype{INTEGER}.
A table giving this correspondence for Fortran and C appears in
Section~\ref{subsec:pt2pt-messagedata}.
There are two exceptions to this last rule: an entry with type name
\type{MPI\_BYTE} or \type{MPI\_PACKED} can be used to match
any byte of storage (on a byte-addressable machine),
irrespective of the datatype of the variable that contains this byte.
The type \type{MPI\_PACKED} is used to send data that has been
explicitly packed, or receive data that will be explicitly unpacked,
see Section~\ref{sec:pt2pt-packing}.
The type \type{MPI\_BYTE} allows one to transfer the binary value of a byte in
memory unchanged.

To summarize, the type matching rules fall into the three categories below\mpitermdefindex{communication!typed, untyped, and packed}.
\begin{itemize}
\item
Communication of typed values (e.g., with datatype different from
\type{MPI\_BYTE}),
where the datatypes of the corresponding entries in the sender
program, in the send call,
in the receive call and in the receiver program
must all match.
\item
Communication of untyped values (e.g., of datatype \type{MPI\_BYTE}), where
both sender and receiver use the
datatype \type{MPI\_BYTE}.  In this case, there are no
requirements on the types of the corresponding entries in the sender and the
receiver programs, nor is it required that they be the same.
\item
Communication involving packed data, where \type{MPI\_PACKED} is used.
\end{itemize}


The following examples illustrate the first two cases.
\label{page:pt2pt-3examples}

\begin{example}
\label{pt2pt-exA}
\Exindex{Fortran}{Datatype matching}{MPI\_Send,MPI\_Recv}%
\exindex{Datatypes!matching}%
Sender and receiver specify matching types.
%%HEADER
%%LANG: FORTRAN
%%FRAGMENT
%%DECL: integer comm, rank, ierr, tag, a(2), b(2)
%%DECL: integer status(MPI_STATUS_SIZE)
%%ENDHEADER
\begin{lstlisting}[language={[MPI]Fortran}]
CALL MPI_COMM_RANK(comm, rank, ierr)
IF (rank .EQ. 0) THEN
   CALL MPI_SEND(a(1), 10, MPI_REAL, 1, tag, comm, ierr)
ELSE IF (rank .EQ. 1) THEN
   CALL MPI_RECV(b(1), 15, MPI_REAL, 0, tag, comm, status, ierr)
END IF
\end{lstlisting}

This code is correct if both \code{a} and \code{b} are real arrays of
size $\ge 10$. (In Fortran, it might be correct to use this code
even if \code{a} or \code{b} have size $< 10$: e.g., when \code{a(1)} can
be equivalenced to an array with ten reals.)
\end{example}

\begin{example}
\label{pt2pt-exB}
\Exindex[Datatype matching (erroneous)]{Fortran}{Datatype matching}{MPI\_Send,MPI\_Recv}%
\exindex{Datatypes!not matching}%
Sender and receiver do not specify matching types.

%%HEADER
%%LANG: FORTRAN
%%FRAGMENT
%%DECL: integer comm, rank, ierr, tag, a(2), b(2)
%%DECL: integer status(MPI_STATUS_SIZE)
%%ENDHEADER
\begin{lstlisting}[language={[MPI]Fortran}]
! ---------------- THIS EXAMPLE IS ERRONEOUS ---------------
CALL MPI_COMM_RANK(comm, rank, ierr)
IF (rank .EQ. 0) THEN
   CALL MPI_SEND(a(1), 10, MPI_REAL, 1, tag, comm, ierr)
ELSE IF (rank .EQ. 1) THEN
   CALL MPI_RECV(b(1), 40, MPI_BYTE, 0, tag, comm, status, ierr)
END IF
\end{lstlisting}

This code is \mpitermni{erroneous}, since sender and receiver do not provide
matching datatype arguments.
\end{example}

\begin{example}
\label{pt2pt-exC}
\Exindex[Using \mpiconst{MPI\_BYTE}]{Fortran}{Datatypes}{MPI\_Send,MPI\_Recv,c:MPI\_BYTE}%
\exindex{Datatypes!untyped}%
Sender and receiver specify communication of untyped values.
%%HEADER
%%LANG: FORTRAN
%%FRAGMENT
%%DECL: integer comm, rank, ierr, tag, a(2), b(2)
%%DECL: integer status(MPI_STATUS_SIZE)
%%ENDHEADER
\begin{lstlisting}[language={[MPI]Fortran}]
CALL MPI_COMM_RANK(comm, rank, ierr)
IF (rank .EQ. 0) THEN
   CALL MPI_SEND(a(1), 40, MPI_BYTE, 1, tag, comm, ierr)
ELSE IF (rank .EQ. 1) THEN
   CALL MPI_RECV(b(1), 60, MPI_BYTE, 0, tag, comm, status, ierr)
END IF
\end{lstlisting}

This code is correct, irrespective of the type and size of \code{a} and
\code{b} (unless this results in an out of bounds memory access).
\end{example}

\begin{users}
If a buffer of type \type{MPI\_BYTE} is passed as an argument to
\mpifunc{MPI\_SEND}, then \MPI/ will send the data stored at contiguous
locations, starting from the address
indicated by the \mpiarg{buf} argument.  This may have unexpected results when
the data layout is not as a casual user would expect it to be.
For example, some Fortran compilers implement variables of type
\ftype{CHARACTER}
as a structure that contains the character length and a pointer to
the actual string.  In such an environment, sending and receiving a Fortran
\ftype{CHARACTER} variable using the \type{MPI\_BYTE} type will not
have the anticipated result of transferring the character string.
For this reason, the user is advised to use typed communication operations
whenever possible.
\end{users}

\subsubsection{Type \texorpdfstring{\type{MPI\_CHARACTER}}{MPI\_CHARACTER}}

The
type \type{MPI\_CHARACTER} matches one character of a Fortran variable of
type \ftype{CHARACTER}, rather than the entire character string stored in the
variable.  Fortran variables of type \ftype{CHARACTER}
or substrings are transferred as if they were arrays of characters.
This is illustrated in the example below.

\begin{example}
\label{pt2pt-exD}
\Exindex[Fortran \ftype{CHARACTER}]{Fortran}{Fortran CHARACTER@Fortran \ftype{CHARACTER}}{MPI\_Send,MPI\_Recv,c:MPI\_CHARACTER}%
Transfer of Fortran \ftype{CHARACTER}s.

%%HEADER
%%LANG: FORTRAN
%%FRAGMENT
%%DECL: integer comm, rank, ierr, tag
%%DECL: integer status(MPI_STATUS_SIZE)
%%ENDHEADER
\begin{lstlisting}[language={[MPI]Fortran}]
CHARACTER*10 a
CHARACTER*10 b

CALL MPI_COMM_RANK(comm, rank, ierr)
IF (rank .EQ. 0) THEN
   CALL MPI_SEND(a, 5, MPI_CHARACTER, 1, tag, comm, ierr)
ELSE IF (rank .EQ. 1) THEN
   CALL MPI_RECV(b(6:10), 5, MPI_CHARACTER, 0, tag, comm, status, ierr)
END IF
\end{lstlisting}

The last five characters of string \code{b} at the \MPI/ process with \code{rank = 1} are replaced by the
first five characters of string \code{a} at the \MPI/ process with \code{rank = 0}.
\end{example}

\begin{rationale}
The alternative choice would be for \type{MPI\_CHARACTER} to
match a character of arbitrary length.  This runs into problems.

A Fortran character variable is a constant length string, with no special
termination symbol.  There is no fixed convention on how to represent
characters, and how to store their length.
Some compilers pass a character argument to
a routine as a pair of arguments, one holding the address of the string and the
other holding the length of string.  Consider the case of an \MPI/
communication call that is passed a communication buffer
with type defined by a derived datatype (Section~\ref{sec:pt2pt-datatype}).
If this communicator buffer contains variables of type \ftype{CHARACTER} then
the information on their length will not be passed to the \MPI/ routine.

This problem forces us to provide explicit information on character
length with the \MPI/ call.  One could add a length parameter to the
type \type{MPI\_CHARACTER}, but this does not add much convenience and the same
functionality can be achieved by defining a suitable derived datatype.
\end{rationale}

\begin{implementors}
Some compilers pass Fortran \ftype{CHARACTER} arguments as a structure with a
length and a pointer to the actual string.  In such an environment, the \MPI/
call needs to dereference the pointer in order to reach the string.
\end{implementors}


\subsection{Data Conversion}
\mpitermtitleindex{data conversion}
\mpitermtitleindex{conversion}
\mpitermtitleindex{semantics!data conversion}
\label{subsec:pt2pt-conversion}

One of the goals of \MPI/ is to support parallel computations across
heterogeneous environments.  Communication in a heterogeneous
environment may require data conversions.
We use the following terminology.
\begin{description}
\item[type conversion]\mpitermdefindex{type!conversion}\mpitermdefindex{conversion!type} changes the datatype of a value, e.g., by rounding a
\ftype{REAL} to an \ftype{INTEGER}.
\item[representation conversion]\mpitermdefindex{representation!conversion}\mpitermdefindex{conversion!representation} changes the binary representation of a value,
e.g., from Hex floating point to IEEE floating point.
\end{description}

The type matching rules imply that \MPI/ communication never entails type
conversion.  On the other hand, \MPI/ requires that a representation
conversion be
performed when a typed value is transferred across environments that use
different representations for the datatype of this value.
\MPI/ does not specify rules for representation conversion.  Such conversion is
expected to preserve integer, logical and character values, and to
convert a floating point value to the nearest value that can be
represented on the target system.

Overflow and underflow exceptions may occur during floating point conversions.
Conversion of integers or characters may also lead to exceptions when a value
that can be represented in one system cannot be represented in the other
system.  An exception
occurring during representation conversion results in a failure of the
communication. An error occurs either in the send operation, or the receive
operation, or both.

If a value sent in a message is untyped (i.e., of type \type{MPI\_BYTE}),
then the binary representation of the byte stored at the receiver is
identical to the binary representation of the byte loaded at the sender.  This
holds true, whether sender and receiver run in the same or in distinct
environments. No representation conversion is required.
(Note that representation conversion may
occur when values of type \type{MPI\_CHARACTER} or
\type{MPI\_CHAR} are transferred, for example, from an
EBCDIC encoding to an ASCII encoding.)

No conversion need occur when an \MPI/ program executes in
a homogeneous system, where all \MPI/ processes run in the same environment.


Consider the three examples, \ref{pt2pt-exA}--\ref{pt2pt-exC}.
The first program is correct, assuming that \code{a} and \code{b} are
\ftype{REAL} arrays of size $\ge 10$.
If the sender and receiver execute in different environments,
then the ten real values that are fetched from the send buffer will
be converted to the representation for reals on the receiver site
before they are stored in the receive buffer.  While the
number of real elements fetched from the send buffer equal the
number of real elements stored in the receive buffer, the number of
bytes stored need not equal the number of bytes loaded. For example, the
sender may use a four byte representation and the receiver an eight
byte representation for reals.

The second program is \mpitermni{erroneous}, and its behavior is undefined.

The third program is correct.  The exact same
sequence of forty bytes that were loaded from the send buffer will be
stored in the receive buffer, even
if sender and receiver run in a different environment.  The message
sent has exactly the same length (in bytes) and the same binary
representation as the message received.  If \code{a} and \code{b} are
of different types,
or if they are of the same type but different data representations are used,
then the bits stored in the receive buffer may encode values that are
different from the values they encoded in the send buffer.


Data representation conversion also applies to the \mpiterm{envelope}\mpitermindex{message!envelope}\mpitermindex{envelope!data representation conversion} of a message:
source, destination and tag are all integers that may need to be converted.

\begin{implementors}
The current definition does not require messages to carry data
type information.  Both sender and receiver provide complete data
type information.
In a heterogeneous environment, one can either use
a machine independent encoding
such as \mpitermindex{XDR}XDR, or have the receiver convert from the sender representation to its
own, or even have the sender do the conversion.

Additional type information might be added to messages
in order to allow the system to detect mismatches between
datatype at sender and receiver.
This might be particularly useful in a slower but safer debug mode.
\end{implementors}

\MPI/ requires support for inter-language communication, e.g.,
if messages are sent
using an \MPI/ procedure from the \MPI/ C language interface
and received
using an \MPI/ procedure from one of the \MPI/ Fortran language interfaces.
The behavior is defined in
\sectionref{sec:misc-lang-interop}.


\section{Communication Modes}
\mpitermtitleindexmainsub{communication}{modes}
\mpitermtitleindex{point-to-point communication!send modes}
\mpitermtitleindex{semantics!point-to-point communication!send modes}
\label{sec:pt2pt-modes}

The send call described in Section~\ref{subsec:pt2pt-basicsend}
is \mpiterm{blocking}:
it does not return until the \mpitermni{message data}\mpitermindex{message!data}
and \mpiterm{envelope}\mpitermindex{message!envelope} have been safely stored away so that the sender is
free to modify
the send buffer.
The message might be copied directly into the matching receive buffer,
or it might be copied into a temporary system buffer.

Message buffering\mpitermindex{message!intermediate buffering}
decouples the send and receive operations.  A blocking send can
complete as soon
as the message was buffered, even if no matching receive has been executed
by the receiver.  On the other hand, message buffering can be expensive,
as it entails additional memory-to-memory copying, and it requires the
allocation of memory for buffering.  \MPI/ offers the choice of
several \mpitermdef{communication modes} that allow one to control the choice of the
communication protocol.

The send call described in Section~\ref{subsec:pt2pt-basicsend}
uses
the \mpitermdefni{standard}\mpitermdefindex{modes!standard}\mpitermdefindex{standard mode send} communication mode.  In this mode,
it is up to \MPI/ to decide whether outgoing
messages will be buffered.  \MPI/ may
buffer outgoing messages.  In such a case, the send call may complete
before a matching receive is invoked.  On the other hand, buffer space may be
unavailable, or \MPI/ may choose not to buffer
outgoing messages, for performance reasons. In this case,
the send call will not complete until a matching receive has been \mpitermni{started}, and
the data has been moved to the receiver.

Thus, a \mpiterm{standard mode send} can be \mpitermni{started}\mpitermindex{started!request objects}\mpitermindex{request objects!started} whether or not a
matching receive has been \mpitermni{started}.  It may \mpitermni{complete}\mpitermindex{completion} before a matching receive
is \mpitermni{started}.  The
standard mode send is \mpiterm{nonlocal}: successful completion of the send
operation may depend on the occurrence of a matching receive.

\begin{rationale}
The reluctance of \MPI/ to mandate whether standard sends are buffering
or not stems from the desire to achieve portable programs.  Since any
system will run out of buffer resources as message sizes are increased,
and some implementations may want to provide little buffering, \MPI/
takes the position that correct (and therefore, portable) programs
do not rely on system buffering in standard mode.  Buffering
may improve the performance
of a correct program, but it doesn't affect the result of the program.
If the user wishes to guarantee a certain amount of buffering, the
user-provided buffer system of Section~\ref{sec:pt2pt-buffer} should be used,
along with the buffered-mode send.
\end{rationale}

There are three additional communication modes.

A \mpitermdefni{buffered}\mpitermdefindex{modes!buffered}\mpitermdefindex{buffered mode send} mode send operation can be started whether or not a
matching receive has been \mpitermni{started}.
It may complete before a matching receive is \mpitermni{started}.  However, unlike
the standard send, this operation is \mpiterm{local}, and its
completion does not depend on the occurrence of a matching receive.  Thus, if a
send is executed and no matching receive is \mpitermni{started}, then \MPI/ must buffer the
outgoing message, so as to allow the send call to complete.  An error will
occur if there is insufficient buffer space.  The amount of available buffer
space is controlled by the user---see Section~\ref{sec:pt2pt-buffer}.
Buffer allocation by the user may be required for the buffered mode to be
effective.

A send that uses the \mpitermdefni{synchronous}\mpitermdefindex{modes!synchronous}\mpitermdefindex{synchronous mode send} mode can be started whether or
not a matching receive was \mpitermni{started}.  However, the send will complete
successfully only if a matching receive is \mpitermni{started}, and the
receive operation has started to receive the message sent by the
synchronous send.
Thus, the completion of a synchronous send not only indicates that the
send buffer can be reused, but
it
also indicates that the receiver has
reached a certain point in its execution, namely that it has started
executing the matching receive.  If both sends and receives are
blocking operations then the use of the synchronous mode provides
synchronous communication semantics: a communication does not complete
at either end before both \MPI/ processes rendezvous at the
communication.  A send executed in this mode is \mpiterm{nonlocal}.

A send that uses the \mpitermdefni{ready}\mpitermdefindex{modes!ready}\mpitermdefindex{ready mode send} communication mode
may be started \emph{only} if the matching receive is already \mpitermni{started}.
Otherwise, the operation is \mpitermni{erroneous}\mpitermindex{erroneous program!ready mode before receive}\mpitermindex{semantics!erroneous program!ready mode before receive} and its outcome is undefined.
On some systems, this allows the removal of a hand-shake
protocol that is otherwise required and results in improved
performance.
The completion of the send operation does not depend on the
status of a matching receive, and merely indicates that the send
buffer can be reused.  A send operation that uses the ready mode has
the same semantics as a standard send operation, or a synchronous send
operation; it is merely that the sender provides additional
information to the system (namely that a matching receive is already
\mpitermni{started}), that can save some overhead.  In a correct program, therefore, a
ready send could be replaced by a standard send with no effect on the
behavior of the program other than performance.

Three additional send functions are provided for the three additional
communication
modes.  The communication mode is indicated by a one letter prefix:
\mpicode{B} for buffered,
\mpicode{S} for synchronous, and
\mpicode{R} for ready.

\begin{mpi-binding}
    function_name("MPI_Bsend")
    parameter(
        "buf", "BUFFER", desc=r"initial address of send buffer", constant=True
    )
    parameter(
        "count",
        "POLYXFER_NUM_ELEM_NNI",
        desc=r"number of elements in send buffer",
    )
    parameter(
        "datatype", "DATATYPE", desc=r"datatype of each send buffer element"
    )
    parameter("dest", "RANK", desc=r"rank of destination")
    parameter("tag", "TAG", desc=r"message tag")
    parameter("comm", "COMMUNICATOR")
\end{mpi-binding}

Send in buffered mode.

According to the definitions in Section~\ref{subsec:procedures},
\mpifunc{MPI\_BSEND} is a completing procedure and the user can re-use
all resources given as arguments, including the \mpitermni{message data buffer}\mpitermindex{message!data}\mpitermindex{message!buffer}.
It is also a local procedure because it returns immediately without
depending on the execution of any \MPI/ procedure in any other \MPI/ process.
\begin{users}
This is one of the exceptions\mpitermindex{semantics!exceptions!completing and local} in which a completing and therefore blocking operation-related procedure is local.
\end{users}

\begin{mpi-binding}
    function_name("MPI_Ssend")
    parameter(
        "buf", "BUFFER", desc=r"initial address of send buffer", constant=True
    )
    parameter(
        "count",
        "POLYXFER_NUM_ELEM_NNI",
        desc=r"number of elements in send buffer",
    )
    parameter(
        "datatype", "DATATYPE", desc=r"datatype of each send buffer element"
    )
    parameter("dest", "RANK", desc=r"rank of destination")
    parameter("tag", "TAG", desc=r"message tag")
    parameter("comm", "COMMUNICATOR")
\end{mpi-binding}

Send in synchronous mode.

\begin{mpi-binding}
    function_name("MPI_Rsend")
    parameter(
        "buf", "BUFFER", desc=r"initial address of send buffer", constant=True
    )
    parameter(
        "count",
        "POLYXFER_NUM_ELEM_NNI",
        desc=r"number of elements in send buffer",
    )
    parameter(
        "datatype", "DATATYPE", desc=r"datatype of each send buffer element"
    )
    parameter("dest", "RANK", desc=r"rank of destination")
    parameter("tag", "TAG", desc=r"message tag")
    parameter("comm", "COMMUNICATOR")
\end{mpi-binding}

Send in ready mode.

There is only one receive operation,
but it matches\mpitermdefindex{matching rules!send modes}
any of the send modes.
The receive procedure described in the last section is \mpiterm{blocking}:
it returns only after the receive buffer contains the newly received
message.  A receive can complete before the matching send
has completed (of course, it can complete only after the matching send
has started).

In a multithreaded implementation of \MPI/, the system may de-schedule a
thread that is blocked on a send or receive
operation, and schedule another thread for execution in the same
address space.
In such a case it is the user's responsibility not to
modify a
communication
buffer until the communication completes.
Otherwise, the outcome of the computation is undefined.


\begin{implementors}
Since a synchronous send cannot complete before a matching receive is \mpitermni{started},
one will not normally buffer messages sent by such an operation.

It is recommended to choose buffering over blocking the sender, whenever
possible, for standard sends.  The programmer can signal a preference
for blocking the sender until a matching receive occurs by using the
synchronous send mode.

A possible communication protocol for the various communication modes\mpitermindex{semantics!point-to-point communication!send modes}\mpitermindex{point-to-point communication!send modes}
is outlined below.

\begin{description}
\item[ready send:]\mpitermindex{ready mode send}\mpitermindex{modes!ready}The message is sent as soon as possible.

\item[synchronous send:]\mpitermindex{synchronous mode send}\mpitermindex{modes!synchronous}The sender sends a request-to-send message.
The receiver stores this request.
When a matching receive is \mpitermni{started}, the receiver sends back a permission-to-send
message, and the sender now sends the message.

\item[standard send:]\mpitermindex{standard mode send}\mpitermindex{modes!standard}First protocol may be used for short messages, and second protocol for
long messages.

\item[buffered send:]\mpitermindex{buffered mode send}\mpitermindex{modes!buffered}The sender copies the message into a buffer and then sends it with a
nonblocking send (using the same protocol as for standard send).
\end{description}

Additional control messages might be needed for flow control and error
recovery.  Of course, there are many other possible protocols.

Ready send can be implemented as\mpitermdefindex{ready mode send!as standard mode send} a standard send. In this case
there will be no performance advantage (or disadvantage) for the use of
ready send.

A standard send can be implemented as\mpitermdefindex{standard mode send!as synchronous mode send} a synchronous send.
In such a case, no data buffering is needed.  However,
users may
expect some buffering.

In a multithreaded environment, the execution of a blocking communication
should block only the executing thread, allowing the thread scheduler to
de-schedule this thread and schedule another thread for execution.
\end{implementors}


\section{Semantics of Point-to-Point Communication}
\mpitermtitleindex{semantics!point-to-point communication}
\label{sec:pt2pt-semantics}

A valid \MPI/ implementation guarantees certain general properties of
point-to-point communication, which are described in this section.

\paraheading{Order.}
Messages are \mpitermdefindex{semantics!point-to-point communication!nonovertaking}\mpitermdef{nonovertaking}:
If a sender sends two messages in succession to the same destination, and
both match the same receive, then this operation cannot receive the
second message if the first one is still pending.
If a receiver posts two receives in succession, and both match the same
message,
then the second receive operation cannot be satisfied by this message, if the
first one is still pending.
This requirement facilitates matching\mpitermindex{matching rules!ordering} of sends to receives.
It guarantees that message-passing code is \mpitermdefindex{semantics!point-to-point communication!deterministic}deterministic, if \MPI/ processes are
single-threaded and the wildcard \mpiconst{MPI\_ANY\_SOURCE} is not used in receives.
(Some of the calls described later, such as \mpifunc{MPI\_CANCEL} or
\mpifunc{MPI\_WAITANY}, are additional sources of nondeterminism.)

If an \MPI/ process has a single thread of execution, then any two communication operations
executed by this \MPI/ process are \mpitermdefindex{semantics!point-to-point communication!ordered}\mpitermdef{ordered}.

\begin{users}
The \MPI/ Forum believes the following paragraph is ambiguous and may clarify the meaning in a future version of the \MPI/ Standard.
\end{users}

On the other hand, if the \MPI/ process is
multithreaded, then the semantics of thread execution may not define
a relative order between two send operations executed by two
distinct threads.  The operations are \mpitermdefindex{semantics!point-to-point communication!logically concurrent}\mpitermdef{logically concurrent}, even if one
physically precedes the other. In such a case, the two messages sent can be
received in any order.  Similarly, if two receive operations that are \mpitermdef{logically concurrent}
receive two successively sent messages, then the two messages can
match the two receives in either order.

\begin{implementors}
The \MPI/ Forum believes the previous paragraph is ambiguous and may clarify the meaning in a future version of the \MPI/ Standard.
\end{implementors}

\begin{example}
\label{pt2pt-exE}
\Exindex{Fortran}{Nonovertaking messages}{MPI\_Bsend,MPI\_Recv}%
\exindex{Nonovertaking messages}%
An example of nonovertaking messages.

%%HEADER
%%LANG: FORTRAN
%%FRAGMENT
%%DECL: integer comm, count, rank, ierr, tag, buf1(1),buf2(1)
%%DECL: integer status(MPI_STATUS_SIZE)
%%ENDHEADER
\begin{lstlisting}[language={[MPI]Fortran}]
CALL MPI_COMM_RANK(comm, rank, ierr)
IF (rank .EQ. 0) THEN
   CALL MPI_BSEND(buf1, count, MPI_REAL, 1, tag, comm, ierr)
   CALL MPI_BSEND(buf2, count, MPI_REAL, 1, tag, comm, ierr)
ELSE IF (rank .EQ. 1) THEN
   CALL MPI_RECV(buf1, count, MPI_REAL, 0, MPI_ANY_TAG, comm, status, &
                 ierr)
   CALL MPI_RECV(buf2, count, MPI_REAL, 0, tag, comm, status, ierr)
END IF
\end{lstlisting}
The message sent by the first send must be received by the first receive, and
the message sent by the second send must be received by the second receive.
\end{example}


\paraheading{Progress.}\mpitermdefindex{semantics!point-to-point communication!progress}\mpitermdefindex{progress!point-to-point communication}
If a pair of
matching send and receive operations have been initiated, then at
least one of these two operations will complete, independently of
other actions in the system: the send operation will
complete, unless the receive is satisfied by another message, and
completes; the
receive operation will complete, unless the message sent is consumed by another
matching receive that was \mpitermni{started} at the same destination \MPI/ process.

\begin{example}
\label{pt2pt-exF}
\Exindex{Fortran}{Message matching}{MPI\_Bsend,MPI\_Ssend,MPI\_Recv}%
\exindex{Intertwined matching pairs}%
\exindex{Progress of matching pairs}%
An example of two, intertwined matching pairs.
%%HEADER
%%LANG: FORTRAN
%%FRAGMENT
%%DECL: integer comm, rank, count, ierr, tag1,tag2,buf1(2),buf2(2)
%%DECL: integer status(MPI_STATUS_SIZE)
%%ENDHEADER
\begin{lstlisting}[language={[MPI]Fortran}]
CALL MPI_COMM_RANK(comm, rank, ierr)
IF (rank .EQ. 0) THEN
   CALL MPI_BSEND(buf1, count, MPI_REAL, 1, tag1, comm, ierr)
   CALL MPI_SSEND(buf2, count, MPI_REAL, 1, tag2, comm, ierr)
ELSE IF (rank .EQ. 1) THEN
   CALL MPI_RECV(buf1, count, MPI_REAL, 0, tag2, comm, status, ierr)
   CALL MPI_RECV(buf2, count, MPI_REAL, 0, tag1, comm, status, ierr)
END IF
\end{lstlisting}
Both \MPI/ processes invoke their first communication call.
Since the first send at the \MPI/ process with \code{rank = 0} uses the buffered mode, it must complete,
irrespective of the state of the other \MPI/ process(es).  Since no matching receive is
\mpitermni{started}, the message will be copied into buffer space.  (If insufficient
buffer space is available, then the program will fail.)
The second send is then invoked.
At that point, a matching pair of send and receive
operation is enabled, and both operations must complete.
Next, the second receive call is invoked, which will be satisfied by the buffered
message. Note that the \MPI/ process with \code{rank = 1} received the messages in the reverse order they
were sent.
\end{example}

\paraheading{Fairness.}
\MPI/ makes no guarantee of \mpitermdefni{fairness}\mpitermdefindex{fairness!not guaranteed}\mpitermdefindex{semantics!point-to-point communication!fairness not guaranteed} in the handling of
communication.  Suppose that a send is \mpitermni{started}.  Then it is possible
that the destination \MPI/ process repeatedly posts a receive that matches this
send, yet the message is never received, because it is overtaken each time by
another message, sent from another source.  Similarly, suppose that a
receive was \mpitermni{started} by a multithreaded \MPI/ process.  Then it is possible that
messages that
match this receive are repeatedly received, yet the receive is never satisfied,
because it is overtaken by other receives \mpitermni{started} at this \MPI/ process (by
other executing threads).  It is the programmer's responsibility to prevent
starvation in such situations.


\paraheading{Resource limitations.}\mpitermdefindex{semantics!point-to-point communication!resource limitations}
Any \mpitermni{pending}\mpitermindex{pending (communication) operation}\mpitermindex{communication!pending}\mpitermindex{operation!pending} communication operation
and \mpitermni{decoupled \MPI/ activity}\mpitermindex{decoupled MPI activities@decoupled \MPI/ activities}
consumes system resources that are limited.
Errors may occur when lack of resources prevent the execution of an \MPI/ call.
High-quality implementations will use a (small) fixed amount of resources for
each \mpitermni{pending} send in the ready or synchronous mode and for each
\mpitermni{pending} receive. However, buffer
space may be consumed to store messages sent in standard mode, and must
be consumed to store
messages sent in buffered mode, when no matching receive is available.
The amount of space available for buffering will be much smaller than program
data memory on many systems.  Then, it will be easy to write programs that
overrun available buffer space.

\MPI/ allows the user to provide buffer memory for messages sent in the buffered
mode.  Furthermore, \MPI/ specifies a detailed operational
model for the use of this
buffer.  An \MPI/ implementation is required
to do no worse than implied by this model.  This allows users to avoid buffer
overflows when they use buffered sends.  Buffer allocation and use is
described
in Section~\ref{sec:pt2pt-buffer}.

A buffered send operation that cannot complete because of a lack of buffer space\mpitermindex{erroneous program!lack of buffer space}\mpitermindex{semantics!erroneous program!lack of buffer space}
is \mpitermni{erroneous}.  When such a situation is detected, an error is signaled that may
cause the program to terminate abnormally.
On the other hand, a standard send
operation that cannot complete because
of lack of buffer space will merely block, waiting for buffer space to become
available or for a matching receive to be \mpitermni{started}.  This behavior is preferable
in many situations.  Consider a situation where a producer repeatedly produces
new values and sends them to a consumer.  Assume that the producer produces
new values faster than the consumer can consume them.  If buffered sends are
used, then a buffer overflow will result.  Additional synchronization has to
be added to the program so as to prevent this from occurring.
If standard sends are used, then the producer will be automatically throttled,
as its send operations will block when buffer space is unavailable.

In some situations, a lack of buffer space leads to deadlock\mpitermindex{semantics!deadlock!buffer space} situations.
This is illustrated by the examples below.

\begin{example}
\label{pt2pt-exG}
\Exindex{Fortran}{Message exchange}{MPI\_Send,MPI\_Recv}%
\exindex{Message exchange (ping-pong)}%
An exchange of messages.

%%HEADER
%%LANG: FORTRAN
%%FRAGMENT
%%DECL: integer comm, count, rank, ierr, tag
%%DECL: integer sendbuf(2), recvbuf(2)
%%DECL: integer status(MPI_STATUS_SIZE)
%%ENDHEADER
\begin{lstlisting}[language={[MPI]Fortran}]
CALL MPI_COMM_RANK(comm, rank, ierr)
IF (rank .EQ. 0) THEN
   CALL MPI_SEND(sendbuf, count, MPI_REAL, 1, tag, comm, ierr)
   CALL MPI_RECV(recvbuf, count, MPI_REAL, 1, tag, comm, status, ierr)
ELSE IF (rank .EQ. 1) THEN
   CALL MPI_RECV(recvbuf, count, MPI_REAL, 0, tag, comm, status, ierr)
   CALL MPI_SEND(sendbuf, count, MPI_REAL, 0, tag, comm, ierr)
END IF
\end{lstlisting}
This program will succeed even if no buffer space for data
is available.  The standard send operation can be replaced, in this example,
with a synchronous send.
\end{example}

\begin{example}
\label{pt2pt-exH}
\Exindex{Fortran}{Errant message exchange}{MPI\_Send,MPI\_Recv}%
\exindex{Deadlock!wrong message exchange}%
An errant attempt to exchange messages.
%%HEADER
%%LANG: FORTRAN
%%FRAGMENT
%%DECL: integer comm, rank, count, ierr, tag
%%DECL: integer sendbuf(10), recvbuf(10)
%%DECL: integer status(MPI_STATUS_SIZE)
%%ENDHEADER
\begin{lstlisting}[language={[MPI]Fortran}]
! ---------------- THIS EXAMPLE IS ERRONEOUS ---------------
CALL MPI_COMM_RANK(comm, rank, ierr)
IF (rank .EQ. 0) THEN
   CALL MPI_RECV(recvbuf, count, MPI_REAL, 1, tag, comm, status, ierr)
   CALL MPI_SEND(sendbuf, count, MPI_REAL, 1, tag, comm, ierr)
ELSE IF (rank .EQ. 1) THEN
   CALL MPI_RECV(recvbuf, count, MPI_REAL, 0, tag, comm, status, ierr)
   CALL MPI_SEND(sendbuf, count, MPI_REAL, 0, tag, comm, ierr)
END IF
\end{lstlisting}
The receive operation of the \MPI/ process with \code{rank = 0} must complete before its send, and
can complete only if the matching send
of the \MPI/ process with \code{rank = 1} is executed. The receive operation of the
\MPI/ process with \code{rank = 1} must complete before its send and
can complete only if the matching send of the \MPI/ process with \code{rank = 0} is executed.
This program will always deadlock.  The same holds for any other send mode.
\end{example}

\begin{example}
\label{pt2pt-exI}
\Exindex{Fortran}{Exchange relies on buffering}{MPI\_Send,MPI\_Recv}%
\exindex{Deadlock!if not buffered}%
An unsafe exchange that relies on \MPI/ to provide sufficient buffering.
%%HEADER
%%LANG: FORTRAN
%%FRAGMENT
%%DECL: integer comm, rank, ierr, tag,count
%%DECL: integer sendbuf(10), recvbuf(10)
%%DECL: integer status(MPI_STATUS_SIZE)
%%ENDHEADER
\begin{lstlisting}[language={[MPI]Fortran}]
! ---------------- THIS EXAMPLE IS ERRONEOUS ---------------
CALL MPI_COMM_RANK(comm, rank, ierr)
IF (rank .EQ. 0) THEN
   CALL MPI_SEND(sendbuf, count, MPI_REAL, 1, tag, comm, ierr)
   CALL MPI_RECV(recvbuf, count, MPI_REAL, 1, tag, comm, status, ierr)
ELSE IF (rank .EQ. 1) THEN
   CALL MPI_SEND(sendbuf, count, MPI_REAL, 0, tag, comm, ierr)
   CALL MPI_RECV(recvbuf, count, MPI_REAL, 0, tag, comm, status, ierr)
END IF
\end{lstlisting}
The message sent by each \MPI/ process has to be copied out before the send operation
completes and the receive operation starts.  For the program to complete, it is
necessary that at least one of the two messages sent be buffered.
Thus, this program can
succeed only if
the communication system can buffer at least \mpiarg{count} words of data.
\end{example}

\begin{users}
If standard mode send operations are used as in Example~\ref{pt2pt-exI}, then a deadlock situation may occur
where both \MPI/ processes are blocked because sufficient buffer space is not
available.
The same will certainly happen, if the synchronous mode is used.
If the buffered mode is used, and not enough buffer space is available,
then the
program will not complete either.  However, rather than a deadlock situation,
we shall have a buffer overflow error.

A portable program using standard mode send operations should not rely on message
buffering for the program to complete without \mpitermindex{deadlock avoidance!send modes}\mpitermni{deadlock}.
All sends in such a portable program can be replaced
with synchronous mode sends and the program will still run correctly.
The buffered send mode can be used for programs that require buffering.

Nonblocking message-passing operations, as described in
Section~\ref{sec:pt2pt-nonblock}, can be used to avoid the need for buffering
outgoing messages.  This can prevent unintentional \mpitermni{serialization}\mpitermindex{serialization!lack of buffer space} or \mpitermni{deadlock}\mpitermindex{deadlock avoidance!lack of buffer space} due to lack of buffer space, and
improves performance, by allowing \mpitermni{overlap}\mpitermindex{communication!overlap with computation}\mpitermindex{communication!overlap with communication} of communication with other communication or with computation,
and avoiding the overheads of allocating buffers and copying messages into
buffers.
\end{users}

\section{Buffer Allocation and Usage}
\mpitermtitleindex{point-to-point communication!buffer allocation}
\mpitermtitleindex{buffered mode send!buffer allocation}
\mpitermtitleindex{automatic buffering!buffer allocation}
\label{sec:pt2pt-buffer}

A user may specify up to one buffer per communicator, up to one buffer per session, and up to one buffer per \MPI/ process to be used for buffering messages sent in buffered mode.  Buffering is done by the sender.

\begin{mpi-binding}
    function_name("MPI_Comm_attach_buffer")
    parameter("comm", "COMMUNICATOR")
    parameter(
        "buffer",
        "BUFFER",
        desc=r"initial buffer address",
        asynchronous=True,
        suppress="f08_intent",
    )
    parameter("size", "POLYNUM_BYTES_NNI", desc=r"buffer size, in bytes")
\end{mpi-binding}

Provides to \MPI/ a communicator-specific buffer in memory. This is to be used for buffering outgoing
messages sent when a buffered mode send is started that uses the communicator \mpiarg{comm}.

If \mpiconstmain{MPI\_BUFFER\_AUTOMATIC} is passed as the argument \mpiarg{buffer}, no explicit buffer is attached; rather, automatic buffering is enabled for all buffered mode communication associated with the communicator \mpiarg{comm} (see Section~\ref{subsec:pt2pt:auto}).
Further, if \mpiconst{MPI\_BUFFER\_AUTOMATIC} is passed as the argument \mpiarg{buffer}, the value of \mpiarg{size} is irrelevant.
Note that in Fortran \mpiconst{MPI\_BUFFER\_AUTOMATIC} is an
object like \mpiconst{MPI\_BOTTOM} (not usable for initialization or
assignment), see Section~\ref{subsec:named-constants}.



\begin{mpi-binding}
    function_name("MPI_Session_attach_buffer")
    parameter("session", "SESSION")
    parameter(
        "buffer",
        "BUFFER",
        desc=r"initial buffer address",
        asynchronous=True,
        suppress="f08_intent",
    )
    parameter("size", "POLYNUM_BYTES_NNI", desc=r"buffer size, in bytes")
\end{mpi-binding}

Provides to \MPI/ a session-specific buffer in memory. This buffer is to be used for buffering outgoing
messages sent when using a communicator that is created from a group that is derived from the session \mpiarg{session}. However, if there is a communicator-specific buffer attached to the particular communicator at the time of the buffered mode send is started, that buffer is used.

If \mpiconst{MPI\_BUFFER\_AUTOMATIC} is passed as the argument \mpiarg{buffer}, no explicit buffer is attached; rather, automatic buffering is enabled for all buffered mode communication associated with the session \mpiarg{session} that is not explicitly covered by a buffer provided at communicator level (see Section~\ref{subsec:pt2pt:auto}).
Further, if \mpiconst{MPI\_BUFFER\_AUTOMATIC} is passed as the argument \mpiarg{buffer}, the value of \mpiarg{size} is irrelevant.
Note that in Fortran \mpiconst{MPI\_BUFFER\_AUTOMATIC} is an
object like \mpiconst{MPI\_BOTTOM} (not usable for initialization or
assignment), see Section~\ref{subsec:named-constants}.

\begin{mpi-binding}
    function_name("MPI_Buffer_attach")
    parameter(
        "buffer",
        "BUFFER",
        desc=r"initial buffer address",
        asynchronous=True,
        suppress="f08_intent",
    )
    parameter("size", "POLYNUM_BYTES_NNI", desc=r"buffer size, in bytes")
\end{mpi-binding}

Provides to \MPI/ an \MPI/ process-specific buffer in memory. This buffer is to be used for buffering outgoing
messages sent when using a communicator to which no communicator-specific buffer is attached or which is derived from a session to which no session-specific buffer is attached at the time the buffered mode send is started.

If \mpiconst{MPI\_BUFFER\_AUTOMATIC} is passed as the argument \mpiarg{buffer}, no explicit buffer is attached; rather, automatic buffering is enabled for all buffered mode communication not explicitly covered by a buffer provided at session or communicator level (see Section~\ref{subsec:pt2pt:auto}).
Further, if \mpiconst{MPI\_BUFFER\_AUTOMATIC} is passed as the argument \mpiarg{buffer}, the value of \mpiarg{size} is irrelevant.
Note that in Fortran \mpiconst{MPI\_BUFFER\_AUTOMATIC} is an
object like \mpiconst{MPI\_BOTTOM} (not usable for initialization or
assignment), see Section~\ref{subsec:named-constants}.


\begin{users}
The use of a global shared buffer can be problematic when used for communication in different libraries, as the buffer represents a shared resource used for all buffered mode communication. Further, with the introduction of the \spm, the use of a single shared buffer violates the concept of resource isolation that is intended with \MPI/ Sessions. It is therefore recommended, especially for libraries and programs using the \spm, to use only communicator-specific or session-specific buffers.
\end{users}

Any of these buffers are used only for messages sent in buffered mode.
Only one \MPI/ process-specific buffer can be attached to an \MPI/ process at a time,
only one session-specific buffer can be attached to a session at a time and
only one communicator-specific buffer can be attached to a communicator at a time.

If automatic buffering is enabled at any level, no other buffer can be attached at that level.

A particular memory region can only be used in one buffer; reusing buffer space
for multiple sessions, communicators and/or the global buffer is erroneous.
Further, only one buffer is used for any one communication
following the rules above; buffer space is not combined, even if two buffers are directly
or indirectly provided  to a communicator to be used for buffered sends.

In C, \mpiarg{buffer} is the starting address of a memory region.
In Fortran, one can pass the first element of a memory region
or a whole array, which must be `simply contiguous'
(for `simply contiguous,' see also
\sectionref{sec:misc-sequence}).

% MPI-4.0: changed the LIS for the "size" parameter from "nonnegative
% integer" to "integer" (to allow it to be MPI_UNDEFINED).
\begin{mpi-binding}
    function_name("MPI_Comm_detach_buffer")
    parameter("comm", "COMMUNICATOR")
    parameter(
        "buffer_addr",
        "C_BUFFER2",
        desc=r"initial buffer address",
        direction="out",
    )
    parameter(
        "size",
        "POLYNUM_BYTES",
        desc=r"buffer size, in bytes",
        direction="out",
    )
\end{mpi-binding}

Detach the communicator-specific buffer currently attached to the communicator.


\begin{mpi-binding}
    function_name("MPI_Session_detach_buffer")
    parameter("session", "SESSION")
    parameter(
        "buffer_addr",
        "C_BUFFER2",
        desc=r"initial buffer address",
        direction="out",
    )
    parameter(
        "size",
        "POLYNUM_BYTES",
        desc=r"buffer size, in bytes",
        direction="out",
    )
\end{mpi-binding}

Detach the session-specific buffer currently attached to the session.

\begin{mpi-binding}
    function_name("MPI_Buffer_detach")
    parameter(
        "buffer_addr",
        "C_BUFFER2",
        desc=r"initial buffer address",
        direction="out",
    )
    parameter(
        "size",
        "POLYNUM_BYTES",
        desc=r"buffer size, in bytes",
        direction="out",
    )
\end{mpi-binding}

Detach the \MPI/ process-specific buffer buffer currently attached to \MPI/.

The procedure calls return the
address and the size of the detached buffer.
If \mpiconst{MPI\_BUFFER\_AUTOMATIC} was used in the corresponding attach procedure, then \mpiconst{MPI\_BUFFER\_AUTOMATIC} is also returned in the detach procedure and the value returned in argument \mpiarg{size} is undefined.
In this case, automatic buffering is disabled upon return from the detach procedure.
When using Fortran \code{mpi\_f08}, the returned value is identical to \code{c\_loc(MPI\_BUFFER\_AUTOMATIC)}. Note that \code{c\_loc()} is an intrinsic in the Fortran \ftype{ISO\_C\_BINDING} module.
\begin{implementors}
In Fortran, the implementation of \mpiconst{MPI\_BUFFER\_AUTOMATIC} must allow the intrinsic \code{c\_loc} to be applied to it.
\end{implementors}

These procedures
will delay their return until all messages currently in the (explicit or automatic) buffer have been transmitted.
Upon return of
these procedures, the user may reuse or deallocate the space taken by the buffer.

If the size of the detached buffer cannot be represented in
\mpiarg{size}, it is set to \mpiconst{MPI\_UNDEFINED}.

% The following text is introducing a new set of routines - flush_buffer
\medskip %%ALLOWLATEX%%

The following \mpiskipfunc{MPI\_XXX\_FLUSH\_BUFFER} procedures, as well as \mpiskipfunc{MPI\_BUFFER\_FLUSH} itself, will not return until all messages currently in the buffer have been transmitted without detaching the buffer.

\begin{rationale}
These flush procedures provide the same functionality as an atomic combination of first detaching the buffer and then attaching it again (but without having to actually execute the detaching and the re-attaching of the buffer), but they may be implemented with less internal overhead.
\end{rationale}

\begin{mpi-binding}
    function_name("MPI_Comm_flush_buffer")
    parameter("comm", "COMMUNICATOR")
\end{mpi-binding}

\mpiskipfunc{MPI\_COMM\_FLUSH\_BUFFER} will not return until all messages currently in the communicator\hyphen/specific buffer of the calling \MPI/ process have been transmitted.

\begin{mpi-binding}
    function_name("MPI_Session_flush_buffer")
    parameter("session", "SESSION")
\end{mpi-binding}

\mpiskipfunc{MPI\_SESSION\_FLUSH\_BUFFER} will not return until all messages currently in the session-specific buffer of the calling \MPI/ process have been transmitted.

\begin{mpi-binding}
    function_name("MPI_Buffer_flush")
\end{mpi-binding}

\mpiskipfunc{MPI\_BUFFER\_FLUSH} will not return until all messages currently in the \MPI/ process-specific buffer of the calling \MPI/ process have been transmitted.


For all \mpiskipfunc{MPI\_XXX\_FLUSH\_BUFFER} procedures, there also exist the following nonblocking variants, which start the respective flush operation. These operations will not complete until all messages currently in the respective  buffer of the calling \MPI/ process have been transmitted.

\begin{mpi-binding}
    function_name("MPI_Comm_iflush_buffer")
    parameter("comm", "COMMUNICATOR")
    parameter("request", "REQUEST", direction="out")
\end{mpi-binding}

\begin{mpi-binding}
    function_name("MPI_Session_iflush_buffer")
    parameter("session", "SESSION")
    parameter("request", "REQUEST", direction="out")
\end{mpi-binding}

\begin{mpi-binding}
    function_name("MPI_Buffer_iflush")
    parameter("request", "REQUEST", direction="out")
\end{mpi-binding}


\begin{example}
\label{pt2pt-exZ}
\Exindex{C}{Attach and detach buffer}{MPI\_Buffer\_attach,MPI\_Buffer\_detach}%
Calls to attach and detach buffers.
%%HEADER
%%LANG: C
%%FRAGMENT
%%ENDHEADER
\begin{lstlisting}[language={[MPI]C}]
#define BUFFSIZE 10000+MPI_BSEND_OVERHEAD
int size;
char *buff;
MPI_Buffer_attach(malloc(BUFFSIZE), BUFFSIZE);
/* a buffer of 10000 bytes can now be used by MPI_Bsend */
/* on all communicators, assuming only one message at a time is sent */
MPI_Buffer_detach(&buff, &size);
/* Buffer size reduced to zero */
MPI_Buffer_attach(buff, size);
/* Buffer of 10000 bytes available again */
\end{lstlisting}
\end{example}

\begin{example}
\label{pt2pt-exY}
\Exindex{C}{Attach and detach communicator-specific buffer}{MPI\_Comm\_attach\_buffer,MPI\_Comm\_detach\_buffer}%
Calls to attach and detach communicator-specific buffers.
%%HEADER
%%LANG: C
%%FRAGMENT
%%ENDHEADER
\begin{lstlisting}[language={[MPI]C}]
#define BUFFSIZE1 10000+MPI_BSEND_OVERHEAD
#define BUFFSIZE2 20000+MPI_BSEND_OVERHEAD
int size;
char *buff1, *buff2;
MPI_Comm world_dup;
MPI_Comm_dup(MPI_COMM_WORLD, &world_dup);
MPI_Comm_attach_buffer(MPI_COMM_WORLD, malloc(BUFFSIZE2), BUFFSIZE2);
MPI_Buffer_attach(malloc(BUFFSIZE1), BUFFSIZE1);
/* a buffer of 20000 bytes can now be used by MPI_Bsend for */
/* communication using MPI_COMM_WORLD,  assuming only one message */
/* at a time is sent a buffer of 10000 bytes can now be used by */
/* MPI_Bsend for communication using any other communicator, */
/* including world_dup assuming only one message at a time is sent */
MPI_Comm_detach_buffer(MPI_COMM_WORLD, &buff1, &size);
MPI_Buffer_detach(&buff2, &size);
/* Both buffers are detached and no specific or MPI process-specific */
/* buffer can be used for further MPI_Bsend */
\end{lstlisting}
\end{example}

\begin{users}
Even though the C procedures
\cfunc{MPI\_Buffer\_attach},
\cfunc{MPI\_Session\_attach\_buffer},
\cfunc{MPI\_Comm\_attach\_buffer},
\cfunc{MPI\_Buffer\_detach},
\cfunc{MPI\_Session\_detach\_buffer} and
\cfunc{MPI\_Comm\_detach\_buffer}
have
an argument of type \code{void*}, these arguments are used
differently: a pointer to the buffer is passed to
\cfunc{MPI\_Buffer\_attach},
\cfunc{MPI\_Session\_attach\_buffer} and
\cfunc{MPI\_Comm\_attach\_buffer};
the address of the pointer is passed to
\cfunc{MPI\_Buffer\_detach},
\cfunc{MPI\_Session\_detach\_buffer} and
\cfunc{MPI\_Comm\_detach\_buffer}, so that
this call can return the pointer value.
In Fortran with the \code{mpi} module or (deprecated) \code{mpif.h}, the type of the \mpiarg{buffer\_addr} argument is
wrongly defined and the argument is therefore unused.
In Fortran with the \code{mpi\_f08} module, the address of the buffer is returned
as \ftype{TYPE(C\_PTR)}, see also
\namedref{Example}{ex:1side-fmalloc-ptr}
about the use of \ftype{C\_PTR} pointers.
\end{users}

\begin{rationale}
In all cases, arguments are defined to be of type \code{void*} (rather than
\code{void*} and \code{void**}, respectively), so as to avoid complex type
casts. E.g., in the last two examples, \code{\&buff}, which is of type
\code{char**}, can be passed as argument to \cfunc{MPI\_Buffer\_detach},
\cfunc{MPI\_Session\_detach\_buffer} and
\cfunc{MPI\_Comm\_detach\_buffer}
without type casting.  If the formal parameter had type \code{void**} then we
would need a type cast before and after each call.
\end{rationale}


\paraheading{General semantics of buffered mode sends.}\mpitermdefindex{semantics!point-to-point communication!buffered mode send}
The statements made in this section describe the behavior of \MPI/
for buffered-mode sends.

When no \MPI/ process-specific buffer is currently (explicitly) attached and if no automatic buffering is enabled, \MPI/ behaves as if a
zero-sized \MPI/ process-specific buffer is (implicitly) attached.

It is erroneous to detach a communicator-specific, session-specific, or \MPI/ process-specific buffer, if no such buffer had been attached using a corresponding attach procedure. This includes attach procedure calls using \mpiconst{MPI\_BUFFER\_AUTOMATIC} as the buffer argument.
It is erroneous to attach a communicator-specific, session-specific, or \MPI/ process-specific buffer, if such buffer had already been attached using a corresponding attach procedure and not yet been detached again. This includes attach procedure calls using \mpiconst{MPI\_BUFFER\_AUTOMATIC} as the buffer argument.
It is erroneous to flush a communicator-specific, session-specific or \MPI/ process-specific buffer, if there is no buffer attached (including automatic buffering).

\mpifunc{MPI\_COMM\_ATTACH\_BUFFER}, \mpifunc{MPI\_SESSION\_ATTACH\_BUFFER}, and \mpifunc{MPI\_BUFFER\_ATTACH} are local.
\mpifunc{MPI\_COMM\_DETACH\_BUFFER}, \mpifunc{MPI\_SESSION\_DETACH\_BUFFER},  \mpifunc{MPI\_BUFFER\_DETACH},\flushline
 \mpifunc{MPI\_COMM\_FLUSH\_BUFFER},
\mpifunc{MPI\_SESSION\_FLUSH\_BUFFER}, and \mpifunc{MPI\_BUFFER\_FLUSH}
are nonlocal;
they must not return before all buffered messages in their related buffers are transmitted,
and they must eventually return when all corresponding receive operations are started
(provided that none are cancelled).

\paraheading{Automatic buffering with buffered mode sends.}\mpitermdefindex{semantics!point-to-point communication!buffered mode send automatic buffering}
\label{subsec:pt2pt:auto}
If the buffer used at the time of buffered mode send is set to the buffer address \mpiconst{MPI\_BUFFER\_AUTOMATIC}, then a buffer of sufficient size is
automatically used by the \MPI/ library.

\begin{users}
When using automatic buffering, the user relinquishes control over buffer management, including allocation and deallocation decisions and timing, to the \MPI/ library. If explicit control is needed over when and how much buffer space is allocated, automatic buffering must not be used.
\end{users}

\begin{implementors}
High-quality implementations of an \MPI/ library should strive to support automatic buffering in a balanced fashion, i.e., providing the right balance between memory allocated for send operations and memory available for the end user.
\end{implementors}

The flush operations for the communicator-specific, session-specific, and \MPI/ process-specific buffers
can also be used for automatic buffering. The flush procedure will not return until all automatically allocated buffers
for the communicator-specific, session-specific, or \MPI/ process-specific buffers, respectively, no longer hold message data and could be deallocated by the \MPI/ library, if it chooses to do so.

\begin{users}
With standard mode send, the limitation of needed buffer space is implemented within
the \MPI/ library through switching from internal buffering to internal synchronous mode.
If the user wants to limit the automatically allocated buffer space for buffered mode send
using automatic buffering, the user may call explicitly the appropriate flush procedure to wait
until automatically allocated buffers are deallocated.
\end{users}

\paraheading{Further rules.}
In the case of an attached buffer (i.e., not using automatic buffering),
the user must provide as much buffering for outgoing messages as would be required if
outgoing message
data were buffered by the sending \MPI/ process, in the specified buffer space,
using a circular, contiguous-space allocation policy.
We outline below a model implementation that defines this policy.
\MPI/ may provide more buffering, and may use a better buffer allocation
algorithm than described below.
On the other hand, \MPI/ may signal an error whenever the
simple buffering allocator described below would run out of space.
\MPI/ must not require more buffer space as described in the model implementation below.



\MPI/ does not provide mechanisms for querying or controlling buffering done by
standard mode sends.  It is expected that vendors will provide such
information
for their implementations.

\begin{rationale}
There is a wide spectrum of possible implementations of buffered communication operations:
buffering can be done at sender, at receiver, or both; buffers can be dedicated
to one sender-receiver pair, or be shared by
all communication operations; buffering can be done in real or
in virtual memory; it can use dedicated memory, or memory shared by other
\MPI/ processes; buffer space may be allocated statically or be changed dynamically;
etc.  It does not seem feasible to provide a portable mechanism for querying
or controlling buffering that would be compatible with all these choices, yet
provide meaningful information.
\end{rationale}


\subsection{Model Implementation of Buffered Mode}
\label{subsec:pt2pt-modelbuffered}

The model implementation uses the packing and unpacking procedures described in
Section~\ref{sec:pt2pt-packing} and the nonblocking communication procedures
described in Section~\ref{sec:pt2pt-nonblock}.

We assume that a circular queue
of pending message entries (PME) is maintained.  Each
entry contains a communication request handle
that identifies a pending nonblocking
send, a pointer to the next entry and the packed message data.  The
entries are stored in successive locations in the buffer.  Free space is
available between the queue tail and the queue head.

A buffered send call results in the execution of the following algorithm:
\begin{itemize}
\item
Traverse sequentially the PME queue from head towards the tail,
deleting all entries for
communication operations that have completed, up to the first entry with an uncompleted
request; update queue head to point to that entry.
\item
Compute the number of bytes, $n$, needed to store an entry for the new message.
An upper bound on $n$ can be computed
as follows: A call to the function\flushline
\mpifunc{MPI\_PACK\_SIZE}\mpicode{(count, datatype,
comm, size)}, with the \mpicode{count}, \mpicode{datatype} and \mpicode{comm} arguments
used in the
\mpifunc{MPI\_BSEND} call, returns an upper bound on the amount of
space needed to buffer the message data (see Section~\ref{sec:pt2pt-packing}).
The \MPI/ constant
\mpiconstmain{MPI\_BSEND\_OVERHEAD} provides an upper bound on the
additional space consumed by the entry (e.g., for pointers or \mpiterm{envelope}\mpitermindex{message!envelope}
information).
\item
Find the next contiguous empty space of $n$ bytes in
buffer (space following queue tail, or space at start of buffer if
queue tail is
too close to end of buffer).  If space is not found then raise buffer
overflow error.
\item
Append to end of PME queue in contiguous space the new entry
that contains request
handle, next pointer and packed message data; \mpifunc{MPI\_PACK} is used to
pack data.
\item
Post nonblocking send (standard mode) for packed data.
\item
Return
\end{itemize}

\section{Nonblocking Communication}
\mpitermtitleindex{point-to-point communication!nonblocking}
\mpitermtitleindex{nonblocking}
\label{sec:pt2pt-nonblock}

\mpitermdefni{Nonblocking communication}\mpitermdefindex{nonblocking!communication}
is important both for reasons of correctness and performance.
For complex communication patterns, the use of only blocking communication
(without buffering) is difficult because the programmer must ensure that
each send is matched with a receive in an order that
avoids \mpitermni{deadlock}\mpitermindex{deadlock avoidance!nonblocking communication}.
For communication patterns that are determined only at run time, this is even
more difficult.
Nonblocking communication can be used to avoid this problem, allowing programmers to
express complex and possibly dynamic communication patterns without needing to
ensure that all sends and receives are issued in an order that prevents
deadlock (see Section~\ref{sec:pt2pt-semantics} and the discussion of ``safe'' programs).
Nonblocking communication also allows for the \mpitermni{overlap}\mpitermindex{communication!overlap with communication}
of communication with different communication operations,
e.g., to prevent the unintentional \mpiterm{serialization} of such operations,
and for the \mpitermni{overlap}\mpitermindex{communication!overlap with computation}
of communication with computation.
Whether an implementation is able to accomplish an effective (from a
performance standpoint) overlap of operations depends on the implementation
itself and the system on which the implementation is running.
Using nonblocking operations \emph{permits} an implementation to overlap\mpitermdefindex{semantics!overlap}
communication with computation, but does not require it to do so.

A nonblocking \mpitermdefni{send start}\mpitermdefindex{send!start}\mpitermindex{send!nonblocking} call \mpitermni{initiates}\mpitermindex{initiation} the send operation, but does not
complete it.  The send start call
can
return before the message was copied out of the send buffer.
A separate \mpitermdefni{send complete}\mpitermdefindex{send!complete}\mpitermindex{completion}
call is needed to complete the communication, i.e., to verify that the
data has been copied out of the send buffer.  With
suitable hardware, the transfer of data out of the sender memory
may proceed concurrently with computations done at the sender after
the send was initiated and before it completed.
Similarly, a nonblocking \mpitermdefni{receive start}\mpitermdefindex{receive!start}\mpitermindex{receive!nonblocking} call \mpitermni{initiates}\mpitermindex{initiation} the receive
operation, but does not complete it.  The call
can
return before
a message is stored into the receive buffer.  A separate
\mpitermdefni{receive complete}\mpitermdefindex{receive!complete}\mpitermindex{completion} call
is needed to complete the receive operation and verify that the data has
been received into the receive buffer.
With suitable hardware, the transfer of data into the receiver memory
may proceed concurrently with computations done after the receive was
initiated and before it completed.  The use of nonblocking receives may also
avoid system buffering and memory-to-memory copying, as information is provided
early on the location of the receive buffer.

Nonblocking send start calls can use the same four modes as blocking sends:
\mpitermni{standard}, \mpitermni{buffered}, \mpitermni{synchronous}, and
\mpitermni{ready}.
These carry the same meaning. 
Sends of all modes, \mpitermni{ready} excepted, can be started whether a matching
receive has been started or not; a nonblocking \mpitermdefni{ready}\mpitermdefindex{ready mode send!nonblocking}
send can be started only if
the matching receive is already started.  In all cases, the send start call
is \mpiterm{local}: it returns immediately\mpitermindex{immediate}, irrespective of the
status of other \MPI/ processes.
If the call causes some system resource to be exhausted, then it will
fail and return an error code.  High-quality
implementations of \MPI/ should ensure that this happens only
in ``pathological'' cases.  That is, an \MPI/ implementation
should be able to support a
large number of \mpitermni{pending}\mpitermindex{pending (communication) operation}\mpitermindex{communication!pending}\mpitermindex{operation!pending} nonblocking operations.

The send-complete call returns no earlier than when all message data has been copied out of the send buffer.
It may carry additional meaning, depending on the send mode.

If the send mode is \mpitermdefni{synchronous}\mpitermdefindex{synchronous mode send!nonblocking},
then the send-complete call is \mpiterm{nonlocal};
the send can complete only if a matching receive has been started
and has been matched with the send.
Note that a synchronous mode send may complete,
if matched by a nonblocking receive,
before the receive complete call occurs.
(It can complete as soon as the sender ``knows'' the transfer will complete,
but before the receiver ``knows'' the transfer will complete.)

If the send mode is \mpitermdefni{buffered}\mpitermdefindex{buffered mode send!nonblocking},
then the send-complete call is \mpiterm{local};
the send must complete irrespective of the status of a matching receive.
If there is no \mpitermni{pending}\mpitermindex{pending (communication) operation}\mpitermindex{communication!pending}\mpitermindex{operation!pending}
receive operation, then the message must be buffered.

If the send mode is \mpitermdefni{standard}\mpitermdefindex{standard mode send!nonblocking},
then the send-complete call can be either \mpiterm{local} or \mpiterm{nonlocal}.
If the message is buffered, it is permitted for the send to complete
before a matching receive is started.
On the other hand, it is permitted for the send not to complete
until a matching receive has been started
and the message has been copied into the receive buffer.

Nonblocking sends can be matched\mpitermindex{matching rules!blocking with nonblocking} with blocking receives, and
vice-versa.

\begin{users}
The completion of a send operation may be delayed for standard mode, and must
be delayed for synchronous mode, until a matching receive has been started.
The use of nonblocking sends in these two cases allows the sender to proceed
ahead of the receiver, so that the computation is more tolerant
of fluctuations in the speeds of the two \MPI/ processes.

Nonblocking sends in the buffered and ready modes have a more limited
impact, e.g., the blocking version of buffered send is capable of
completing regardless of when a matching receive call is made. However,
separating the start from the completion of these sends still gives
some opportunity for optimization within the \MPI/ library. For example,
starting a buffered send gives an implementation more flexibility in
determining if and how the message is buffered. There are also advantages
for both nonblocking buffered and ready modes when data copying can
be done concurrently with computation.

The message-passing model implies that communication is initiated by
the sender.
The communication will generally have lower overhead if a receive is
already \mpitermni{started} when the sender initiates the communication (data can be moved
directly to the receive buffer, and there is no need to queue a pending send
request).  However,
a receive operation can complete only after the matching send has \mpitermni{started}.
The use of nonblocking receives allows one to achieve lower communication overheads
without blocking the receiver while it waits for the send.
\end{users}

\subsection{Communication Request Objects}
\mpitermtitleindexmainsub{nonblocking}{request objects}
\label{subsec:pt2pt-commobject}

Nonblocking communication operations use opaque \mpitermdefni{request} objects to
identify communication operations and match the operation that
initiates the communication with the operation that terminates it.
These are system objects that are accessed via a handle.
A request object identifies various properties of a communication
operation, such as the send mode, the communication
buffer that is associated with it, its
context, the tag and destination arguments to be used for a send, or the tag
and source arguments to be used for a receive.  In addition, this object stores
information about the status of the \mpitermni{pending}\mpitermindex{pending (communication) operation}\mpitermindex{communication!pending}\mpitermindex{operation!pending} communication operation.

\subsection{Communication Initiation}
\mpitermtitleindexmainsub{nonblocking}{initiation}
\label{subsec:pt2pt-commstart}

For the functions defined in this section, we use the same naming conventions as for blocking communication: a
prefix of \mpicode{B}, \mpicode{S}, or \mpicode{R} is used for
\mpitermni{buffered}\mpitermindex{buffered mode send!nonblocking},
\mpitermni{synchronous}\mpitermindex{synchronous mode send!nonblocking}, or
\mpitermni{ready}\mpitermindex{ready mode send!nonblocking} mode.
In addition, for these functions a prefix of \mpicode{I} (for \mpiterm{immediate} and \mpiterm{incomplete}) indicates
that the call is nonblocking.

\cdeclindex{MPI\_Request}
\begin{mpi-binding}
    function_name("MPI_Isend")
    parameter(
        "buf",
        "BUFFER",
        asynchronous=True,
        constant=True,
        desc=r"initial address of send buffer",
    )
    parameter(
        "count",
        "POLYXFER_NUM_ELEM_NNI",
        desc=r"number of elements in send buffer",
    )
    parameter(
        "datatype", "DATATYPE", desc=r"datatype of each send buffer element"
    )
    parameter("dest", "RANK", desc=r"rank of destination")
    parameter("tag", "TAG", desc=r"message tag")
    parameter("comm", "COMMUNICATOR")
    parameter("request", "REQUEST", direction="out")
\end{mpi-binding}


Start\mpitermindex{send!start} a standard mode\mpitermindex{standard mode send!nonblocking} nonblocking send.


\cdeclindex{MPI\_Request}
\begin{mpi-binding}
    function_name("MPI_Ibsend")
    parameter(
        "buf",
        "BUFFER",
        constant=True,
        asynchronous=True,
        desc=r"initial address of send buffer",
    )
    parameter(
        "count",
        "POLYXFER_NUM_ELEM_NNI",
        desc=r"number of elements in send buffer",
    )
    parameter(
        "datatype", "DATATYPE", desc=r"datatype of each send buffer element"
    )
    parameter("dest", "RANK", desc=r"rank of destination")
    parameter("tag", "TAG", desc=r"message tag")
    parameter("comm", "COMMUNICATOR")
    parameter("request", "REQUEST", direction="out")
\end{mpi-binding}

Start\mpitermindex{send!start} a buffered mode\mpitermindex{buffered mode send!nonblocking} nonblocking send.



\cdeclindex{MPI\_Request}
\begin{mpi-binding}
    function_name("MPI_Issend")
    parameter(
        "buf",
        "BUFFER",
        constant=True,
        asynchronous=True,
        desc=r"initial address of send buffer",
    )
    parameter(
        "count",
        "POLYXFER_NUM_ELEM_NNI",
        desc=r"number of elements in send buffer",
    )
    parameter(
        "datatype", "DATATYPE", desc=r"datatype of each send buffer element"
    )
    parameter("dest", "RANK", desc=r"rank of destination")
    parameter("tag", "TAG", desc=r"message tag")
    parameter("comm", "COMMUNICATOR")
    parameter("request", "REQUEST", direction="out")
\end{mpi-binding}

Start\mpitermindex{send!start} a synchronous mode\mpitermindex{synchronous mode send!nonblocking} nonblocking send.

\cdeclindex{MPI\_Request}
\begin{mpi-binding}
    function_name("MPI_Irsend")
    parameter(
        "buf",
        "BUFFER",
        constant=True,
        asynchronous=True,
        desc=r"initial address of send buffer",
    )
    parameter(
        "count",
        "POLYXFER_NUM_ELEM_NNI",
        desc=r"number of elements in send buffer",
    )
    parameter(
        "datatype", "DATATYPE", desc=r"datatype of each send buffer element"
    )
    parameter("dest", "RANK", desc=r"rank of destination")
    parameter("tag", "TAG", desc=r"message tag")
    parameter("comm", "COMMUNICATOR")
    parameter("request", "REQUEST", direction="out")
\end{mpi-binding}

Start\mpitermindex{send!start} a ready mode\mpitermindex{ready mode send!nonblocking} nonblocking send.


\cdeclindex{MPI\_Request}
\begin{mpi-binding}
    function_name("MPI_Irecv")
    parameter(
        "buf",
        "BUFFER",
        direction="out",
        asynchronous=True,
        suppress="f08_intent",
        desc=r"initial address of receive buffer",
    )
    parameter(
        "count",
        "POLYXFER_NUM_ELEM_NNI",
        desc=r"number of elements in receive buffer",
    )
    parameter(
        "datatype", "DATATYPE", desc=r"datatype of each receive buffer element"
    )
    parameter(
        "source", "RANK", desc=r"rank of source or \mpiconst{MPI_ANY_SOURCE}"
    )
    parameter("tag", "TAG", desc=r"message tag or \mpiconst{MPI_ANY_TAG}")
    parameter("comm", "COMMUNICATOR")
    parameter("request", "REQUEST", direction="out")
\end{mpi-binding}

Start\mpitermindex{receive!start} a nonblocking receive\mpitermindex{receive!nonblocking}.

\cdeclindex{MPI\_Request}
\begin{mpi-binding}
    function_name("MPI_Isendrecv")

    parameter(
        "sendbuf",
        "BUFFER",
        constant=True,
        asynchronous=True,
        desc=r"initial address of send buffer",
    )
    parameter(
        "sendcount",
        "POLYXFER_NUM_ELEM_NNI",
        desc=r"number of elements in send buffer",
    )
    parameter(
        "sendtype", "DATATYPE", desc=r"datatype of each send buffer element"
    )
    parameter("dest", "RANK", desc=r"rank of destination")
    parameter("sendtag", "TAG", desc=r"send tag")

    parameter(
        "recvbuf",
        "BUFFER",
        direction="out",
        asynchronous=True,
        suppress="f08_intent",
        desc=r"initial address of receive buffer",
    )
    parameter(
        "recvcount",
        "POLYXFER_NUM_ELEM_NNI",
        desc=r"number of elements in receive buffer",
    )
    parameter(
        "recvtype", "DATATYPE", desc=r"datatype of each receive buffer element"
    )
    parameter(
        "source", "RANK", desc=r"rank of source or \mpiconst{MPI_ANY_SOURCE}"
    )
    parameter(
        "recvtag", "TAG", desc=r"receive tag or \mpiconst{MPI_ANY_TAG}"
    )

    parameter("comm", "COMMUNICATOR")
    parameter("request", "REQUEST", direction="out")
\end{mpi-binding}

Initiate\mpitermindex{send-receive!start} a nonblocking communication request for a \mpitermni{send and receive}\mpitermindex{send-receive!nonblocking} operation.

\cdeclindex{MPI\_Request}
\begin{mpi-binding}
    function_name("MPI_Isendrecv_replace")

    parameter(
        "buf",
        "BUFFER",
        direction="inout",
        asynchronous=True,
        desc=r"initial address of send and receive buffer",
    )
    parameter(
        "count",
        "POLYXFER_NUM_ELEM_NNI",
        desc=r"number of elements in send and receive buffer",
    )
    parameter(
        "datatype", "DATATYPE", desc=r"type of elements in send and receive buffer"
    )

    parameter("dest", "RANK", desc=r"rank of destination")
    parameter("sendtag", "TAG", desc=r"send message tag")
    parameter( "source", "RANK", desc=r"rank of source or \mpiconst{MPI_ANY_SOURCE}")
    parameter( "recvtag", "TAG", desc=r"receive message tag or \mpiconst{MPI_ANY_TAG}")

    parameter("comm", "COMMUNICATOR")
    parameter("request", "REQUEST", direction="out")
\end{mpi-binding}

Initiate\mpitermindex{send-receive!start} a nonblocking communication request for a \mpitermni{send and receive}\mpitermindex{send-receive!nonblocking} operation.
The same buffer is used both for the send and for the receive, so that the
message sent is replaced by the message received.

These calls allocate a communication
request object and associate it with the request handle (the argument
\mpiarg{request}).
The request can be used later to
query the status of the communication or wait for its completion.

A nonblocking send call indicates that the
system may start copying data out of the send buffer.
The sender should not
modify
any part of the send
buffer after a nonblocking send operation is called, until the send completes.

A nonblocking receive call indicates that the system may start
writing data into the receive buffer.  The receiver should not access
any part of the
receive buffer after a nonblocking receive operation is called,
until the receive completes.

%%
%% See comments elsewhere in this file as to why the commented out language
%% is both incorrect and unnecessary.
\begin{users}
  To prevent problems with the argument copying and register optimization done
  by Fortran compilers, please note the hints in
Sections~\ref{sec:misc-problems}--\ref{sec:f90-problems:comparison-with-C}.
%% especially in
%% Sections~\ref{sec:misc-sequence} and~\ref{sec:f90-problems:vector-subscripts}
%% on pages~\pageref{sec:misc-sequence}-\pageref{sec:f90-problems:vector-subscripts}
%% about ``{\sf Problems Due to Data Copying and Sequence Association with Subscript Triplets}''
%% and ``{\sf Vector Subscripts}'',
%% and in Sections~\ref{sec:misc-register} to~\ref{sec:f90-problems:perm-data-movements}
%% on pages~\pageref{sec:misc-register} to~\pageref{sec:f90-problems:perm-data-movements}
%% about ``{\sf Optimization Problems}'', ``{\sf Code Movements and Register Optimization}'',
%% ``{\sf Temporary Data Movements}'' and ``{\sf Permanent Data Movements}''.
\end{users}


\subsection{Communication Completion}
\mpitermtitleindexmainsub{nonblocking}{completion}
\mpitermtitleindex{semantics!nonblocking completion}
\label{subsec:pt2pt-commend}

The functions \mpifunc{MPI\_WAIT} and \mpifunc{MPI\_TEST} are used to complete a
nonblocking communication.  The \mpiterm{completion} of a send operation indicates that
the sender is now free
to update the send buffer
(the send operation itself leaves the content
of the send buffer unchanged). It does not indicate that the
message has been received,
rather, it may have been buffered by the communication
subsystem.  However, if a \mpiterm{synchronous mode send}
was used, the \mpiterm{completion} of the
send operation indicates that a matching receive was \mpitermindex{initiation}\mpitermni{initiated}, and that the
message will eventually be received by this matching receive.

The \mpiterm{completion} of a receive operation indicates that the receive buffer
contains the received message, the receiver is now free to access it,
and that the status object is set.  It does
not indicate that the matching send operation has \mpitermni{completed} (but indicates, of
course, that the send was \mpitermindex{initiation}\mpitermni{initiated}).

We shall use the following terminology:
A \mpitermdef{null handle}\mpitermindex{request objects!null handle} is a handle with
value\flushline
\mpiconstmain{MPI\_REQUEST\_NULL}.
A \mpitermni{persistent communication request}\mpitermindex{persistent communication request}
and the handle to it are \mpitermdef{inactive}\mpitermdefindex{persistent communication request!inactive}
if the request is not associated with any ongoing
communication (see \sectionref{sec:pt2pt-persistent}).
A handle is \mpitermdef{active}\mpitermdefindex{persistent communication request!active} if it is neither \mpitermni{null} nor \mpiterm{inactive}.
An
\mpitermdef{empty}\mpitermdefindex{status!empty} status is a status that is set to return
\mpiarg{tag}\mpicode{ = }\mpiconst{MPI\_ANY\_TAG},
\mpiarg{source}\mpicode{ = }\mpiconst{MPI\_ANY\_SOURCE},
\mpiarg{error}\mpicode{ = }\error{MPI\_SUCCESS},
and is also internally configured so that calls to
\mpifunc{MPI\_GET\_COUNT} and \mpifunc{MPI\_GET\_ELEMENTS} return
\mpiarg{count}\mpicode{ = 0}
and \mpifunc{MPI\_TEST\_CANCELLED} returns \mpicode{false}.
We set a status variable to \mpiterm{empty} when the value returned by it is not
significant. Status is set in this
way so as to prevent errors due to accesses of stale information.

The fields in a \mpiarg{status} object returned by a call to
\mpifunc{MPI\_WAIT}, \mpifunc{MPI\_TEST}, or any of the other derived functions
(\mpicode{MPI\_}\{\mpicode{TEST$|$WAIT}\}\{\mpicode{ALL$|$SOME$|$ANY}\}),
where the \mpiarg{request} corresponds to a send call\mpitermindex{status!for send operation}, are undefined, with two
exceptions: The error status field will contain valid information if the wait
or test call returned with \error{MPI\_ERR\_IN\_STATUS}; and the returned
status can be queried by the call \mpifunc{MPI\_TEST\_CANCELLED}.

Error codes belonging to the error class \error{MPI\_ERR\_IN\_STATUS} should
be returned only by the \MPI/ completion functions that take arrays of
\mpictype{MPI\_Status}.
For the functions that take a single \mpictype{MPI\_Status} argument, the
error code is returned by the function, and the value of the \mpiconst{MPI\_ERROR}
field in the \mpictype{MPI\_Status} argument is undefined (see~\ref{subsec:pt2pt-status}).

\cdeclmainindex{MPI\_Request}\cdeclmainindex{TYPE(MPI\_Request)}
\cdeclindex{MPI\_Status}
\begin{mpi-binding}
    function_name("MPI_Wait")
    parameter("request", "REQUEST", desc=r"request", direction="inout")
    parameter("status", "STATUS", direction="out")
\end{mpi-binding}

A call to \mpifunc{MPI\_WAIT} returns when the operation
identified by \mpiarg{request} is \mpitermindex{completion}\mpitermni{complete}.  If the request is an \mpiterm{active} \mpiterm{persistent communication request},
it is marked \mpiterm{inactive}.  Any other type of request
is deallocated
and the request handle is set to \mpiconst{MPI\_REQUEST\_NULL}.
\mpifunc{MPI\_WAIT} is in general a \mpiterm{nonlocal} procedure.
When the operation represented by the \mpiarg{request} is
\mpitermni{enabled}\mpitermindex{enabled operation state}\mpitermindex{state!operation!enabled}\mpitermindex{operation!state!enabled}
then a call to \mpifunc{MPI\_WAIT} is a \mpiterm{local} procedure call.

The call returns, in \mpiarg{status}, information on
the completed operation.  The content of the status object for a receive
operation can be accessed as
described in Section~\ref{subsec:pt2pt-status}.
The status object for a send operation may be queried by
a call to \mpifunc{MPI\_TEST\_CANCELLED}
(see Section~\ref{sec:pt2pt-probe}).

One is allowed to call \mpifunc{MPI\_WAIT} with a \mpitermni{null} or \mpitermni{inactive}
\mpiarg{request} argument.
In this case the procedure returns immediately with \mpiterm{empty}\mpitermindex{status!empty} \mpiarg{status}.

\begin{users}
Successful return of \mpifunc{MPI\_WAIT} after a \mpifunc{MPI\_IBSEND} implies
that the user send buffer can be reused---i.e., data has been sent
out or copied into a
buffer attached with \mpifunc{MPI\_BUFFER\_ATTACH}, \mpifunc{MPI\_COMM\_ATTACH\_BUFFER} or \mpifunc{MPI\_SESSION\_ATTACH\_BUFFER}.
Further, at this point,
we can no longer \mpiterm{cancel}\mpitermindex{message!cancel}\mpitermindex{semantics!point-to-point communication!cancel} the send (see Section~\ref{sec:pt2pt-probe}).
If a matching receive is never
\mpitermni{started}, then the buffer cannot be freed.  This runs somewhat
counter to the stated goal of \mpifunc{MPI\_CANCEL} (always being able to free
program space that was committed to the communication subsystem).
\end{users}

\begin{implementors}
In a multithreaded environment, a call to
\mpifunc{MPI\_WAIT} should block only the calling thread, allowing the thread
scheduler to schedule another thread for execution.
\end{implementors}

\cdeclindex{MPI\_Request}
\cdeclindex{MPI\_Status}
\begin{mpi-binding}
    function_name("MPI_Test")
    parameter("request", "REQUEST", direction="inout")
    parameter(
        "flag",
        "LOGICAL",
        desc=r"\mpicode{true} if operation completed",
        direction="out",
    )
    parameter("status", "STATUS", direction="out")
\end{mpi-binding}

A call to \mpifunc{MPI\_TEST} returns \mpiarg{flag}\mpicode{ = true} if the
operation identified by \mpiarg{request} is \mpitermindex{completion}\mpitermni{complete}.  In such a case, the
status object is set to contain information on the completed
operation.  If the request is an \mpiterm{active} \mpiterm{persistent communication request},
it is marked as \mpiterm{inactive}.  Any other type of request
is deallocated and the request handle is set to
\mpiconst{MPI\_REQUEST\_NULL}. The call returns
\mpiarg{flag}\mpicode{ = false} if the operation identified by
\mpiarg{request} is not complete.
In this case, the value
of the status object is undefined.
\mpifunc{MPI\_TEST} is a \mpiterm{local} procedure.

The return status object for a receive operation carries information that
can be accessed as described in
Section~\ref{subsec:pt2pt-status}.
The status object for a send operation carries information that can be
accessed by
a call to \mpifunc{MPI\_TEST\_CANCELLED}
(see Section~\ref{sec:pt2pt-probe}).

One
is allowed to call \mpifunc{MPI\_TEST} with a \mpitermni{null} or \mpiterm{inactive} \mpiarg{request}
argument. In such a case the procedure returns with \mpiarg{flag}\mpicode{ = true} and
\mpiterm{empty}\mpitermindex{status!empty} \mpiarg{status}.

The procedures \mpifunc{MPI\_WAIT} and \mpifunc{MPI\_TEST} can be used to
complete any request-based nonblocking\mpitermindex{nonblocking!completion} or persistent\mpitermindex{persistent communication request!completion} operation.

\begin{users}
The use of
the nonblocking \mpifunc{MPI\_TEST} call allows the user to
schedule alternative activities within a single thread of execution.
An event-driven thread scheduler can be emulated with periodic calls to
\mpifunc{MPI\_TEST}.
\end{users}


\begin{example}
\label{pt2pt-exJ}
\Exindex{Fortran}{Nonblocking point-to-point}{MPI\_Isend,MPI\_Irecv,MPI\_Wait}%
\exindex{Nonblocking operations}%
Simple usage of nonblocking operations and \mpifunc{MPI\_WAIT}.
%%HEADER
%%LANG: FORTRAN
%%FRAGMENT
%%DECL: integer comm, rank, ierr, tag, request
%%DECL: real a(10)
%%DECL: integer status(MPI_STATUS_SIZE)
%%SUBST:\*\*.*:tag=1
%%ENDHEADER
\begin{lstlisting}[language={[MPI]Fortran}]
CALL MPI_COMM_RANK(comm, rank, ierr)
IF (rank .EQ. 0) THEN
   CALL MPI_ISEND(a(1), 10, MPI_REAL, 1, tag, comm, request, ierr)
   ! **** do some computation to mask latency ****
   CALL MPI_WAIT(request, status, ierr)
ELSE IF (rank .EQ. 1) THEN
   CALL MPI_IRECV(a(1), 15, MPI_REAL, 0, tag, comm, request, ierr)
   ! **** do some computation to mask latency ****
   CALL MPI_WAIT(request, status, ierr)
END IF
\end{lstlisting}
\end{example}

A request object can be \mpitermni{freed}\mpitermindex{request objects!freeing} using the following \MPI/ procedure.

\cdeclindex{MPI\_Request}
\begin{mpi-binding}
    function_name("MPI_Request_free")
    parameter("request", "REQUEST", direction="inout")
\end{mpi-binding}

\mpifunc{MPI\_REQUEST\_FREE} is a \mpiterm{local} procedure. Upon successful return,
\mpifunc{MPI\_REQUEST\_FREE} sets \mpiarg{request} to \mpiconst{MPI\_REQUEST\_NULL}.
For an \mpiterm{inactive} \mpiargshort{request} representing any type of \MPI/ operation,
\mpifunc{MPI\_REQUEST\_FREE} shall do the \mpiterm{freeing stage} of the associated
operation during its execution.

For a \mpiarg{request} representing a \mpiterm{nonblocking} point-to-point
or a \mpitermindex{persistent communication request}\mpitermni{persistent} point-to-point operation, it is permitted (although
strongly discouraged) to call \mpifunc{MPI\_REQUEST\_FREE} when the
\mpiarg{request} is \mpiterm{active}. In this special case, \mpifunc{MPI\_REQUEST\_FREE}
will only mark the request for freeing and \MPI/ will actually do the
\mpiterm{freeing stage} of the operation associated with the request later.

The use of this procedure for generalized requests is described in
\sectionref{sec:ei-gr}.

Calling \mpifunc{MPI\_REQUEST\_FREE} with an \mpiterm{active} \mpiarg{request} representing any
other type of \MPI/ operation (e.g., any partitioned operation (see Chapter~\ref{chap:part}),
any collective operation (see Chapter~\ref{chap:coll}), any I/O
operation (see Chapter~\ref{chap:io-2}), or any
request-based \RMA/ operation (see Chapter~\ref{chap:one-side-2})) is \mpitermindex{erroneous program!freeing active request}\mpitermindex{semantics!erroneous program!freeing active request}\mpitermni{erroneous}.

\begin{rationale}
% restrict the next sentence to point-to-point...
For point-to-point operations, the \mpifunc{MPI\_REQUEST\_FREE} mechanism is provided for
reasons of performance and convenience on the sending side.
\end{rationale}

\begin{users}
Once a request is freed by a call to \mpifunc{MPI\_REQUEST\_FREE}, it is not possible
to check for the successful completion of the associated communication
with calls to \mpifunc{MPI\_WAIT} or \mpifunc{MPI\_TEST}. Also, if an error occurs subsequently
during the communication, an error code cannot be returned to the
user---such an error must be treated as \mpitermindex{error handling!fatal after request free}fatal. An active receive request should
never be freed as the receiver will have no way to verify that the receive
has completed and the receive buffer can be reused.
\end{users}

\begin{example}
\label{pt2pt-exK}
\Exindex{Fortran}{Nonblocking send and receive with request free}{MPI\_Request\_free,MPI\_Isend,MPI\_Irecv,MPI\_Wait}%
\exindex{Nonblocking operations}%
An example using \mpifunc{MPI\_REQUEST\_FREE}.

%%HEADER
%%LANG: FORTRAN
%%FRAGMENT
%%DECL: integer rank, req, ierr, outval, inval, i, n
%%DECL: integer status(MPI_STATUS_SIZE)
%%ENDHEADER
\begin{lstlisting}[language={[MPI]Fortran},basicstyle=\lstsmall]
CALL MPI_COMM_RANK(MPI_COMM_WORLD, rank, ierr)
IF (rank .EQ. 0) THEN
   DO i=1,n
      CALL MPI_ISEND(outval, 1, MPI_REAL, 1, 0, MPI_COMM_WORLD, req, ierr)
      CALL MPI_REQUEST_FREE(req, ierr)
      CALL MPI_IRECV(inval, 1, MPI_REAL, 1, 0, MPI_COMM_WORLD, req, ierr)
      CALL MPI_WAIT(req, status, ierr)
   END DO
ELSE IF (rank .EQ. 1) THEN
   CALL MPI_IRECV(inval, 1, MPI_REAL, 0, 0, MPI_COMM_WORLD, req, ierr)
   CALL MPI_WAIT(req, status, ierr)
   DO I=1,n-1
      CALL MPI_ISEND(outval, 1, MPI_REAL, 0, 0, MPI_COMM_WORLD, req, ierr)
      CALL MPI_REQUEST_FREE(req, ierr)
      CALL MPI_IRECV(inval, 1, MPI_REAL, 0, 0, MPI_COMM_WORLD, req, ierr)
      CALL MPI_WAIT(req, status, ierr)
   END DO
   CALL MPI_ISEND(outval, 1, MPI_REAL, 0, 0, MPI_COMM_WORLD, req, ierr)
   CALL MPI_WAIT(req, status, ierr)
END IF
\end{lstlisting}
\end{example}

\subsection{Semantics of Nonblocking Communication Operations}
\mpitermtitleindex{semantics!nonblocking communication operations}
\label{subsec:pt2pt-semantics}



The semantics of nonblocking communication operations are defined by suitably extending the
definitions in Section~\ref{sec:pt2pt-semantics}.

\paraheading{Order.}
Nonblocking communication operations are \mpitermdef{ordered}\mpitermdefindex{semantics!point-to-point communication!ordered} according to the execution
order of the calls that \mpitermni{initiate}\mpitermindex{initiation}\mpitermindex{request objects!initiation} the communication.
The \mpitermdef{nonovertaking}\mpitermdefindex{semantics!point-to-point communication!nonovertaking}
requirement of Section~\ref{sec:pt2pt-semantics} is extended to
nonblocking communication, with this definition of order being used.

\begin{example}
\label{pt2pt-exL}
\Exindex{Fortran}{Message ordering for nonblocking operations}{MPI\_Isend,MPI\_Irecv,MPI\_Wait}%
\exindex{Nonblocking operations!message ordering}%
Message ordering for nonblocking operations.
%%HEADER
%%LANG: FORTRAN
%%FRAGMENT
%%DECL: integer comm, rank, ierr, tag, a(2), b(2), r1, r2
%%DECL: integer status(MPI_STATUS_SIZE)
%%ENDHEADER
\begin{lstlisting}[language={[MPI]Fortran}]
CALL MPI_COMM_RANK(comm, rank, ierr)
IF (rank .EQ. 0) THEN
   CALL MPI_ISEND(a, 1, MPI_REAL, 1, 0, comm, r1, ierr)
   CALL MPI_ISEND(b, 1, MPI_REAL, 1, 0, comm, r2, ierr)
ELSE IF (rank .EQ. 1) THEN
   CALL MPI_IRECV(a, 1, MPI_REAL, 0, MPI_ANY_TAG, comm, r1, ierr)
   CALL MPI_IRECV(b, 1, MPI_REAL, 0, 0, comm, r2, ierr)
END IF
CALL MPI_WAIT(r1, status, ierr)
CALL MPI_WAIT(r2, status, ierr)
\end{lstlisting}
The first send will match the
first receive, even
if both messages are sent before either receive is executed.
\end{example}

\paraheading{Progress.}\mpitermdefindex{progress!point-to-point communication}\mpitermdefindex{semantics!point-to-point communication!progress}
A call to \mpifunc{MPI\_WAIT} that \mpitermni{completes}\mpitermindex{completion}\mpitermindex{request objects!completion} a receive will eventually
terminate and return if a matching send has been \mpitermni{started}\mpitermindex{request objects!started}\mpitermindex{started!request objects}, unless
the send is satisfied by another receive.  In particular, if the matching send
is
\mpitermni{nonblocking},
then the receive should \mpitermni{complete} even if no call is executed by the
sender to \mpitermni{complete} the send.  Similarly, a call to \mpifunc{MPI\_WAIT} that
\mpitermni{completes} a send will eventually
return if a matching receive has been \mpitermni{started}, unless
the receive is satisfied by another send, and even if no call is executed to
\mpitermni{complete} the receive.


\begin{example}
\label{pt2pt-exM}
\Exindex{Fortran}{Progress semantics}{MPI\_Send,MPI\_Ssend,MPI\_Recv,MPI\_Irecv,MPI\_Wait}%
\exindex{Nonblocking operations!progress}%
An illustration of progress semantics.
%%HEADER
%%LANG: FORTRAN
%%FRAGMENT
%%DECL: integer comm, rank, ierr, tag, a(2), b(2), r
%%DECL: integer status(MPI_STATUS_SIZE)
%%ENDHEADER
\begin{lstlisting}[language={[MPI]Fortran}]
CALL MPI_COMM_RANK(comm, rank, ierr)
IF (rank .EQ. 0) THEN
   CALL MPI_SSEND(a, 1, MPI_REAL, 1, 0, comm, ierr)
   CALL MPI_SEND(b, 1, MPI_REAL, 1, 1, comm, ierr)
ELSE IF (rank .EQ. 1) THEN
   CALL MPI_IRECV(a, 1, MPI_REAL, 0, 0, comm, r, ierr)
   CALL MPI_RECV(b, 1, MPI_REAL, 0, 1, comm, status, ierr)
   CALL MPI_WAIT(r, status, ierr)
END IF
\end{lstlisting}

This code should not deadlock\mpitermindex{deadlock avoidance!nonblocking communication} in a correct \MPI/ implementation.
The first synchronous send must complete once the
matching (nonblocking) receive is \mpitermni{started}, even though
the completing wait call has not yet been reached.
Thus, the sending \MPI/ process will continue and execute
the second send procedure, allowing the receiving \MPI/ process to complete execution.
\end{example}

If an \mpifunc{MPI\_TEST} that \mpitermni{completes} a receive is repeatedly called
with the same
arguments, and
a matching send has been \mpitermni{started}, then the call will eventually
return \mpiarg{flag}\mpicode{ = true}, unless the send is satisfied by another
receive.
If an \mpifunc{MPI\_TEST} that \mpitermni{completes} a send is repeatedly called with the
same arguments, and
a matching receive has been \mpitermni{started}, then the call will eventually
return \mpiarg{flag}\mpicode{ = true}, unless the receive is satisfied by another
send.
See also \sectionref{subsec:terms:progress} on \mpitermni{progress}.

\subsection{Multiple Completions}
\mpitermtitleindex{multiple completions}
\mpitermtitleindex{completion!multiple}
\mpitermtitleindex{request objects!multiple completions}
\label{subsec:pt2pt-multiple}

It is convenient to be able to wait for the \mpitermni{completion} of any, some, or all the
operations in a list, rather than having to wait for a specific message.
A call to \mpifunc{MPI\_WAITANY} or \mpifunc{MPI\_TESTANY} can be used to
wait for the
\mpitermni{completion} of one out of several operations. A call to \mpifunc{MPI\_WAITALL}
or \mpifunc{MPI\_TESTALL} can be
used to wait for all \mpitermni{pending}\mpitermindex{pending (communication) operation}\mpitermindex{communication!pending}\mpitermindex{operation!pending} operations in a list. A call to
\mpifunc{MPI\_WAITSOME} or \mpifunc{MPI\_TESTSOME} can be used to \mpitermni{complete} all
enabled operations in a list.

\cdeclindex{MPI\_Request}
\cdeclindex{MPI\_Status}
\begin{mpi-binding}
    function_name("MPI_Waitany")
    parameter("count", "ARRAY_LENGTH_NNI", desc=r"list length")
    parameter(
        "array_of_requests",
        "REQUEST",
        desc=r"array of requests",
        direction="inout",
        length="count",
    )
    parameter(
        "index",
        "INDEX",
        desc=r"index of handle for operation that completed",
        direction="out",
    )
    parameter("status", "STATUS", direction="out")
\end{mpi-binding}

Does not return until 
one of the operations associated with the \mpitermni{active}
requests in the array has \mpitermni{completed}\mpitermindex{completion!multiple}\mpitermindex{multiple completions}.
If more than one operation is
enabled and can \mpitermni{complete}, one is arbitrarily chosen.
Returns in \mpiarg{index} the index
of that request in the array and returns in \mpiargshort{status} the status of the
completing operation.
(The array is indexed from zero in C, and from one in Fortran.)
If the request
is an \mpitermni{active} \mpiterm{persistent communication request}, it is marked \mpitermni{inactive}.  Any other type of
request
is deallocated and the request handle is set to \mpiconst{MPI\_REQUEST\_NULL}.

The \mpiarg{array\_of\_requests} list may contain \mpitermni{null}\mpitermindex{null handle}\mpitermindex{request objects!null handle} or \mpitermni{inactive}
handles.
If the list contains no \mpitermni{active} handles (list has length zero or all
entries are \mpitermni{null} or \mpiterm{inactive}),
then the call returns immediately with
\mpiarg{index}\mpicode{ = }\mpiconst{MPI\_UNDEFINED}, and an \mpiterm{empty}\mpitermindex{status!empty} \mpiarg{status}.

The execution of
\mpifunc{MPI\_WAITANY}
with an array containing multiple entries
has the same effect as the execution of
\mpifunc{MPI\_WAIT}
with the array entry indicated by the output value of \mpiarg{index}
(unless the output value of \mpiarg{index} is \mpiconst{MPI\_UNDEFINED}).
\mpifunc{MPI\_WAITANY} with an array containing one \mpitermni{active} entry
is equivalent to \mpifunc{MPI\_WAIT}.

\cdeclindex{MPI\_Request}
\cdeclindex{MPI\_Status}
\begin{mpi-binding}
    function_name("MPI_Testany")
    parameter("count", "ARRAY_LENGTH_NNI", desc=r"list length")
    parameter(
        "array_of_requests",
        "REQUEST",
        desc=r"array of requests",
        direction="inout",
        length="count",
    )
    parameter(
        "index",
        "INDEX",
        desc=(
            r"index of operation that completed or \mpiconst{MPI_UNDEFINED} if "
            r"none completed"
        ),
        direction="out",
    )
    parameter(
        "flag",
        "LOGICAL",
        desc=r"\mpicode{true} if one of the operations is complete",
        direction="out",
    )
    parameter("status", "STATUS", direction="out")
\end{mpi-binding}

Tests for \mpitermni{completion}\mpitermindex{completion!multiple}\mpitermindex{multiple completions} of
either one or none of the operations associated with \mpitermni{active} handles.
In the former case, it returns \mpiarg{flag}\mpicode{ = true},
returns in \mpiarg{index} the index of this request in the array,
and returns in \mpiarg{status} the status of that operation.  If
the request is an \mpitermni{active} \mpiterm{persistent communication request}, it is marked as \mpiterm{inactive}.
Any other type of request
is deallocated
and the handle is set to \mpiconst{MPI\_REQUEST\_NULL}.
(The array is indexed from zero in C, and from one in Fortran.)
In the latter case (no operation \mpitermni{completed}),
it returns \mpiarg{flag}\mpicode{ = false}, returns a value
of \mpiconst{MPI\_UNDEFINED} in \mpiarg{index} and \mpiarg{status} is
undefined.

The array may contain \mpitermni{null}\mpitermindex{null handle}\mpitermindex{request objects!null handle} or inactive handles.
If the
array contains no \mpitermni{active} handles then the call returns
\mpitermni{immediately}\mpitermindex{immediate} with \mpiarg{flag}\mpicode{ = true},
\mpiarg{index}\mpicode{ = }\mpiconst{MPI\_UNDEFINED}, and an \mpitermni{empty}\mpitermindex{status!empty} \mpiarg{status}.

If the array of requests contains \mpitermni{active} handles then
the execution of
\mpifunc{MPI\_TESTANY}
has the same effect as the execution of
\mpifunc{MPI\_TEST}
with each of the \mpitermni{active} handles in the array
in some arbitrary order, until one call returns \mpiarg{flag}\mpicode{ = true}, or
all return \mpiarg{flag}\mpicode{ = false}.  In the former case, \mpiarg{index} is set
to indicate which array element returned \mpiarg{flag}\mpicode{ = true}
and in the latter case, it is set to \mpiconst{MPI\_UNDEFINED}.
\mpifunc{MPI\_TESTANY} with an array containing one \mpitermni{active} entry
is equivalent to \mpifunc{MPI\_TEST}.

\cdeclindex{MPI\_Request}
\cdeclindex{MPI\_Status}
\begin{mpi-binding}
    function_name("MPI_Waitall")
    parameter("count", "ARRAY_LENGTH_NNI", desc=r"list length")
    parameter(
        "array_of_requests",
        "REQUEST",
        desc=r"array of requests",
        direction="inout",
        length="count",
    )
    parameter(
        "array_of_statuses",
        "STATUS",
        desc=r"array of status objects",
        direction="out",
        length="*",
        pointer=False,
    )
\end{mpi-binding}


Does not return until 
all communication operations associated with \mpitermni{active} handles
in the list \mpitermni{complete}\mpitermindex{completion!multiple}\mpitermindex{multiple completions}, and returns the status of all these operations
(this includes the case where no handle in the list is \mpitermni{active}).
Both arrays have the same number of valid entries.  The \code{i}-th entry in
\mpiarg{array\_of\_statuses} is set to the return status of the
\code{i}-th operation.
\mpitermni{Active} \mpitermni{persistent requests}\mpitermindex{persistent communication request} are marked \mpitermni{inactive}.  Requests of any other type
are
deallocated and the corresponding handles in the array are set to
\mpiconst{MPI\_REQUEST\_NULL}.
The list may contain \mpitermni{null}\mpitermindex{null handle}\mpitermindex{request objects!null handle} or \mpitermni{inactive} handles.
The call sets to \mpitermni{empty}\mpitermindex{status!empty} the status of each such entry.

The error-free execution of
\mpifunc{MPI\_WAITALL}
has the same effect as the execution of
\mpifunc{MPI\_WAIT}
for each of the array elements
in some arbitrary order.
\mpifunc{MPI\_WAITALL} with an array of length one
is equivalent to \mpifunc{MPI\_WAIT}.

When one or more of the communication operations \mpitermni{completed} by a
call to \mpifunc{MPI\_WAITALL} fail, it is
desirable to return specific information on each
communication.  The function \mpifunc{MPI\_WAITALL} will return in such
case the error code \error{MPI\_ERR\_IN\_STATUS} and will set the
error field of each status to a specific error code.  This code will be
\error{MPI\_SUCCESS}, if the specific communication \mpitermni{completed}; it will
be another specific error code, if it failed;
or it can be \errormain{MPI\_ERR\_PENDING} if it has neither failed nor \mpitermni{completed}.
The function \mpifunc{MPI\_WAITALL} will return \error{MPI\_SUCCESS} if no request
had an error,
or will return another error code if it failed
for other reasons (such as invalid arguments).  In such cases, it will
not update the error fields of the statuses.

\begin{rationale}
This design streamlines error handling\mpitermindex{multiple completions!error handling}\mpitermindex{error handling!multiple completions} in the application.
The application code need only test the (single) function result to
determine if an error has occurred.  It needs to check each individual
status only when an error occurred.
\end{rationale}

\cdeclindex{MPI\_Request}
\cdeclindex{MPI\_Status}
\begin{mpi-binding}
    function_name("MPI_Testall")
    parameter("count", "ARRAY_LENGTH_NNI", desc=r"list length")
    parameter(
        "array_of_requests",
        "REQUEST",
        desc=r"array of requests",
        direction="inout",
        length="count",
    )
    parameter("flag", "LOGICAL", desc=r"\mpicode{true} if all of the operations are complete", direction="out")
    parameter(
        "array_of_statuses",
        "STATUS",
        desc=r"array of status objects",
        direction="out",
        length="*",
        pointer=False,
    )
\end{mpi-binding}

Returns \mpiarg{flag}\mpicode{ = true}
if all communication operations associated
with \mpitermni{active} handles in the array have \mpitermni{completed}\mpitermindex{completion!multiple}\mpitermindex{multiple completions} (this includes the
case where no handle in the list is \mpitermni{active}).
In this case, each status entry that corresponds to an
\mpitermni{active} request
is set to the status of the corresponding
operation.  \mpitermni{Active} \mpitermni{persistent requests}\mpitermindex{persistent communication request} are marked \mpitermni{inactive}.
Requests of any other type are deallocated and the corresponding handles in
the array are
set to \mpiconst{MPI\_REQUEST\_NULL}.
Each status entry that corresponds to a \mpitermni{null}\mpitermindex{null handle}\mpitermindex{request objects!null handle} or \mpitermni{inactive}
handle is set to \mpitermni{empty}\mpitermindex{status!empty}.

Otherwise,
\mpiarg{flag}\mpicode{ = false} is returned, no request is modified
and the values of the status entries are undefined.
This is a \mpiterm{local} procedure.

Errors that occurred during the execution of \mpifunc{MPI\_TESTALL}
are handled in the same manner as errors in \mpifunc{MPI\_WAITALL}.

\cdeclindex{MPI\_Request}
\cdeclindex{MPI\_Status}
\begin{mpi-binding}
    function_name("MPI_Waitsome")
    parameter("incount", "ARRAY_LENGTH_NNI", desc=r"length of array_of_requests")
    parameter(
        "array_of_requests",
        "REQUEST",
        desc=r"array of requests",
        direction="inout",
        length="incount",
    )
    parameter(
        "outcount",
        "ARRAY_LENGTH",
        desc=r"number of completed requests",
        direction="out",
    )
    parameter(
        "array_of_indices",
        "INDEX",
        desc=r"array of indices of operations that completed",
        direction="out",
        length="*",
    )
    parameter(
        "array_of_statuses",
        "STATUS",
        desc=r"array of status objects for operations that completed",
        direction="out",
        length="*",
        pointer=False,
    )
\end{mpi-binding}

Does not return until
at least one of the operations associated with \mpitermni{active}
handles in the list have \mpitermni{completed}\mpitermindex{completion!multiple}\mpitermindex{multiple completions}.
Returns in \mpiarg{outcount} the number of requests from the list
\mpiarg{array\_of\_requests} that have \mpitermni{completed}.  Returns in the first
\mpiarg{outcount} locations of the array \mpiarg{array\_of\_indices}
the indices of these operations (index within the
array \mpiarg{array\_of\_requests}; the array is indexed from zero in
C and from one in
Fortran).  Returns in the first \mpiarg{outcount}
locations of the array \mpiarg{array\_of\_statuses}
the status for these \mpitermni{completed} operations.
\mpitermni{Completed} \mpitermni{active} \mpitermni{persistent requests}\mpitermindex{persistent communication request} are marked as \mpitermni{inactive}.  Any other
type or request that \mpitermni{completed}
is deallocated, and the
associated handle is set to \mpiconst{MPI\_REQUEST\_NULL}.

If the list contains no \mpitermni{active} handles, then the
call returns \mpitermni{immediately}\mpitermindex{immediate} with \mpiarg{outcount}\mpicode{ = }\mpiconst{MPI\_UNDEFINED}.

When one or more of the communication operations \mpitermni{completed} by
\mpifunc{MPI\_WAITSOME} fails, then it is desirable to return specific
information on each communication.
The arguments \mpiarg{outcount},
\mpiarg{array\_of\_indices} and \mpiarg{array\_of\_statuses} will be
adjusted to indicate \mpitermni{completion} of all communication operations that have
succeeded or failed.  The call will return the error code
\error{MPI\_ERR\_IN\_STATUS} and the error field of each status
returned will be set to indicate success or to indicate the specific error
that occurred.  The call will return \error{MPI\_SUCCESS} if no request
resulted in an error,
and will return another error code if it failed
for other reasons (such as invalid arguments).  In such cases, it will
not update the error fields\mpitermindex{error handling!multiple completions}\mpitermindex{multiple completions!error handling} of the statuses.

\cdeclindex{MPI\_Request}
\cdeclindex{MPI\_Status}
\begin{mpi-binding}
    function_name("MPI_Testsome")
    parameter("incount", "ARRAY_LENGTH_NNI", desc=r"length of array_of_requests")
    parameter(
        "array_of_requests",
        "REQUEST",
        desc=r"array of requests",
        direction="inout",
        length="incount",
    )
    parameter(
        "outcount",
        "ARRAY_LENGTH",
        desc=r"number of completed requests",
        direction="out",
    )
    parameter(
        "array_of_indices",
        "INDEX",
        desc=r"array of indices of operations that completed",
        direction="out",
        length="*",
    )
    parameter(
        "array_of_statuses",
        "STATUS",
        desc=r"array of status objects for operations that completed",
        direction="out",
        length="*",
        pointer=False,
    )
\end{mpi-binding}

This procedure behaves like \mpifunc{MPI\_WAITSOME}, except that it returns
\mpitermni{immediately}\mpitermindex{immediate}. If no operation has completed it
returns \mpiarg{outcount}\mpicode{ = 0}.
If there is no \mpitermni{active} handle in the list it
returns \mpiarg{outcount}\mpicode{ = }\mpiconst{MPI\_UNDEFINED}.

\mpifunc{MPI\_TESTSOME} is a \mpiterm{local} procedure, which returns
\mpitermni{immediately}\mpitermindex{immediate}, whereas
\mpifunc{MPI\_WAITSOME} will
not return until
a communication \mpitermni{completes}, if it was
passed a list that contains at least one \mpitermni{active} handle.  Both calls fulfill a
\mpitermdefni{fairness requirement}\mpitermdefindex{fairness!requirement}\mpitermdefindex{semantics!point-to-point communication!fairness requirement}: If a request for a receive repeatedly
appears in a list of requests passed to \mpifunc{MPI\_WAITSOME} or
\mpifunc{MPI\_TESTSOME}, and a matching send has been \mpitermni{started}, then the receive
will eventually succeed, unless the send is satisfied by another receive; and
similarly for send requests.

Errors that occur during the execution of \mpifunc{MPI\_TESTSOME} are
handled as for \linebreak \mpifunc{MPI\_WAITSOME}.

\begin{users}
The use of \mpifunc{MPI\_TESTSOME} is likely to be more efficient than the use
of \mpifunc{MPI\_TESTANY}. The former returns information on all
\mpitermni{completed} communication operations, with the latter, a new call is required for
each communication that completes.

A server with multiple clients can use \mpifunc{MPI\_WAITSOME} so as not to
starve any client.  Clients send messages to the server with service
requests. The server calls \mpifunc{MPI\_WAITSOME} with one receive request
for each client, and then handles all receives that completed.
If a call to \mpifunc{MPI\_WAITANY} is used instead, then one client
could starve while requests from another client always sneak in first.
\end{users}

\begin{implementors}
\mpifunc{MPI\_TESTSOME} should \mpitermni{complete} as many
\mpitermni{pending}\mpitermindex{pending (communication) operation}\mpitermindex{communication!pending}\mpitermindex{operation!pending} communication operations of the \mpiarg{array\_of\_requests} as possible.
\end{implementors}



\begin{example}
\label{pt2pt-exN}
\Exindex{Fortran}{Client-server}{MPI\_Isend,MPI\_Irecv,MPI\_Wait,MPI\_Waitany}%
\exindex{Client-server code}%
Client-server code (starvation can occur).
%%HEADER
%%LANG: FORTRAN90
%%FRAGMENT
%%DECL: integer comm, rank, ierr, tag, a(2,2), b(2)
%%DECL: integer n, size, i, request_list(10), index
%%DECL: integer status(MPI_STATUS_SIZE), request
%%DECL: interface
%%DECL: subroutine do_service(a)
%%DECL: integer a
%%DECL: end subroutine do_service
%%DECL: end interface
%%ENDHEADER
\begin{lstlisting}[language={[MPI]Fortran}]
CALL MPI_COMM_SIZE(comm, size, ierr)
CALL MPI_COMM_RANK(comm, rank, ierr)
IF (rank .GT. 0) THEN         ! client code
   DO WHILE(.TRUE.)
      CALL MPI_ISEND(a, n, MPI_REAL, 0, tag, comm, request, ierr)
      CALL MPI_WAIT(request, status, ierr)
   END DO
ELSE         ! rank=0 -- server code
   DO i=1,size-1
      CALL MPI_IRECV(a(1,i), n, MPI_REAL, i, tag, &
                     comm, request_list(i), ierr)
   END DO
   DO WHILE(.TRUE.)
      CALL MPI_WAITANY(size-1, request_list, index, status, ierr)
      CALL DO_SERVICE(a(1,index))  ! handle one message
      CALL MPI_IRECV(a(1, index), n, MPI_REAL, index, tag, &
                     comm, request_list(index), ierr)
   END DO
END IF
\end{lstlisting}
\end{example}
\label{page:pt2pt-waitany}

\begin{example}
\label{pt2pt-exO}
\Exindex{Fortran}{Client server code with waitsome}{MPI\_Isend,MPI\_Irecv,MPI\_Wait,MPI\_Waitsome}%
\exindex{Client-server code}%
Same code, using \mpifunc{MPI\_WAITSOME}.
%%HEADER
%%LANG: FORTRAN
%%FRAGMENT
%%DECL: integer comm, rank, ierr, tag, a(2,2), b(2)
%%DECL: integer status(MPI_STATUS_SIZE), numdone
%%DECL: integer request_list(10),i, n, request
%%DECL: integer indices(10), size
%%DECL: integer statuses(MPI_STATUS_SIZE,10)
%%DECL: interface
%%DECL: subroutine do_service(a)
%%DECL: integer a
%%DECL: end subroutine do_service
%%DECL: end interface
%%ENDHEADER
\begin{lstlisting}[language={[MPI]Fortran}]
CALL MPI_COMM_SIZE(comm, size, ierr)
CALL MPI_COMM_RANK(comm, rank, ierr)
IF (rank .GT. 0) THEN         ! client code
   DO WHILE(.TRUE.)
      CALL MPI_ISEND(a, n, MPI_REAL, 0, tag, comm, request, ierr)
      CALL MPI_WAIT(request, status, ierr)
   END DO
ELSE         ! rank=0 -- server code
   DO i=1,size-1
      CALL MPI_IRECV(a(1,i), n, MPI_REAL, i, tag, &
                     comm, request_list(i), ierr)
   END DO
   DO WHILE(.TRUE.)
      CALL MPI_WAITSOME(size, request_list, numdone, &
                        indices, statuses, ierr)
      DO i=1,numdone
         CALL DO_SERVICE(a(1, indices(i)))
         CALL MPI_IRECV(a(1, indices(i)), n, MPI_REAL, 0, tag, &
                        comm, request_list(indices(i)), ierr)
      END DO
   END DO
END IF
\end{lstlisting}
\end{example}



\subsection{Non-Destructive Test of \texorpdfstring{\mpiarg{status}}{status}}
\mpitermtitleindex{status!test}
\label{subsec:pt2pt-teststatus}

These procedures are useful for accessing the information associated with a
request, without \mpitermni{freeing}\mpitermindex{freeing stage} the request (in case the user is expected to access
it later).  It allows one to layer libraries more conveniently,
since multiple layers of software may access the same \mpitermni{completed} request and
extract from it the status information.

\cdeclindex{MPI\_Request}
\cdeclindex{MPI\_Status}
\begin{mpi-binding}
    function_name("MPI_Request_get_status")
    parameter("request", "REQUEST", desc=r"request")
    parameter(
        "flag",
        "LOGICAL",
        direction="out",
        desc=r"boolean flag, same as from \mpifunc{MPI_TEST}",
    )
    parameter(
        "status",
        "STATUS",
        direction="out",
        desc=r"status object if flag is true",
    )
\end{mpi-binding}

Sets \mpiarg{flag}\code{ = true} if the operation is \mpitermni{complete}, and, if so, returns in
status the request status.  However, unlike test or wait, it does not
deallocate or \mpitermni{inactivate} the request; a subsequent call to test, wait or free
must be executed with that request.  It sets \mpiarg{flag}\code{ = false} if the
operation is not \mpitermni{complete}.

One is allowed to call \mpifunc{MPI\_REQUEST\_GET\_STATUS} with a \mpitermni{null}\mpitermindex{null handle}\mpitermindex{request objects!null handle} or \mpitermni{inactive}
request argument. In such a case the procedure returns with \mpiarg{flag}\code{ = true} and \mpitermni{empty}\mpitermindex{status!empty} status.

The \mpitermni{progress}\mpitermindex{progress!point-to-point communication}\mpitermindex{semantics!point-to-point communication!progress}
rule for \mpifunc{MPI\_TEST}, as described in Section~\ref{subsec:pt2pt-semantics}, also applies to \mpifunc{MPI\_REQUEST\_GET\_STATUS}.

\cdeclindex{MPI\_Request}
\cdeclindex{MPI\_Status}
\begin{mpi-binding}
    function_name("MPI_Request_get_status_any")
    parameter("count", "ARRAY_LENGTH_NNI", desc=r"list length")
    parameter(
        "array_of_requests",
        "REQUEST",
        desc=r"array of requests",
        direction="in", constant=True,
        length="count",
    )
    parameter(
        "index",
        "INDEX",
        desc=(
	    r"index of operation that completed or \mpiconst{MPI_UNDEFINED} if "
	    r"none completed"
        ),
        direction="out",
    )
    parameter(
        "flag",
        "LOGICAL",
        direction="out",
        desc=r"\mpicode{true} if one of the operations is complete",
    )
    parameter(
        "status",
        "STATUS",
        direction="out",
        desc=r"status object if flag is true",
    )
\end{mpi-binding}

Tests for \mpitermni{completion}\mpitermindex{completion!multiple}\mpitermindex{multiple completions} of
either one or none of the operations associated with \mpitermni{active} handles.
In the former case, it returns \mpiarg{flag}\mpicode{ = true},
returns in \mpiarg{index} the index of this request in the array,
and returns in \mpiarg{status} the status of that operation.
(The array is indexed from zero in C, and from one in Fortran.)
In the latter case (no operation \mpitermni{completed}),
it returns \mpiarg{flag}\mpicode{ = false}, returns a value
of \mpiconst{MPI\_UNDEFINED} in \mpiarg{index} and \mpiarg{status} is
undefined.

The array may contain \mpitermni{null}\mpitermindex{null handle}\mpitermindex{request objects!null handle} or inactive handles.
If the
array contains no \mpitermni{active} handles then the call returns
\mpitermni{immediately}\mpitermindex{immediate} with \mpiarg{flag}\mpicode{ = true},
\mpiarg{index}\mpicode{ = }\mpiconst{MPI\_UNDEFINED}, and an \mpitermni{empty}\mpitermindex{status!empty} \mpiarg{status}.

If the array of requests contains active handles then the execution of \mpifunc{MPI\_REQUEST\_GET\_STATUS\_ANY}
has the same effect as the execution of \mpifunc{MPI\_REQUEST\_GET\_STATUS} with each of the active
array elements in some arbitrary order, until one call returns \mpiarg{flag}\mpicode{ = true}, or all return \mpiarg{flag}\mpicode{ = false}. In the
former case, \mpiarg{index} is set to indicate which array element returned \mpiarg{flag}\mpicode{ = true} and in the latter
case, it is set to \mpiconst{MPI\_UNDEFINED}. \mpifunc{MPI\_REQUEST\_GET\_STATUS\_ANY} with an array containing
one request is equivalent to \mpifunc{MPI\_REQUEST\_GET\_STATUS}.

\cdeclindex{MPI\_Request}
\cdeclindex{MPI\_Status}
\begin{mpi-binding}
    function_name("MPI_Request_get_status_all")
    parameter("count", "ARRAY_LENGTH_NNI", desc=r"list length")
    parameter(
        "array_of_requests",
        "REQUEST",
        desc=r"array of requests",
        direction="in", constant=True,
        length="count"
    )
    parameter(
        "flag",
        "LOGICAL",
        direction="out",
        desc=r"true if all of the operations are complete",
    )
    parameter(
        "array_of_statuses",
        "STATUS",
        direction="out",
        length="*",
        pointer=False,
        desc=r"array of status objects",
    )
\end{mpi-binding}

\mpifunc{MPI\_REQUEST\_GET\_STATUS\_ALL} returns \mpiarg{flag}\mpicode{ = true}
if all communication operations associated
with \mpitermni{active} handles in the array have \mpitermni{completed}\mpitermindex{completion!multiple}\mpitermindex{multiple completions} (this includes the
case where all handles in the list are \mpitermni{inactive} or \mpiconst{MPI\_REQUEST\_NULL}).
In this case, each status entry that corresponds to an
\mpitermni{active} request
is set to the status of the corresponding
operation. Unlike test or wait, it does not
deallocate or \mpitermni{inactivate} the requests; a subsequent call to test, wait or free
should be executed with each of those requests.

Each status entry that corresponds to a \mpitermni{null}\mpitermindex{null handle}\mpitermindex{request objects!null handle} or \mpitermni{inactive}
handle is set to \mpitermni{empty}\mpitermindex{status!empty}.

Otherwise,
\mpiarg{flag}\mpicode{ = false} is returned and the values of the status entries are undefined.

The \mpitermni{progress}\mpitermindex{progress!point-to-point communication}\mpitermindex{semantics!point-to-point communication!progress}
rule for \mpifunc{MPI\_TEST}, as described in Section~\ref{subsec:pt2pt-semantics}, also applies to \mpifunc{MPI\_REQUEST\_GET\_STATUS\_ALL}.


\cdeclindex{MPI\_Request}
\cdeclindex{MPI\_Status}
\begin{mpi-binding}
    function_name("MPI_Request_get_status_some")
    parameter("incount", "ARRAY_LENGTH_NNI", desc=r"length of array_of_requests")
    parameter(
        "array_of_requests",
        "REQUEST",
        desc=r"array of requests",
        direction="in", constant=True,
        length="incount",
    )
    parameter(
        "outcount",
        "ARRAY_LENGTH",
        desc=r"number of completed requests",
        direction="out",
    )
    parameter(
        "array_of_indices",
        "INDEX",
        desc=r"array of indices of operations that completed",
        direction="out",
        length="*",
    )
    parameter(
        "array_of_statuses",
        "STATUS",
        direction="out",
        desc=r"array of status objects for operations that completed",
        length="*",
        pointer=False
    )
\end{mpi-binding}

\mpifunc{MPI\_REQUEST\_GET\_STATUS\_SOME} returns in \mpiarg{outcount} the number of requests from the list
\mpiarg{array\_of\_requests} that have \mpitermni{completed}.  Returns in the first
\mpiarg{outcount} locations of the array \mpiarg{array\_of\_indices}
the indices of these operations within the
array \mpiarg{array\_of\_requests}; the array is indexed from zero in
C and from one in
Fortran.  Returns in the first \mpiarg{outcount}
locations of the array \mpiarg{array\_of\_statuses}
the status for these \mpitermni{completed} operations.
However, unlike test or wait, it does not
deallocate or \mpitermni{inactivate} any requests in \mpiarg{array\_of\_requests};
a subsequent call to test, wait or free
should be executed with each completed request. If no operation in \mpiarg{array\_of\_requests}
is complete, it returns \mpiarg{outcount}\mpicode{ = 0}. If all operations in \mpiarg{array\_of\_requests} are either
\mpiconst{MPI\_REQUEST\_NULL} or \mpitermni{inactive}, \mpiarg{outcount} will be set to \mpiconst{MPI\_UNDEFINED}.
The \mpitermni{progress}\mpitermindex{progress!point-to-point communication}\mpitermindex{semantics!point-to-point communication!progress}
rule for \mpifunc{MPI\_TEST}, as described in Section~\ref{subsec:pt2pt-semantics}, also applies to \mpifunc{MPI\_REQUEST\_GET\_STATUS\_SOME}.

% This procedure behaves like \mpifunc{MPI\_TESTSOME}, except that it does not free or inactivate any request.
% Testsome has this additional text, which is not necessary after above sentence?
% If no operation has completed it
% returns \mpiarg{outcount}\mpicode{ = 0}.
% If there is no \mpitermni{active} handle in the list it
% returns \mpiarg{outcount}\mpicode{ = }\mpiconst{MPI\_UNDEFINED}.

Like \mpifunc{MPI\_WAITSOME} and \mpifunc{MPI\_TESTSOME},
\mpifunc{MPI\_REQUEST\_GET\_STATUS\_SOME} fulfills a
\mpitermdefni{fairness requirement}\mpitermdefindex{fairness!requirement}\mpitermdefindex{semantics!point-to-point communication!fairness requirement}: If a request for a receive repeatedly
appears in a list of requests passed to \mpifunc{MPI\_REQUEST\_GET\_STATUS\_SOME},
\mpifunc{MPI\_WAITSOME}, or \linebreak\mpifunc{MPI\_TESTSOME} and
a matching send has been \mpitermni{started}, then the receive
will eventually succeed, unless the send is satisfied by another receive; and
similarly for send requests.

%Errors that occur during the execution of \mpifunc{MPI\_REQUEST\_GET\_STATUS\_SOME} are
%handled as for \mpifunc{MPI\_WAITSOME}.

\begin{implementors}
\mpifunc{MPI\_REQUEST\_GET\_STATUS\_SOME} should \mpitermni{complete} as many pending communication operations as
possible.
\end{implementors}

\begin{users}
\mpifunc{MPI\_REQUEST\_GET\_STATUS\_ANY}, \vlgb\mpifunc{MPI\_REQUEST\_GET\_STATUS\_SOME}, and \mpifunc{MPI\_REQUEST\_GET\_STATUS\_ALL}
offer tradeoffs between precision and speed, as do the corrsponding TEST and WAIT functions. The ANY variants are fast, but imprecise and unfair.
The ALL variants will provide all-or-nothing information and/or completion, which can limit their applicability. The SOME variants, because
of their precision and fairness guarantee, will typically be the slowest on a per-call basis.
\end{users}

\section{Probe and Cancel}
\label{sec:pt2pt-probe}

The \mpifunc{MPI\_PROBE}, \mpifunc{MPI\_IPROBE},
\mpifunc{MPI\_MPROBE}, and \mpifunc{MPI\_IMPROBE}
procedures allow incoming messages to be checked for,
without
actually receiving them.  The user can then decide how to receive them, based
on the information returned by the \mpitermdef{probe} (basically, the information
returned by \mpiarg{status}).  In particular, the user may allocate memory for
the receive buffer, according to the length of the probed message.

The \mpifunc{MPI\_CANCEL} procedure allows
\mpitermni{pending}\mpitermindex{pending (communication) operation}\mpitermindex{communication!pending}\mpitermindex{operation!pending}
communication operations to be
\mpitermdef{cancelled}\mpitermdefindex{cancel}\mpitermdefindex{message!cancel}\mpitermdefindex{semantics!point-to-point communication!cancel}.
This
is required for cleanup.  Posting a send or a receive ties up user resources
(send or receive buffers), and a \mpitermni{cancel} may be needed to free these resources
gracefully.

\mpitermni{Cancelling}\mpitermindex{cancel}\mpitermindex{message!cancel}\mpitermindex{semantics!point-to-point communication!cancel} a send request by calling \mpifunc{MPI\_CANCEL} is deprecated.
\mpitermni{Cancelling} a send-recv request by calling \mpifunc{MPI\_CANCEL} is not allowed.

\subsection{Probe}
\mpitermtitleindex{probe}
\mpitermtitleindex{point-to-point communication!probe}
\mpitermtitleindex{semantics!point-to-point communication!probe}
\label{subsec:pt2pt-probe}

\cdeclindex{MPI\_Status}
\begin{mpi-binding}
    function_name("MPI_Iprobe")
    parameter(
        "source", "RANK", desc=r"rank of source or \mpiconst{MPI_ANY_SOURCE}"
    )
    parameter("tag", "TAG", desc=r"message tag or \mpiconst{MPI_ANY_TAG}")
    parameter("comm", "COMMUNICATOR")
    parameter("flag", "LOGICAL", desc=r"\mpicode{true} if there is a matching message that can be received", direction="out")
    parameter("status", "STATUS", direction="out")
\end{mpi-binding}

\mpifunc{MPI\_IPROBE}
returns \mpiarg{flag}\mpicode{ = true}
if there is a message that can be received
and that matches the pattern specified by the
arguments \mpiarg{source}, \mpiarg{tag}, and \mpiarg{comm}.
The call matches the same message
that would have been received by a call to \mpifunc{MPI\_RECV} with the same argument values for \mpiarg{source}, \mpiarg{tag},
\mpiarg{comm}, and \mpiarg{status} executed at the same point in the program, and returns in
\mpiarg{status} the same
value that would have been returned by \mpifunc{MPI\_RECV}.
Otherwise, the call returns \mpiarg{flag}\mpicode{ = false}, and leaves \mpiarg{status}
undefined.

If \mpifunc{MPI\_IPROBE} returns \mpiarg{flag}\mpicode{ = true},
then the content of the status object can be subsequently
accessed as described in Section~\ref{subsec:pt2pt-status} to find the
source, tag, and length of the probed message.

\mpifunc{MPI\_IPROBE} is a \mpiterm{local} procedure since its return does not depend
on \MPI/ calls in other \MPI/ processes, which is marked with the prefix \mpicode{I} (for \mpiterm{immediate}).

A subsequent receive executed with the same communicator, and the source and
tag returned in status by \mpifunc{MPI\_IPROBE} will receive the message that
was matched\mpitermindex{matching rules!probe} by the probe, if no other intervening receive occurs after the
probe, and the send is not successfully \mpiterm{cancelled}\mpitermindex{cancel}\mpitermindex{message!cancel}\mpitermindex{semantics!point-to-point communication!cancel} before the receive.
If the receiving \MPI/ process is multithreaded, it is the user's
responsibility to ensure that the last condition holds.

The \mpiarg{source} argument of \mpifunc{MPI\_IPROBE} can be
\mpiconst{MPI\_ANY\_SOURCE}, and the \mpiargshort{tag} argument can be
\mpiconst{MPI\_ANY\_TAG}, so that one can \mpiterm{probe} for \mpitermni{messages} from an arbitrary
source and/or with
an arbitrary tag.  However, a specific communication context
must be provided with the \mpiarg{comm} argument.

It is not necessary to receive a message immediately after it has been
probed for, and the
same message may be probed for several times before it is received.

A probe with \mpiconst{MPI\_PROC\_NULL} as source returns \mpiarg{flag}\mpicode{ = true}, and the status object returns \mpiarg{source}\mpicode{ = }\mpiconst{MPI\_PROC\_NULL},
\mpiarg{tag}\mpicode{ = }\mpiconst{MPI\_ANY\_TAG}, and \mpiarg{count}\mpicode{ = 0};
see \sectionref{sec:pt2pt-nullproc}.

\cdeclindex{MPI\_Status}
\begin{mpi-binding}
    function_name("MPI_Probe")
    parameter(
        "source", "RANK", desc=r"rank of source or \mpiconst{MPI_ANY_SOURCE}"
    )
    parameter("tag", "TAG", desc=r"message tag or \mpiconst{MPI_ANY_TAG}")
    parameter("comm", "COMMUNICATOR")
    parameter("status", "STATUS", direction="out")
\end{mpi-binding}

\mpifunc{MPI\_PROBE} behaves like \mpifunc{MPI\_IPROBE} except that it is a \mpiterm{nonlocal}
call that returns only after a matching message has been found.

The \MPI/ implementations of \mpifunc{MPI\_PROBE} and \mpifunc{MPI\_IPROBE} need
to guarantee \mpitermni{progress}\mpitermindex{progress!point-to-point communication}\mpitermindex{probe!progress}\mpitermindex{semantics!point-to-point communication!progress}:
if a call to \mpifunc{MPI\_PROBE} has been issued by an \MPI/ process, and a send that
matches the probe has been \mpitermni{initiated}\mpitermindex{initiation} by some \MPI/ process, then the call to
\mpifunc{MPI\_PROBE} will return, unless the message is received by another
concurrent receive operation (that is executed by another thread at the probing
\MPI/ process).

Similarly, if an \MPI/ process repeatedly calls
\mpifunc{MPI\_IPROBE} and a matching message has been issued,
then
\mpifunc{MPI\_IPROBE} will eventually return \mpiarg{flag}\mpicode{ = true}
unless the message is received by another concurrent receive
operation or matched by a concurrent \mpiterm{matching probe}.
See also \sectionref{subsec:terms:progress} on \mpitermni{progress}.

\begin{example}
\label{pt2pt-exP}
\Exindex{Fortran}{Client-server with probe}{MPI\_Send,MPI\_Recv,MPI\_Probe}%
\exindex{Client-server code!with probe}%
Use probe to wait for an incoming message.

%%HEADER
%%LANG: FORTRAN
%%FRAGMENT
%%DECL: integer comm, rank, ierr, tag, i, x
%%DECL: integer status(MPI_STATUS_SIZE)
%%ENDHEADER
\begin{lstlisting}[language={[MPI]Fortran}]
    CALL MPI_COMM_RANK(comm, rank, ierr)
    IF (rank .EQ. 0) THEN
       CALL MPI_SEND(i, 1, MPI_INTEGER, 2, 0, comm, ierr)
    ELSE IF (rank .EQ. 1) THEN
       CALL MPI_SEND(x, 1, MPI_REAL, 2, 0, comm, ierr)
    ELSE IF (rank .EQ. 2) THEN
       DO i=1,2
          CALL MPI_PROBE(MPI_ANY_SOURCE, 0, comm, status, ierr)
          IF (status(MPI_SOURCE) .EQ. 0) THEN
100          CALL MPI_RECV(i, 1, MPI_INTEGER, 0, 0, comm, status, ierr)
          ELSE
200          CALL MPI_RECV(x, 1, MPI_REAL, 1, 0, comm, status, ierr)
          END IF
       END DO
    END IF
\end{lstlisting}
Each message is received with the right type.
\end{example}


\begin{example}
\label{pt2pt-exQ}
\Exindex{Fortran}{Client-server (with error)}{MPI\_Send,MPI\_Recv,MPI\_Probe}%
\exindex{Client-server code!with probe (wrong)}%
A similar program to the previous example, but now it
has a problem.

%%HEADER
%%LANG: FORTRAN
%%FRAGMENT
%%DECL: integer comm, rank, ierr, i, x
%%DECL: integer status(MPI_STATUS_SIZE)
%%ENDHEADER
\begin{lstlisting}[language={[MPI]Fortran}]
! ---------------- THIS EXAMPLE IS ERRONEOUS ---------------
    CALL MPI_COMM_RANK(comm, rank, ierr)
    IF (rank .EQ. 0) THEN
       CALL MPI_SEND(i, 1, MPI_INTEGER, 2, 0, comm, ierr)
    ELSE IF (rank .EQ. 1) THEN
       CALL MPI_SEND(x, 1, MPI_REAL, 2, 0, comm, ierr)
    ELSE IF (rank .EQ. 2) THEN
       DO i=1,2
          CALL MPI_PROBE(MPI_ANY_SOURCE, 0, comm, status, ierr)
          IF (status(MPI_SOURCE) .EQ. 0) THEN
100          CALL MPI_RECV(i, 1, MPI_INTEGER, MPI_ANY_SOURCE, &
                           0, comm, status, ierr)
          ELSE
200          CALL MPI_RECV(x, 1, MPI_REAL, MPI_ANY_SOURCE, &
                           0, comm, status, ierr)
          END IF
       END DO
    END IF
\end{lstlisting}

In Example~\ref{pt2pt-exQ},
the two receive calls in statements labeled 100 and 200 in
Example~\ref{pt2pt-exP} are
slightly modified,
using \mpiconst{MPI\_ANY\_SOURCE} as the \mpiarg{source}
argument.
The program is now incorrect: the receive operation may receive a message that
is distinct from the message probed by the preceding call to
\mpifunc{MPI\_PROBE}.
\end{example}

\begin{users}
In a multithreaded \MPI/ program, \mpifunc{MPI\_PROBE} and
\mpifunc{MPI\_IPROBE} might need special care.
If a thread \mpitermni{probes} for a message and then immediately posts a matching
receive, the receive may match a message other than that found by the probe since another thread could concurrently receive that original
message~\cite{mprobe-paper}.
\mpifunc{MPI\_MPROBE} and \mpifunc{MPI\_IMPROBE} solve this problem by
matching the incoming message so that it may only be received with
\mpifunc{MPI\_MRECV} or \mpifunc{MPI\_IMRECV} on the corresponding \mpiterm{message handle}\mpitermindex{message!handle}.
\end{users}

\begin{implementors}
A call to \mpifunc{MPI\_PROBE} will match the
message that would have been received by a call to \mpifunc{MPI\_RECV}
with the same argmument values for \mpiarg{source}, \mpiarg{tag}, \mpiarg{comm}, and \mpiarg{status} executed at the same point.
Suppose
that this message has source \mpicode{s}, tag \mpicode{t} and communicator
\mpicode{c}.  If the
tag argument in the probe call has value \mpiconst{MPI\_ANY\_TAG},
then the message probed
will be the earliest\mpitermindex{matching rules!probe} pending message from source \mpicode{s}
with communicator \mpicode{c} and any tag; in any case, the message probed will be the
earliest pending message from source \mpicode{s} with tag \mpicode{t} and
communicator
\mpicode{c} (this is
the message that would have been received, so as to preserve message order).
This message continues as the earliest pending message
from source \mpicode{s} with tag \mpicode{t} and communicator \mpicode{c}, until it is received.
A receive operation subsequent to the probe that uses the same communicator as the
probe and uses the tag and source values returned by the probe, must receive
this message, unless it has already been received by another receive operation.
\end{implementors}

\subsection{Matching Probe}
\mpitermtitleindex{matching probe}
\mpitermtitleindex{probe!matching}
\mpitermtitleindex{point-to-point communication!matching probe}
\mpitermtitleindex{semantics!point-to-point communication!matching probe}
\label{sec:matching-probe}

The function \mpifunc{MPI\_PROBE} checks for incoming \mpitermni{messages} without
receiving them. Since the list of incoming \mpitermni{messages} is global among the
threads of each \MPI/ process, it can be hard to use this functionality
in threaded environments \cite{mprobe-paper,mprobe-proposal-rev4}.

Like \mpifunc{MPI\_PROBE} and \mpifunc{MPI\_IPROBE}, the
\mpitermdef{matching probe} operation (%
\mpifunc{MPI\_MPROBE} and \mpifunc{MPI\_IMPROBE} procedures) allow incoming
\mpitermni{messages} to be queried without actually receiving them, except that
\mpifunc{MPI\_MPROBE} and \mpifunc{MPI\_IMPROBE} provide a mechanism to receive the specific
\mpiterm{message} that was matched regardless of other intervening probe or
receive operations. This gives the application an opportunity to decide
how to receive the message, based on the information returned by
the probe. In particular, the user may allocate memory for the receive
buffer, according to the length of the probed message.

\cdeclindex{MPI\_Status}
\cdeclmainindex{MPI\_Message}\cdeclmainindex{TYPE(MPI\_Message)}
\begin{mpi-binding}
    function_name("MPI_Improbe")
    parameter(
        "source", "RANK", desc=r"rank of source or \mpiconst{MPI_ANY_SOURCE}"
    )
    parameter("tag", "TAG", desc=r"message tag or \mpiconst{MPI_ANY_TAG}")
    parameter("comm", "COMMUNICATOR")
    parameter("flag", "LOGICAL", desc=r"\mpicode{true} if there is a matching message that can be received", direction="out")
    parameter("message", "MESSAGE", desc=r"returned message", direction="out")
    parameter("status", "STATUS", direction="out")
\end{mpi-binding}

\mpifunc{MPI\_IMPROBE} returns
\mpiarg{flag}\mpicode{ = true} if there is a message that can be received and that
matches the pattern specified by the arguments \mpiarg{source},
\mpiarg{tag}, and \mpiarg{comm}. The call matches the same message that
would have been received by a call to \mpifunc{MPI\_RECV} with the same argument values for \mpiarg{source}, \mpiarg{tag},
\mpiarg{comm}, and \mpiarg{status} executed at the same point in the program and returns
in \mpiarg{status} the same value that would have been returned by
\mpifunc{MPI\_RECV}. In addition, it returns in \mpiarg{message} a \mpitermdef{message handle}\mpitermdefindex{message!handle}
to the matched message.  Otherwise, the call returns \mpiarg{flag}\mpicode{ = false},
and leaves \mpiarg{status} and \mpiarg{message} undefined.

\mpifunc{MPI\_IMPROBE} is a \mpiterm{local} procedure.
According to the definitions in Section~\ref{subsec:procedures} and
in contrast to \mpifunc{MPI\_IPROBE}, it is a \mpiterm{nonblocking} procedure
because it is the \mpiterm{initialization} of a \mpiterm{matched receive} operation.

A \mpiterm{matched receive}\mpitermindex{receive!matched} (\mpifunc{MPI\_MRECV} or \mpifunc{MPI\_IMRECV}) executed with the \mpiterm{message handle}\mpitermindex{message!handle} will receive the
message that was matched by the \mpiterm{matching probe}. Unlike \mpifunc{MPI\_IPROBE}, no
other probe or receive operation may match the message returned by
\mpifunc{MPI\_IMPROBE}. Each \mpiterm{message handle} returned by \mpifunc{MPI\_IMPROBE} must
be received with either \mpifunc{MPI\_MRECV} or \mpifunc{MPI\_IMRECV}.

The \mpiarg{source} argument of \mpifunc{MPI\_IMPROBE} can be
\mpiconst{MPI\_ANY\_SOURCE}, and the \mpishortarg{tag} argument can be \mpiconst{MPI\_ANY\_TAG},
so that one can \mpiterm{probe} for \mpitermni{messages} from an arbitrary source and/or
with an arbitrary tag. However, a specific communication context must
be provided with the \mpiarg{comm} argument.

A \mpiterm{synchronous mode send} operation that is matched with \mpifunc{MPI\_IMPROBE}
or \mpifunc{MPI\_MPROBE} will \mpitermni{complete} successfully only if both a \mpiterm{matching receive}
is \mpitermni{started} with \mpifunc{MPI\_MRECV} or \mpifunc{MPI\_IMRECV},
and the \mpiterm{matching receive} operation has \mpitermni{started}\mpitermindex{started!request objects}\mpitermindex{request objects!started} to receive the message sent by the
\mpiterm{synchronous mode send}.

There is a special \mpitermdefni{predefined message handle}\mpitermdefindex{message handle!predefined}\mpitermdefindex{message!predefined handle}: \mpiconstmain{MPI\_MESSAGE\_NO\_PROC},
which is a message that has \mpiconst{MPI\_PROC\_NULL} as its source.
The predefined constant \mpiconstmain{MPI\_MESSAGE\_NULL} is the value used for \mpitermdefni{invalid message handles}\mpitermdefindex{message handle!invalid}\mpitermdefindex{message!invalid handle}.

A \mpiterm{matching probe} with \mpiarg{source}\mpicode{ = }\mpiconst{MPI\_PROC\_NULL} returns \mpiarg{flag}\mpicode{ = true},
\mpiarg{message}\mpicode{ = }\mpiconst{MPI\_MESSAGE\_NO\_PROC}, and the status object
returns \mpiarg{source}\mpicode{ = }\mpiconst{MPI\_PROC\_NULL}, \mpiarg{tag}\mpicode{ = }\mpiconst{MPI\_ANY\_TAG},
and \mpiarg{count}\mpicode{ = 0}; see Section~\ref{sec:pt2pt-nullproc}.
It is not necessary to call \mpifunc{MPI\_MRECV} or \mpifunc{MPI\_IMRECV}
with \mpiconst{MPI\_MESSAGE\_NO\_PROC}, but it is not \mpitermni{erroneous} to do so.

\begin{rationale}
\mpiconst{MPI\_MESSAGE\_NO\_PROC} was chosen instead of
\mpiconstskip{MPI\_MESSAGE\_PROC\_NULL} % \mpiconstskip is used instead of \mpiconstmain, because this is an invalid name, i.e., not an MPI constant
to avoid possible confusion as another null handle constant.
\end{rationale}

\cdeclindex{MPI\_Status}
\begin{mpi-binding}
    function_name("MPI_Mprobe")
    parameter(
        "source", "RANK", desc=r"rank of source or \mpiconst{MPI_ANY_SOURCE}"
    )
    parameter("tag", "TAG", desc=r"message tag or \mpiconst{MPI_ANY_TAG}")
    parameter("comm", "COMMUNICATOR")
    parameter("message", "MESSAGE", desc=r"returned message", direction="out")
    parameter("status", "STATUS", direction="out")
\end{mpi-binding}


\mpifunc{MPI\_MPROBE} behaves like \mpifunc{MPI\_IMPROBE} except that it is a
\mpiterm{blocking} call that returns only after a matching message has been
found.

The implementation of \mpifunc{MPI\_MPROBE} and \mpifunc{MPI\_IMPROBE} needs
to guarantee \mpitermni{progress}\mpitermindex{progress!point-to-point communication}\mpitermindex{matched probe!progress}\mpitermindex{semantics!point-to-point communication!progress} in the same way as in the case of
\mpifunc{MPI\_PROBE} and \mpifunc{MPI\_IPROBE}.
See also \sectionref{subsec:terms:progress} on \mpitermni{progress}.

According to the definitions in Section~\ref{subsec:procedures}, \mpifunc{MPI\_MPROBE}
is \mpiterm{incomplete}. It is also a \mpiterm{nonlocal} procedure.
\begin{users}
This is one of the exceptions\mpitermindex{semantics!exceptions!incomplete and nonlocal} in which \mpiterm{incomplete} procedures are \mpiterm{nonlocal}.
\end{users}

\subsection{Matched Receives}
\mpitermtitleindex{matched receive}
\mpitermtitleindex{receive!matched}
\mpitermtitleindex{point-to-point communication!matched receive}
\mpitermtitleindex{semantics!point-to-point communication!matched receive}
\label{sec:matched-receive}

The \mpitermdef{matched receive}\mpitermdefindex{receive!matched} operation (\mpifunc{MPI\_MRECV} and \mpifunc{MPI\_IMRECV} procedures) receive
\mpitermni{messages} that have been previously matched by a \mpiterm{matching probe} operation
(Section~\ref{sec:matching-probe}).

\cdeclindex{MPI\_Status}
\begin{mpi-binding}
    function_name("MPI_Mrecv")
    parameter(
        "buf",
        "BUFFER",
        desc=r"initial address of receive buffer",
        direction="out",
    )
    parameter(
        "count",
        "POLYXFER_NUM_ELEM_NNI",
        desc=r"number of elements in receive buffer",
    )
    parameter(
        "datatype", "DATATYPE", desc=r"datatype of each receive buffer element"
    )
    parameter("message", "MESSAGE", direction="inout", desc=r"message")
    parameter("status", "STATUS", direction="out")
\end{mpi-binding}

This call receives a message matched by a \mpiterm{matching probe} operation
(Section~\ref{sec:matching-probe}).

The \mpitermni{receive buffer}\mpitermindex{receive!buffer} consists of the storage containing \mpiarg{count}
consecutive elements of the type specified by \mpiarg{datatype},
starting at address \mpiarg{buf}.
The length of the received message
must be less than or equal to the length of the receive buffer. An
overflow error occurs if all incoming data does not fit, without
truncation, into the receive buffer.

If the message is shorter than the receive buffer, then only those
locations corresponding to the (shorter) message are modified.

On return from this function, the \mpiterm{message handle}\mpitermindex{message!handle} is set to
\mpiconst{MPI\_MESSAGE\_NULL}. All errors that occur during the execution of
this operation are handled according to the error handler set for the
communicator used in the matching probe call that produced the message
handle.

If \mpifunc{MPI\_MRECV} is called with
\mpiconst{MPI\_MESSAGE\_NO\_PROC}
as the
message argument, the call returns immediately with the status object
set to \mpiarg{source}\mpicode{ = }\mpiconst{MPI\_PROC\_NULL}, \mpiarg{tag}\mpicode{ = }\mpiconst{MPI\_ANY\_TAG}, and
\mpiarg{count}\mpicode{ = 0}.
This is consistent with the status object produced by
a call to \mpifunc{MPI\_RECV} or to \mpifunc{MPI\_PROBE}
with \mpiarg{source}\mpicode{ = }\mpiconst{MPI\_PROC\_NULL}
(see Section~\ref{sec:pt2pt-nullproc}).
A call to \mpifunc{MPI\_MRECV} with \mpiconst{MPI\_MESSAGE\_NULL} is \mpitermni{erroneous}\mpitermindex{erroneous program!invalid matched receive}\mpitermindex{semantics!erroneous program!invalid matched receive}.


\begin{mpi-binding}
    function_name("MPI_Imrecv")
    parameter(
        "buf",
        "BUFFER",
        direction="out",
        asynchronous=True,
        desc=r"initial address of receive buffer",
    )
    parameter(
        "count",
        "POLYXFER_NUM_ELEM_NNI",
        desc=r"number of elements in receive buffer",
    )
    parameter(
        "datatype", "DATATYPE", desc=r"datatype of each receive buffer element"
    )
    parameter("message", "MESSAGE", direction="inout", desc=r"message")
    parameter("request", "REQUEST", direction="out")
\end{mpi-binding}

\mpifunc{MPI\_IMRECV} is the nonblocking variant of \mpifunc{MPI\_MRECV}
and starts a nonblocking receive of a matched message. Completion
semantics are similar to \mpifuncshort{MPI\_IRECV} as described in
Section~\ref{subsec:pt2pt-commstart}. On return from this function,
the \mpiterm{message handle}\mpitermindex{message!handle} is set to \mpiconst{MPI\_MESSAGE\_NULL}.

If \mpifunc{MPI\_IMRECV} is called with \mpiconst{MPI\_MESSAGE\_NO\_PROC}
as the message argument, the call returns immediately with a request
object that, when completed, will yield a status object set to
\mpiarg{source}\mpicode{ = }\mpiconst{MPI\_PROC\_NULL}, \mpiarg{tag}\mpicode{ = }\mpiconst{MPI\_ANY\_TAG},
and \mpiarg{count}\mpicode{ = 0}, as if a receive from \mpiconst{MPI\_PROC\_NULL}
was issued (see Section~\ref{sec:pt2pt-nullproc}).
A call to \mpifunc{MPI\_IMRECV} with \mpiconst{MPI\_MESSAGE\_NULL} is \mpitermni{erroneous}\mpitermindex{erroneous program!invalid matched receive}\mpitermindex{semantics!erroneous program!invalid matched receive}.

\begin{implementors}
If reception of a matched message is started with
\mpifunc{MPI\_IMRECV}, then it is possible to \mpiterm{cancel}\mpitermindex{message!cancel}\mpitermindex{semantics!point-to-point communication!cancel} the returned
request with \mpifunc{MPI\_CANCEL}. If \mpifunc{MPI\_CANCEL} succeeds, the
matched message must be found by a subsequent message probe
(\mpifunc{MPI\_PROBE}, \mpifunc{MPI\_IPROBE}, \mpifunc{MPI\_MPROBE}, or
\mpifunc{MPI\_IMPROBE}), received by a subsequent receive operation or
\mpiterm{cancelled} by the sender.
See Section~\ref{sec:cancel} for details about \mpifunc{MPI\_CANCEL}.
The \mpitermni{cancellation} of operations initiated with \mpifunc{MPI\_IMRECV} may
fail.
\end{implementors}

\subsection{Cancel}
\mpitermtitleindex{cancel}
\mpitermtitleindex{point-to-point communication!cancel}
\mpitermtitleindex{semantics!point-to-point communication!cancel}
\label{sec:cancel}

\cdeclindex{MPI\_Request}
\begin{mpi-binding}
    function_name("MPI_Cancel")
    parameter("request", "REQUEST", direction="in", pointer=True)
\end{mpi-binding}

A call to \mpifunc{MPI\_CANCEL} marks for \mpitermni{cancellation} a \mpitermni{pending}\mpitermindex{pending (communication) operation}\mpitermindex{communication!pending}\mpitermindex{operation!pending},
\mpiterm{nonblocking} 
communication operation (send or receive).
\mpitermni{Cancelling} a send request by calling \mpifunc{MPI\_CANCEL} is deprecated.
The \mpiterm{cancel}\mpitermindex{message!cancel}\mpitermindex{semantics!point-to-point communication!cancel} call is
\mpiterm{local}. It returns \mpitermni{immediately}\mpitermindex{immediate}, possibly before the
communication is actually \mpitermni{cancelled}.
It is still necessary to call \mpifunc{MPI\_REQUEST\_FREE},
\mpifunc{MPI\_WAIT} or \mpifunc{MPI\_TEST} (or any of the derived procedures)
with the \mpiterm{cancelled} request as argument after the call to \mpifunc{MPI\_CANCEL}.
If a communication is marked for \mpitermni{cancellation}, then a \mpifunc{MPI\_WAIT}
call for that communication is guaranteed to return, irrespective of
the activities of other \MPI/ processes (i.e., \mpifunc{MPI\_WAIT} behaves as a
\mpiterm{local} function); similarly if \mpifunc{MPI\_TEST} is
repeatedly called for a \mpitermni{cancelled} communication,
then \mpifunc{MPI\_TEST} will eventually return \mpiarg{flag}\mpicode{ = true}.

\mpifunc{MPI\_CANCEL} can be used to \mpiterm{cancel}\mpitermindex{message!cancel}\mpitermindex{semantics!point-to-point communication!cancel} a communication that uses
a \mpiterm{persistent communication request} (see Section~\ref{sec:pt2pt-persistent}), in
the same way as it is described above for nonblocking operations.
\mpitermni{Cancelling} a persistent send request by calling \mpifunc{MPI\_CANCEL} is deprecated.
A successful \mpitermni{cancellation} \mpitermni{cancels}
the \mpiterm{active} communication, but not the request itself.  After the call to
\mpifunc{MPI\_CANCEL} and the subsequent call to \mpifunc{MPI\_WAIT} or
\mpifunc{MPI\_TEST}, the request becomes \mpiterm{inactive} and
can be activated for a new communication.

The successful
\mpitermni{cancellation} of a \mpiterm{buffered mode send} frees the buffer space occupied by
the pending message.
\mpitermni{Cancelling} a \mpiterm{buffered mode send} request by calling \mpifunc{MPI\_CANCEL} is deprecated.

Either the \mpitermni{cancellation} succeeds, or the communication succeeds, but
not both.
If a send is marked for \mpitermni{cancellation}, which is deprecated, then it must be the case that
either the send \mpitermni{completes} normally, in which case the
message sent was received at the destination, or that the send is
successfully
\mpiterm{cancelled}, in which case no part of the message was received at the
destination.  Then, any matching\mpitermindex{matching rules!cancel} receive has to be satisfied by another send.
If a receive is marked for \mpitermni{cancellation}, then it must be the case that
either the receive \mpitermni{completes} normally, or that the receive is
successfully \mpiterm{cancelled}, in which case no part of the receive buffer
is altered.  Then, any matching\mpitermindex{matching rules!cancel} send has to be satisfied by another receive.

If the operation has been
\mpiterm{cancelled}, then information to that effect will be returned in the
status argument of the operation that \mpitermni{completes} the communication.

\begin{rationale}
Although the \IN\ request handle parameter should not need to be passed
by reference, the C binding has listed the argument type as \mpictype{MPI\_Request}\code{*} since
\MPIIDOTO/. This function signature therefore cannot be changed without breaking
existing \MPI/ applications.
\end{rationale}

\cdeclindex{MPI\_Status}
\begin{mpi-binding}
    function_name("MPI_Test_cancelled")
    parameter("status", "STATUS", constant=True)
    parameter("flag", "LOGICAL", desc=r"\mpicode{true} if the operation has been cancelled", direction="out")
\end{mpi-binding}

Returns \mpiarg{flag}\mpicode{ = true} if the communication associated with the
status object was \mpiterm{cancelled} successfully.  In such a case, all
other fields of \mpiarg{status} (such as \mpiarg{count} or \mpiarg{tag}) are
undefined.  Returns \mpiarg{flag}\mpicode{ = false}, otherwise.  If a receive
operation might be \mpiterm{cancelled} then one should call \mpifunc{MPI\_TEST\_CANCELLED}
first, to check whether the operation was
\mpiterm{cancelled}, before checking on the other fields of the return status.

\begin{users}
\mpitermni{Cancel}\mpitermindex{cancel}\mpitermindex{message!cancel}\mpitermindex{semantics!point-to-point communication!cancel} can be an expensive operation that should be used only exceptionally.
\end{users}

\begin{implementors}
If a send operation uses an ``eager'' protocol (data is transferred to
the receiver
before a matching receive is \mpitermni{started}), then the \mpitermni{cancellation} of this send
may require communication with the intended receiver in order to free
allocated
buffers.  On some systems this may require an interrupt to the
intended receiver.
Note that, while communication may be needed to implement
\mpifunc{MPI\_CANCEL},
this is still a \mpiterm{local} procedure, since its completion does not
depend on the code executed by other \MPI/ processes.  If processing is required on
another \MPI/ process, this should be transparent to the application (hence the need
for an interrupt and an interrupt handler).
See also \sectionref{subsec:terms:progress} on \mpitermni{progress}\mpitermindex{progress!cancellation}.
\end{implementors}

\section{Persistent Communication Requests}
\mpitermtitleindex{persistent communication request}
\mpitermtitleindex{point-to-point communication!persistent}
\mpitermtitleindex{semantics!point-to-point communication!persistent}
\label{sec:pt2pt-persistent}

Often a communication with the same argument list (with the exception of the buffer contents) is repeatedly
executed within the inner loop of a parallel computation.  In such a
situation, it may be possible to optimize the communication by
binding the list of communication arguments to a \mpiterm{persistent communication request}
once and then repeatedly using
the request to \mpitermni{start}\mpitermindex{started!request objects}\mpitermindex{request objects!started}\mpitermindex{persistent communication request!starting} and \mpitermni{complete}\mpitermindex{completion}\mpitermindex{persistent communication request!completion} operations. %was 'messages'
In the case of point-to-point communication, the
\mpiterm{persistent communication request} thus created can be thought of as a
communication port or a ``half-channel.''
It does not provide the full functionality of a conventional channel,
since there is no binding of the send port to the receive port. This
construct allows the reduction of the overhead for communication
between the \MPI/ process and communication controller, but not of the overhead for
communication between one communication controller and another.
It is not necessary that messages sent with a persistent point-to-point request be received
by a receive operation using a persistent point-to-point request, or vice versa.

There are also persistent collective communication operations defined in
Section~\ref{sec:coll-persistent} and Section~\ref{sec:sparsecoll-pnbc}.
The remainder of this section covers the point-to-point
persistent \mpiterm{initialization} operations and the start\mpitermindex{starting procedure}
routines, which are used for persistent point-to-point, partitioned point-to-point, and
persistent collective communication operations.

A point-to-point \mpitermdef{persistent communication request} is created using one of the
five
following calls.  These point-to-point persistent \mpiterm{initialization} calls involve no communication.

\cdeclindex{MPI\_Request}
\begin{mpi-binding}
    function_name("MPI_Send_init")
    parameter(
        "buf",
        "BUFFER",
        constant=True,
        asynchronous=True,
        desc=r"initial address of send buffer",
    )
    parameter("count", "POLYXFER_NUM_ELEM_NNI", desc=r"number of elements sent")
    parameter("datatype", "DATATYPE", desc=r"type of each element")
    parameter("dest", "RANK", desc=r"rank of destination")
    parameter("tag", "TAG", desc=r"message tag")
    parameter("comm", "COMMUNICATOR")
    parameter("request", "REQUEST", direction="out")
\end{mpi-binding}

Creates a \mpiterm{persistent communication request}
for a \mpiterm{standard mode send}\mpitermindex{send!persistent} operation.

\cdeclindex{MPI\_Request}
\begin{mpi-binding}
    function_name("MPI_Bsend_init")
    parameter(
        "buf",
        "BUFFER",
        constant=True,
        asynchronous=True,
        desc=r"initial address of send buffer",
    )
    parameter("count", "POLYXFER_NUM_ELEM_NNI", desc=r"number of elements sent")
    parameter("datatype", "DATATYPE", desc=r"type of each element")
    parameter("dest", "RANK", desc=r"rank of destination")
    parameter("tag", "TAG", desc=r"message tag")
    parameter("comm", "COMMUNICATOR")
    parameter("request", "REQUEST", direction="out")
\end{mpi-binding}

Creates a \mpiterm{persistent communication request}
for a \mpiterm{buffered mode send} operation.

\cdeclindex{MPI\_Request}
\begin{mpi-binding}
    function_name("MPI_Ssend_init")
    parameter(
        "buf",
        "BUFFER",
        constant=True,
        asynchronous=True,
        desc=r"initial address of send buffer",
    )
    parameter("count", "POLYXFER_NUM_ELEM_NNI", desc=r"number of elements sent")
    parameter("datatype", "DATATYPE", desc=r"type of each element")
    parameter("dest", "RANK", desc=r"rank of destination")
    parameter("tag", "TAG", desc=r"message tag")
    parameter("comm", "COMMUNICATOR")
    parameter("request", "REQUEST", direction="out")
\end{mpi-binding}

Creates a \mpiterm{persistent communication request}
for a \mpiterm{synchronous mode send}\mpitermindex{synchronous mode send!persistent} operation.

\cdeclindex{MPI\_Request}
\begin{mpi-binding}
    function_name("MPI_Rsend_init")
    parameter(
        "buf",
        "BUFFER",
        constant=True,
        asynchronous=True,
        desc=r"initial address of send buffer",
    )
    parameter("count", "POLYXFER_NUM_ELEM_NNI", desc=r"number of elements sent")
    parameter("datatype", "DATATYPE", desc=r"type of each element")
    parameter("dest", "RANK", desc=r"rank of destination")
    parameter("tag", "TAG", desc=r"message tag")
    parameter("comm", "COMMUNICATOR")
    parameter("request", "REQUEST", direction="out")
\end{mpi-binding}

Creates a \mpiterm{persistent communication request}
for a \mpiterm{ready mode send}\mpitermindex{ready mode send!persistent} operation.

\cdeclindex{MPI\_Request}
\begin{mpi-binding}
    function_name("MPI_Recv_init")
    parameter(
        "buf",
        "BUFFER",
        direction="out",
        asynchronous=True,
        desc=r"initial address of receive buffer",
    )
    parameter(
        "count", "POLYXFER_NUM_ELEM_NNI", desc=r"number of elements received"
    )
    parameter("datatype", "DATATYPE", desc=r"type of each element")
    parameter(
        "source", "RANK", desc=r"rank of source or \mpiconst{MPI_ANY_SOURCE}"
    )
    parameter("tag", "TAG", desc=r"message tag or \mpiconst{MPI_ANY_TAG}")
    parameter("comm", "COMMUNICATOR")
    parameter("request", "REQUEST", direction="out")
\end{mpi-binding}

Creates a \mpiterm{persistent communication request}
for a receive operation\mpitermindex{receive!persistent}.
The argument \mpiarg{buf} is marked as \mpiarg{OUT}
because the user gives permission to write on the receive buffer\mpitermindex{receive!buffer} by passing the
argument to \mpifunc{MPI\_RECV\_INIT}.

A \mpiterm{persistent communication request} is \mpiterm{inactive} after it was created---no
active communication is attached to the request.

A communication that uses a \mpiterm{persistent communication request}
is \mpitermni{started}\mpitermindex{started!request objects}\mpitermindex{request objects!started}\mpitermindex{persistent communication request!starting} by the function \mpifunc{MPI\_START}.

\cdeclindex{MPI\_Request}
\begin{mpi-binding}
    function_name("MPI_Start")
    parameter("request", "REQUEST", direction="inout")
\end{mpi-binding}

The argument, \mpiarg{request}, is a handle returned by
any of the \mpiterm{initialization} procedures
for persistent point-to-point communication (the previous five procedures), or
for partitioned point-to-point communication (see Chapter~\ref{chap:part}), or
for persistent collective communication (see Sections~\ref{sec:coll-persistent} and~\ref{sec:sparsecoll-pnbc}).
The associated request should be \mpiterm{inactive}.
The request becomes \mpiterm{active} once the call is made.

If the request is for a \mpiterm{ready mode send} operation, then
a matching receive operation should be \mpitermni{started} before the call is made.  The
communication buffer must not be
modified
after the call, and
until the operation \mpitermni{completes}.

The call is \mpiterm{local}, with similar semantics to the nonblocking
communication operations described in
Section~\ref{sec:pt2pt-nonblock}.  That is,
a call to \mpifunc{MPI\_START} with a
request created by \mpifunc{MPI\_SEND\_INIT}
starts a
communication in the same manner as a call to \mpifunc{MPI\_ISEND};
a call to \mpifunc{MPI\_START} with a
request created by \mpifunc{MPI\_BSEND\_INIT}
starts a
communication in the same manner as a call to
\mpifunc{MPI\_IBSEND}; and so on.

\cdeclindex{MPI\_Request}
\begin{mpi-binding}
    function_name("MPI_Startall")
    parameter("count", "ARRAY_LENGTH_NNI", desc=r"list length")
    parameter(
        "array_of_requests",
        "REQUEST",
        direction="inout",
        length="count",
        desc=r"array of requests",
    )
\end{mpi-binding}

The execution of
\mpifunc{MPI\_STARTALL}
has the same effect as the execution of
\mpifunc{MPI\_START}
for each of the array elements
in some arbitrary order.
\mpifunc{MPI\_STARTALL} with an array of length one
is equivalent to \mpifunc{MPI\_START}.

A communication started with a call to \mpifunc{MPI\_START} or
\mpifunc{MPI\_STARTALL} is
completed by a call to \mpifunc{MPI\_WAIT}, \mpifunc{MPI\_TEST}, or
one of the derived functions described in
Section~\ref{subsec:pt2pt-multiple}.  The request becomes \mpiterm{inactive} after
successful completion of such call.  The request is not deallocated
and it can be activated anew by an \mpifunc{MPI\_START} or
\mpifunc{MPI\_STARTALL} call.

A \mpiterm{persistent communication request} is deallocated by a call to
\mpifunc{MPI\_REQUEST\_FREE}
(Section~\ref{subsec:pt2pt-commend}).
%
The call to \mpifunc{MPI\_REQUEST\_FREE} can occur at any point in the program
after the persistent request was created.  However, the request will be
deallocated only after it becomes \mpiterm{inactive}.
\mpitermni{Active} receive requests should not be \mpitermni{freed}\mpitermindex{request objects!freeing}. Otherwise, it will not be
possible to check that the receive has \mpitermni{completed}.
\mpitermni{Collective} operation requests
(defined in
Section~\ref{sec:nbcoll} and Section~\ref{sec:sparsecoll-nbc}
for nonblocking collective operations, and
Section~\ref{sec:coll-persistent} and Section~\ref{sec:sparsecoll-pnbc}
for persistent collective operations)
must not be \mpitermni{freed} while \mpiterm{active}.
It is preferable, in general, to free requests when they are inactive.  If this
rule is followed, then the functions
described in this section will be invoked
in a sequence of the form,
\[
\textbf{Create \ (Start \ Complete)$^*$ \ Free}
\]
where
$*$ indicates zero or more repetitions.
If the same \mpiterm{persistent communication request} is used in several concurrent
threads, it is the user's responsibility to coordinate calls so that the
correct sequence is obeyed.

\mpitermni{Inactive persistent requests} are not automatically \mpitermni{freed}\mpitermindex{persistent communication request!inactive!freeing}
when the associated communicator is disconnected (via \mpifunc{MPI\_COMM\_DISCONNECT}, see \sectionref{subsec:disconnect})
or the associated \wpm{} or \spm{} is finalized (via \mpifunc{MPI\_FINALIZE}, see \sectionref{sec:inquiry-finalize}, or \mpifunc{MPI\_SESSION\_FINALIZE}, see \sectionref{sec:session-init-finalize}).
In these situations, any further use of the request handle is erroneous.
In particular, freeing associated inactive request handles after such a communicator disconnect or finalization is then impossible.

\begin{users}
Persistent request handles may bind internal resources
such as \MPI/ buffers in shared memory for providing efficient communication.
Therefore, it is highly recommended to explicitly free inactive request handles,
using \mpifunc{MPI\_REQUEST\_FREE},  when they are no longer in use,
and in particular before freeing or disconnecting the associated communicator
with \mpifunc{MPI\_COMM\_FREE} or \mpifunc{MPI\_COMM\_DISCONNECT}
or finalizing the associated session with \mpifunc{MPI\_SESSION\_FINALIZE}.
\end{users}

A send operation \mpitermni{started}\mpitermindex{started!request objects}\mpitermindex{request objects!started}
with \mpifunc{MPI\_START} can be \mpitermni{matched}\mpitermindex{matching rules!persistent point-to-point} with
any receive operation and, likewise, a receive operation \mpitermni{started}
with \mpifunc{MPI\_START} can receive messages generated by any send
operation.

%% Note that a previous version of this included what it hoped were the
%% titles of the sections.  That hope was misplaced, as the section
%% titles had been changed.  Another example of the folly in trying to
%% keep multiple references manually synchronized.  The titles are really
%% not necessary here, as the first sentance covers the general issues.
%% Attempting to enumerate them all is also a mistake, as it leaves out
%% some that are also important.
\begin{users}
  To prevent problems with the argument copying and register optimization done
  by Fortran compilers, please note the hints in
Sections~\ref{sec:misc-problems}--\ref{sec:f90-problems:comparison-with-C}.
%% especially in
%% Sections~\ref{sec:misc-sequence} and~\ref{sec:f90-problems:vector-subscripts}
%% on pages~\pageref{sec:misc-sequence}-\pageref{sec:f90-problems:vector-subscripts}
%% about ``{\sf Problems Due to Data Copying and Sequence Association with Subscript Triplets}''
%% and ``{\sf Vector Subscripts}'',
%% and in Sections~\ref{sec:misc-register} to~\ref{sec:f90-problems:perm-data-movements}
%% on pages~\pageref{sec:misc-register} to~\pageref{sec:f90-problems:perm-data-movements}
%% about ``{\sf Optimization Problems}'', ``{\sf Code Movements and Register Optimization}'',
%% ``{\sf Temporary Data Movements}'' and ``{\sf Permanent Data Movements}''.
\end{users}

\section{Null \texorpdfstring{\MPI/}{MPI} Processes}
\mpitermtitleindex{null processes}
\mpitermtitleindex{matching rules!null process}
\label{sec:pt2pt-nullproc}

In many instances, it is convenient to specify a ``dummy'' source or
destination
for communication.  This simplifies the code that is needed for dealing with
boundaries, for example, in the case of a noncircular shift done with calls to
send-receive\mpifuncindex{MPI\_SENDRECV}\mpifuncindex{MPI\_SENDRECV\_REPLACE}.

The special value \mpiconstmain{MPI\_PROC\_NULL} can be used
instead of a rank wherever a
source or a destination argument is required in a call.  A communication
with \mpiconst{MPI\_PROC\_NULL} has no effect.
A send to \mpiconst{MPI\_PROC\_NULL}
succeeds and returns as soon as possible.
A receive from \mpiconst{MPI\_PROC\_NULL} succeeds and
returns as soon as possible with no modifications to the receive buffer.
When a receive with \mpiarg{source}\mpicode{ = }\mpiconst{MPI\_PROC\_NULL} is executed then
the status object returns \mpiarg{source}\mpicode{ = }\mpiconst{MPI\_PROC\_NULL},
\mpiarg{tag}\mpicode{ = }\mpiconst{MPI\_ANY\_TAG} and
\mpiarg{count}\mpicode{ = 0}.
A probe or matching probe with \mpiarg{source}\mpicode{ = }\mpiconst{MPI\_PROC\_NULL} succeeds and
returns as soon as possible, and the status object returns
\mpiarg{source}\mpicode{ = }\mpiconst{MPI\_PROC\_NULL},
\mpiarg{tag}\mpicode{ = }\mpiconst{MPI\_ANY\_TAG} and
\mpiarg{count}\mpicode{ = 0}.
A matching probe (cf.~Section~\ref{sec:matching-probe}) with
\mpiarg{source}\mpicode{ = }\mpiconst{MPI\_PROC\_NULL} returns \mpiarg{flag}\mpicode{ = }\mpicode{true},
\mpiarg{message}\mpicode{ = }\mpiconst{MPI\_MESSAGE\_NO\_PROC}, and the status object returns
\mpiarg{source}\mpicode{ = }\mpiconst{MPI\_PROC\_NULL}, \mpiarg{tag}\mpicode{ = }\mpiconst{MPI\_ANY\_TAG}, and
\mpiarg{count}\mpicode{ = 0}.
